% Lectura y comentario : « El medioambiente y nosotros » (residuos, recogida, reciclaje) — Espagnol Tle A — Cours signé OnBuch+
% Programme officiel MINESEC. Compiler avec : tectonic E07-medioambiente-residuos-reciclaje.tex
\newcommand{\DOCMATIERE}{Espagnol}
\newcommand{\DOCNIVEAU}{Tle A}
\newcommand{\DOCMODULE}{Módulo 2 : Environnement, santé et bien-être}
\newcommand{\DOCLECON}{E07}
\newcommand{\DOCTITRE}{Lectura y comentario : « El medioambiente y nosotros » (residuos, recogida, reciclaje)}
% OnBuch+ — préambule « cours signés » ENRICHI (compilé avec tectonic / XeLaTeX)
% Système visuel « Pop » d'OnBuch+ : fond crème (jamais blanc), bordures encre
% épaisses, ombres « dures » décalées, titres Archivo Black en capitales.
% Version MATHÉMATIQUES : graphes (pgfplots/tikz), boîtes pédagogiques
% (définition, propriété, méthode, attention, le savais-tu, exercice résolu,
% exercice, TP, bilan), page de garde et signature OnBuch+ ancrée en bas.
%
% Macros attendues dans main.tex AVANT \input{preamble} :
%   \DOCMATIERE  \DOCNIVEAU  \DOCMODULE  \DOCLECON (numéro)  \DOCTITRE
\documentclass[11pt]{article}

\usepackage{fontspec}
\usepackage[spanish,french]{babel} % français = langue principale ; l'espagnol via \esp{..} / espagnol
\usepackage[a4paper,margin=2cm,top=3.4cm,bottom=2.8cm,headheight=1.4cm,footskip=1.3cm]{geometry}
\usepackage{amsmath,amssymb,amsfonts}
\usepackage{mathtools} % \xmapsto, \coloneqq... souvent utilisés par les modèles
\usepackage{cancel} % simplifications barrées (bilans de forces)
% \abs{z} (module, valeur absolue) et \norm{u} (norme), utilisés par les modèles
\providecommand{\abs}{}\let\abs\relax\DeclarePairedDelimiter{\abs}{\lvert}{\rvert}
\providecommand{\norm}{}\let\norm\relax\DeclarePairedDelimiter{\norm}{\lVert}{\rVert}
\usepackage[table]{xcolor} % [table] : fournit aussi \rowcolors (alternance de lignes)
\usepackage{pagecolor}
\usepackage{tikz}
\usetikzlibrary{arrows.meta,positioning,calc,shapes.geometric,shapes.misc,shapes.arrows,decorations.pathmorphing,decorations.pathreplacing,decorations.markings,patterns,shadows,mindmap,trees,fit,backgrounds,matrix,intersections,angles,quotes}
\usepackage{pgfplots}
\usepackage[version=4]{mhchem} % équations nucléaires \ce{^{235}_{92}U}
\usepackage[european,siunitx]{circuitikz} % schémas électriques
\pgfplotsset{compat=1.18}
\usepgfplotslibrary{fillbetween,groupplots} % aires entre courbes (calcul intégral)
\usepackage{enumitem}
\usepackage{fancyhdr}
% [most] charge toutes les bibliothèques tcolorbox (listings, minted, raster,
% documentation...) et gonfle inutilement la mémoire TeX sur un document long
% (~300 boîtes) : on ne charge que ce qui est réellement utilisé.
\usepackage[skins,breakable]{tcolorbox}
\usepackage{graphicx}
\usepackage{colortbl}
\usepackage{booktabs}
\usepackage{tabularx}
\usepackage{diagbox} % case d'en-tête diagonale des tableaux à double entrée
% Colonnes extensibles centrée / gauche / droite (C, L, R), souvent utilisées par les modèles
\newcolumntype{C}{>{\centering\arraybackslash}X}
\newcolumntype{L}{>{\raggedright\arraybackslash}X}
\newcolumntype{R}{>{\raggedleft\arraybackslash}X}
\usepackage{array}
\usepackage{multicol}
\usepackage{siunitx}
% parse-numbers=false : accepte aussi \SI{2,5 \times 10^{-3}}{...}
\sisetup{locale=FR,output-decimal-marker={,},group-minimum-digits=4,parse-numbers=false}
\DeclareSIUnit\ohm{\ensuremath{\symup{\Omega}}} % U+2126 absent de la police
\DeclareSIUnit\division{div} % réglages d'oscilloscope (ms/div, V/div)
\DeclareSIUnit\tr{tr} % tours (vitesse de rotation, ex. \SI{1500}{\tr\per\minute})
\DeclareSIUnit\molar{M} % concentration molaire (ex. \SI{3}{\molar}, \SI{1}{\micro\molar})
\DeclareSIUnit\an{an} % année (le modèle confond parfois avec \year, un compteur TeX) — ex. \SI{5}{\kilo\meter\per\an}
\DeclareSIUnit\FCFA{FCFA} % franc CFA (ex. \SI{500}{\FCFA\per\kilo\gram})
\DeclareSIUnit\degreeBrix{\textdegree Brix} % teneur en sucre d'un sirop/jus (réfractométrie)
\DeclareSIUnit\month{mois} % le modèle utilise parfois \month, un compteur TeX, comme unité de durée
\usepackage{qrcode}
% \degree : commande fréquemment utilisée par les modèles pour un angle ou une
% température en dehors de \SI{}{} (non fournie par les paquets chargés).
\providecommand{\degree}{^{\circ}}
% \qed (fin de démonstration), utilisé par les modèles sans amsthm
\providecommand{\qed}{\hfill$\square$}
% \barre : double barre des valeurs interdites dans un tableau de variations (tkz-tab)
\providecommand{\barre}{\ensuremath{\|}}
% environnement proof (amsthm non chargé) utilisé par les modèles
\ifdefined\proof\else\newenvironment{proof}[1][Démonstration]{\par\noindent\textit{#1.}\ }{\qed\par}\fi
\usepackage{lastpage}
\usepackage{hyperref}
\hypersetup{hidelinks}

% ============================================================
% Palette « Pop » OnBuch+ — mêmes hex que la classe P (plus_ui.dart)
% ============================================================
\definecolor{poporange}{HTML}{F59321}
\definecolor{popdark}{HTML}{E07A0C}
\definecolor{popcream}{HTML}{FFF6E8}
\definecolor{popink}{HTML}{1C1714}
\definecolor{poppurple}{HTML}{7A5AE0}
\definecolor{popgreen}{HTML}{2E9E6A}
\definecolor{poppink}{HTML}{E0457B}
\definecolor{popblue}{HTML}{2F7DE1}
\definecolor{popgold}{HTML}{C98A12}
\definecolor{popmuted}{HTML}{8A7F76}
\definecolor{popline}{HTML}{EADFD2}
\definecolor{popgray}{HTML}{8A7F76}
\colorlet{popgrey}{popgray}
% Alias tolérants (le modèle nomme parfois les couleurs différemment)
\colorlet{popwhite}{white}
\colorlet{popblack}{popink}
\colorlet{popred}{poppink}
\colorlet{popyellow}{popgold}
% Teintes claires (fonds de tableaux/encadrés) — le modèle nomme parfois
% les couleurs avec un suffixe L (ex. popblueL) sans les avoir définies.
\colorlet{poporangeL}{poporange!20}
\colorlet{popdarkL}{popdark!20}
\colorlet{poppurpleL}{poppurple!20}
\colorlet{popgreenL}{popgreen!20}
\colorlet{poppinkL}{poppink!20}
\colorlet{popblueL}{popblue!20}
\colorlet{popgoldL}{popgold!20}
\colorlet{popredL}{popred!20}
\colorlet{popgrayL}{popgray!20}
% Teintes claires pour fonds de tableaux / zones
\colorlet{popblueL}{popblue!10!white}
\colorlet{poporangeL}{poporange!14!white}
\colorlet{popgreenL}{popgreen!12!white}
\colorlet{poppurpleL}{poppurple!12!white}
\colorlet{poppinkL}{poppink!12!white}
\colorlet{popgoldL}{popgold!14!white}

\pagecolor{popcream}

% ============================================================
% Typographie
% ============================================================
\defaultfontfeatures{Ligatures=TeX}
\setmainfont{PlusJakartaSans-Medium.ttf}[Path=../../../fonts/,BoldFont=PlusJakartaSans-Bold.ttf]
% Maths en sans-serif (Fira Math) avec chiffres et lettres droites de Plus
% Jakarta, pour que les formules se fondent dans le texte.
\usepackage[math-style=ISO,bold-style=ISO]{unicode-math}
\setmathfont{FiraMath-Regular.otf}[Path=../../../fonts/]
\setmathfont{PlusJakartaSans-Medium.ttf}[Path=../../../fonts/,range={up/num,up/{Latin,latin},\mathup}]
\usepackage{icomma} % virgule décimale sans espace en mode math
% Caractères mathématiques Unicode absents des polices de texte (Plus Jakarta,
% Archivo Black) : rendus par leur équivalent mathématique, en texte comme en
% formule — sinon ils s'affichent en carré vide (ex. « Arithmétique dans ℤ »).
\usepackage{newunicodechar}
\newunicodechar{ℝ}{\ensuremath{\mathbb{R}}}
\newunicodechar{ℤ}{\ensuremath{\mathbb{Z}}}
\newunicodechar{ℂ}{\ensuremath{\mathbb{C}}}
\newunicodechar{ℕ}{\ensuremath{\mathbb{N}}}
\newunicodechar{ℚ}{\ensuremath{\mathbb{Q}}}
\newunicodechar{𝔻}{\ensuremath{\mathbb{D}}}
\newunicodechar{α}{\ensuremath{\alpha}}
\newunicodechar{β}{\ensuremath{\beta}}
\newunicodechar{γ}{\ensuremath{\gamma}}
\newunicodechar{δ}{\ensuremath{\delta}}
\newunicodechar{ε}{\ensuremath{\varepsilon}}
\newunicodechar{θ}{\ensuremath{\theta}}
\newunicodechar{λ}{\ensuremath{\lambda}}
\newunicodechar{μ}{\ensuremath{\mu}}
\newunicodechar{σ}{\ensuremath{\sigma}}
\newunicodechar{τ}{\ensuremath{\tau}}
\newunicodechar{φ}{\ensuremath{\varphi}}
\newunicodechar{ψ}{\ensuremath{\psi}}
\newunicodechar{ω}{\ensuremath{\omega}}
\newunicodechar{Σ}{\ensuremath{\Sigma}}
% majuscules produites par \MakeUppercase dans les titres (α-aminés -> Α)
\newunicodechar{Α}{\ensuremath{\alpha}}
\newunicodechar{Β}{\ensuremath{\beta}}
\newunicodechar{Γ}{\ensuremath{\Gamma}}
\newunicodechar{Δ}{\ensuremath{\Delta}}
\newunicodechar{Ε}{\ensuremath{\varepsilon}}
\newunicodechar{Θ}{\ensuremath{\Theta}}
\newunicodechar{Λ}{\ensuremath{\Lambda}}
\newunicodechar{Μ}{\ensuremath{\mu}}
\newunicodechar{Τ}{\ensuremath{\tau}}
\newunicodechar{Φ}{\ensuremath{\Phi}}
\newunicodechar{Ψ}{\ensuremath{\Psi}}
\newunicodechar{Ω}{\ensuremath{\Omega}}
\newunicodechar{↦}{\ensuremath{\mapsto}}
\newunicodechar{∈}{\ensuremath{\in}}
\newunicodechar{∉}{\ensuremath{\notin}}
\newunicodechar{∧}{\ensuremath{\wedge}}
\newunicodechar{⇔}{\ensuremath{\Leftrightarrow}}
\newunicodechar{⟺}{\ensuremath{\Longleftrightarrow}}
\newunicodechar{⇒}{\ensuremath{\Rightarrow}}
\newunicodechar{⟹}{\ensuremath{\Longrightarrow}}
\newunicodechar{∘}{\ensuremath{\circ}}
\newunicodechar{∪}{\ensuremath{\cup}}
\newunicodechar{∩}{\ensuremath{\cap}}
\newunicodechar{≡}{\ensuremath{\equiv}}
\newunicodechar{⊂}{\ensuremath{\subset}}
\newunicodechar{⊥}{\ensuremath{\perp}}
\newunicodechar{∃}{\ensuremath{\exists}}
\newunicodechar{∀}{\ensuremath{\forall}}
\newunicodechar{∛}{\ensuremath{\sqrt[3]{\,}}}
\newunicodechar{∜}{\ensuremath{\sqrt[4]{\,}}}
\newunicodechar{⃗}{\ensuremath{{}^{\to}}}
\newunicodechar{ⁿ}{\ifmmode^{n}\else\textsuperscript{\ensuremath{n}}\fi}
\newunicodechar{ˣ}{\ifmmode^{x}\else\textsuperscript{\ensuremath{x}}\fi}
\newunicodechar{ᵇ}{\ifmmode^{b}\else\textsuperscript{\ensuremath{b}}\fi}
\newunicodechar{ʳ}{\ifmmode^{r}\else\textsuperscript{\ensuremath{r}}\fi}
\newunicodechar{ᵘ}{\ifmmode^{u}\else\textsuperscript{\ensuremath{u}}\fi}
\newunicodechar{ʸ}{\ifmmode^{y}\else\textsuperscript{\ensuremath{y}}\fi}
\newunicodechar{ᵃ}{\ifmmode^{a}\else\textsuperscript{\ensuremath{a}}\fi}
\newunicodechar{ᵉ}{\ifmmode^{e}\else\textsuperscript{\ensuremath{e}}\fi}
\newunicodechar{ᵏ}{\ifmmode^{k}\else\textsuperscript{\ensuremath{k}}\fi}
\newunicodechar{ᵖ}{\ifmmode^{p}\else\textsuperscript{\ensuremath{p}}\fi}
\newunicodechar{ᴺ}{\ifmmode^{N}\else\textsuperscript{\ensuremath{N}}\fi}
\newunicodechar{ᶜ}{\ifmmode^{c}\else\textsuperscript{\ensuremath{c}}\fi}
\newunicodechar{ʰ}{\ifmmode^{h}\else\textsuperscript{\ensuremath{h}}\fi}
\newunicodechar{ᵠ}{\ifmmode^{q}\else\textsuperscript{\ensuremath{q}}\fi}
\newunicodechar{ₙ}{\ifmmode_{n}\else\textsubscript{\ensuremath{n}}\fi}
\newunicodechar{ₐ}{\ifmmode_{a}\else\textsubscript{\ensuremath{a}}\fi}
\newunicodechar{ᵢ}{\ifmmode_{i}\else\textsubscript{\ensuremath{i}}\fi}
\newunicodechar{ᵣ}{\ifmmode_{r}\else\textsubscript{\ensuremath{r}}\fi}
\newunicodechar{ₖ}{\ifmmode_{k}\else\textsubscript{\ensuremath{k}}\fi}
\newunicodechar{ᵦ}{\ifmmode_{\beta}\else\textsubscript{\ensuremath{\beta}}\fi}
\newunicodechar{ₓ}{\ifmmode_{x}\else\textsubscript{\ensuremath{x}}\fi}
\newfontfamily\popdisplay{ArchivoBlack-Regular.ttf}[Path=../../../fonts/]
\newfontfamily\poplogo{SpaceGrotesk-Bold.ttf}[Path=../../../fonts/]
\newfontfamily\popsemi{PlusJakartaSans-SemiBold.ttf}[Path=../../../fonts/]
\newfontfamily\popxbold{PlusJakartaSans-ExtraBold.ttf}[Path=../../../fonts/]
\newfontfamily\popxboldls{PlusJakartaSans-ExtraBold.ttf}[Path=../../../fonts/,LetterSpace=3.0]

\setlist[enumerate]{leftmargin=1.7em,itemsep=5pt,topsep=4pt}
\setlist[itemize]{leftmargin=1.7em,itemsep=4pt,topsep=4pt,label={\color{poporange}\textbullet}}
\setlength{\parindent}{0pt}
\setlength{\parskip}{5pt}
\renewcommand{\arraystretch}{1.25}

% pgfplots : style Pop par défaut
\pgfplotsset{
  every axis/.append style={axis line style={popink,line width=1pt},tick style={popink},
    grid style={popline,line width=0.5pt},label style={font=\small},tick label style={font=\footnotesize}},
  popcurve/.style={poporange,line width=1.8pt,no markers,smooth},
}
% Même style disponible dans un \draw TikZ simple (hors pgfplots)
\tikzset{popcurve/.style={poporange,line width=1.8pt}}
\tikzset{popgrid/.style={popline,very thin}}
\colorlet{popcurve}{poporange} % le nom du style est parfois utilisé comme couleur

% ============================================================
% Pill / squiggle
% ============================================================
\newcommand{\poppill}[3]{%
  \tikz[baseline=-0.55ex]\node[draw=popink,line width=1.2pt,rounded corners=2.6pt,%
    fill=#2,inner xsep=6.5pt,inner ysep=3pt]{%
    \popxboldls\fontsize{7}{7}\selectfont\color{#3}\MakeUppercase{#1}};%
}
\newcommand{\popsquiggle}[1]{%
  \tikz[baseline=-0.4ex]{
    \draw[#1,line width=2.6pt,line cap=round]
      plot [smooth] coordinates {(0,0.09) (0.34,0.01) (0.68,0.21) (1.02,0.01) (1.36,0.10)};
  }%
}
% Mot-clé surligné (marqueur orange sous le texte)
\newcommand{\cle}[1]{{\popxbold\color{popdark}#1}}

% ============================================================
% En-tête / pied
% ============================================================
\pagestyle{fancy}
\fancyhf{}
\renewcommand{\headrulewidth}{0pt}
\renewcommand{\footrulewidth}{0pt}
\lhead{\raisebox{-0.15ex}{\poplogo\fontsize{13}{13}\selectfont\color{popink}OnBuch\color{poporange}+}}
\rhead{\poppill{\DOCMATIERE\ \textbullet\ \DOCNIVEAU\ \textbullet\ Leçon \DOCLECON}{white}{popink}}
\lfoot{\popxbold\fontsize{7.6}{7.6}\selectfont\color{popmuted}\MakeUppercase{Cours signé OnBuch+}\ \textbullet\ {\popsemi\DOCTITRE}}
\rfoot{\tikz[baseline=-0.6ex]\node[rounded corners=8.5pt,fill=popink,minimum height=17pt,minimum width=17pt,inner xsep=5pt,inner sep=0]{\popxbold\fontsize{8}{8}\selectfont\color{popcream}\thepage\,/\,\pageref*{LastPage}};}

% ============================================================
% Titres
% ============================================================
\newcommand{\chaptitre}[2]{%
  \begin{tcolorbox}[enhanced,colback=poporange,colframe=popink,boxrule=2.6pt,arc=11pt,
      drop shadow=popink,left=15pt,right=15pt,top=12pt,bottom=11pt,before skip=3pt,after skip=18pt]
    \poppill{#2}{popink}{popcream}\par\vspace{9pt}
    {\raggedright\hyphenpenalty=10000\exhyphenpenalty=10000\popdisplay\fontsize{22}{24}\selectfont\color{popcream}\MakeUppercase{#1}\par}
  \end{tcolorbox}
}
% Section numérotée : pastille orange avec le numéro + titre Archivo
\newcounter{popsec}
\newcommand{\coursec}[1]{%
  \stepcounter{popsec}\setcounter{popsub}{0}%
  \par\vspace{16pt}\noindent
  \begin{minipage}{\linewidth}
  \tikz[baseline=-0.7ex]\node[draw=popink,line width=1.6pt,fill=poporange,rounded corners=4pt,minimum size=22pt,inner sep=0]{\popdisplay\fontsize{12}{12}\selectfont\color{popcream}\thepopsec};%
  \hspace{8pt}{\popdisplay\fontsize{15}{17}\selectfont\color{popink}\MakeUppercase{#1}}\par
  \vspace{4pt}\noindent{\color{poporange}\rule{36pt}{3.6pt}}
  \end{minipage}\par\nopagebreak\vspace{8pt}
}
\newcounter{popsub}
\newcommand{\courssub}[1]{%
  \stepcounter{popsub}%
  \par\vspace{9pt}\noindent{\popxbold\fontsize{11.5}{13}\selectfont\color{popink}{\color{poporange}\thepopsec.\thepopsub}\hspace{6pt}#1}\par\nopagebreak\vspace{4pt}
}

% ============================================================
% Boîtes pédagogiques — même recette : fond blanc, bordure épaisse colorée,
% ombre dure encre, badge plein. Titre optionnel [#1].
% ============================================================
\tcbset{popbase/.style={breakable,enhanced,colback=white,boxrule=2.3pt,arc=9pt,
  shadow={3pt}{-3pt}{0pt}{popink},left=14pt,right=14pt,top=11pt,bottom=11pt,before skip=11pt,after skip=15pt,
  fontupper=\popsemi\fontsize{10.6}{14.5}\selectfont\color{popink},
  fontlower=\popsemi\fontsize{10.6}{14.5}\selectfont\color{popink}}}
% Coche dessinée (le glyphe ✓ n'existe pas dans les polices du document)
\newcommand{\popcheck}[1]{\tikz[baseline=-0.5ex]\draw[#1,line width=1.6pt,line cap=round,line join=round](-0.13,0.0)--(-0.04,-0.09)--(0.13,0.1);}
\providecommand{\coche}{\popcheck{popgreen}}
\providecommand{\resultat}[1]{\par\smallskip\noindent\fcolorbox{popgreen}{popgreenL}{\parbox{\dimexpr\linewidth-2\fboxsep-2\fboxrule\relax}{\textbf{Résultat.} #1}}\par\smallskip}
\newcommand{\popboxhead}[3]{\poppill{#1}{#2}{white}\ifx\relax#3\relax\else\hspace{6pt}{\popxbold\fontsize{10.6}{12}\selectfont\color{popink}#3}\fi\par\vspace{7pt}}

\newtcolorbox{aretenir}[1][]{popbase,colframe=popgreen,before upper={\popboxhead{À retenir}{popgreen}{#1}}}
% Texte en espagnol (document de lecture, dialogue, poème, extrait) : cadre
% d'étude. Un titre optionnel (source, auteur) s'affiche dans l'en-tête.
\newtcolorbox{textebox}[1][]{popbase,colframe=popblue,before upper={\popboxhead{Texto}{popblue}{#1}\begin{otherlanguage}{spanish}},after upper={\end{otherlanguage}}}
% Marques ✓ / ✗ (forme correcte / fautive) souvent utilisées par les modèles
\providecommand{\cross}{\textcolor{poppink}{\ensuremath{\times}}}
\providecommand{\xmark}{\cross}\providecommand{\crossmark}{\cross}\providecommand{\cmark}{\ensuremath{\checkmark}}
% Ligne de réponse / justification à compléter (inventée par les modèles)
\providecommand{\justification}[1]{\par\noindent\textit{Justification~:} \dotfill\par}
% \textipa (package tipa non chargé) : la police ne couvre pas l'API -> simple passage
\providecommand{\textipa}[1]{#1}
% Soulignement (ulem non chargé) : \uline, \ul
\providecommand{\uline}[1]{\underline{#1}}\providecommand{\ul}[1]{\underline{#1}}
% Mot ou phrase en espagnol à l'intérieur d'un paragraphe français
\newcommand{\esp}[1]{\ifmmode\text{\textit{#1}}\else\begingroup\selectlanguage{spanish}\itshape #1\endgroup\fi}
\newtcolorbox{exemplebox}[1][]{popbase,colframe=poporange,before upper={\popboxhead{Exemple}{poporange}{#1}}}
\newtcolorbox{definition}[1][]{popbase,colframe=popblue,before upper={\popboxhead{Définition}{popblue}{#1}}}
\newtcolorbox{propriete}[1][]{popbase,colframe=poppurple,before upper={\popboxhead{Propriété}{poppurple}{#1}}}
\newtcolorbox{methode}[1][]{popbase,colframe=popgold,before upper={\popboxhead{Méthode}{popgold}{#1}}}
\newtcolorbox{attention}[1][]{popbase,colframe=poppink,before upper={\popboxhead{Attention}{poppink}{#1}}}
\newtcolorbox{experience}[1][]{popbase,colframe=popblue,colback=popblueL,before upper={\popboxhead{Expérience}{popblue}{#1}}}
% « Le savais-tu ? » : carte violette pleine (contexte, culture, Cameroun)
\newtcolorbox{savaistu}[1][]{popbase,colback=poppurple,colframe=popink,
  fontupper=\popsemi\fontsize{10.4}{14.2}\selectfont\color{white},
  before upper={\poppill{Le savais-tu\,?}{white}{poppurple}\ifx\relax#1\relax\else\hspace{6pt}{\popxbold\color{white}#1}\fi\par\vspace{7pt}}}
% Exercice résolu : énoncé en haut, \tcblower puis la solution sur fond crème
\newcounter{popexo}\newcounter{popexr}
\newtcolorbox{exoresolu}[1][]{popbase,colframe=poporange,colbacklower=popcream,
  segmentation style={popink,line width=1.2pt,dashed},
  before upper={\stepcounter{popexr}\popboxhead{Exercice résolu \thepopexr}{poporange}{#1}},
  before lower={\poppill{Solution}{popink}{popcream}\par\vspace{6pt}}}
% Exercice (non résolu en place) — la correction est dans la partie Corrigés
\newtcolorbox{exercice}[1][]{popbase,colframe=popink,
  before upper={\stepcounter{popexo}\popboxhead{Exercice \thepopexo}{popink}{#1}}}
% Corrigé
\newtcolorbox{corrige}[1][]{popbase,colframe=popgreen,colback=popgreenL,
  before upper={\popboxhead{Corrigé}{popgreen}{#1}}}
% Objectifs / prérequis / bilan
\newtcolorbox{objectifs}{popbase,colframe=popink,colback=poporangeL,
  before upper={\popboxhead{Objectifs de la leçon}{popink}{}}}
\newtcolorbox{prerequis}{popbase,colframe=popink,colback=popblueL,
  before upper={\popboxhead{Prérequis}{popblue}{}}}
\newtcolorbox{bilan}{popbase,colframe=popink,colback=white,boxrule=2.8pt,
  before upper={\popboxhead{Fiche bilan}{poporange}{L'essentiel à connaître pour l'examen}}}
% Figure encadrée avec légende
\newtcolorbox{popfigure}[1][]{enhanced,breakable=false,colback=white,colframe=popline,boxrule=1.4pt,arc=8pt,
  left=10pt,right=10pt,top=10pt,bottom=8pt,before skip=10pt,after skip=12pt,halign=center,
  lower separated=false,halign lower=center,
  fontlower=\popsemi\fontsize{9}{11.5}\selectfont\color{popmuted},
  #1}
\newcommand{\legende}[1]{\par\vspace{6pt}{\popsemi\fontsize{9}{11.5}\selectfont\color{popmuted}#1}}

% ============================================================
% Signature OnBuch+
% ============================================================
\newcommand{\poplogoinline}[1]{{\poplogo\fontsize{#1}{#1}\selectfont\color{popink}OnBuch\color{poporange}+}}
\newcommand{\signatureonbuch}{%
  \begin{tcolorbox}[enhanced,colback=popink,colframe=popink,boxrule=2.2pt,arc=10pt,
      drop shadow=poporange,left=14pt,right=14pt,top=10pt,bottom=10pt,before skip=0pt,after skip=0pt]
    \tikz[baseline=-0.7ex]\node[circle,fill=popgreen,draw=popcream,line width=1.3pt,minimum size=22pt,inner sep=0]{\popcheck{white}};%
    \hspace{9pt}\begin{minipage}[c]{0.62\linewidth}
      {\popxbold\fontsize{12}{13}\selectfont\color{popcream}Signé\ \textbullet\ {\poplogo OnBuch\color{poporange}+}}\par\vspace{2pt}
      {\raggedright\popsemi\fontsize{8}{10.5}\selectfont\color{popline}\MakeUppercase{Équipe pédagogique OnBuch+}\\\MakeUppercase{Conforme au programme officiel MINESEC}\par}
    \end{minipage}\hfill
    {\poplogo\fontsize{22}{22}\selectfont\color{popcream}OnBuch\color{poporange}+}
  \end{tcolorbox}%
}

% ============================================================
% Page de garde
% #1 titre · #2 matière · #3 niveau · #4 module · #5 n° leçon · #6 durée ·
% #7 description / objectifs (texte ou itemize)
% ============================================================
\newcommand{\pagegarde}[7]{%
  \thispagestyle{empty}
  \newgeometry{a4paper,margin=1.7cm,top=1.6cm,bottom=1.4cm}%
  \begin{tikzpicture}[remember picture,overlay]
    \fill[poporange!18] ([xshift=-3.2cm,yshift=-3.6cm]current page.north east) circle (4.6cm);
    \fill[poppurple!14] ([xshift=2.4cm,yshift=7.6cm]current page.south west) circle (3.8cm);
  \end{tikzpicture}%
  {\poplogo\fontsize{30}{30}\selectfont\color{popink}OnBuch\color{poporange}+}\hfill
  \raisebox{0.6ex}{\poppill{Cours signé \textbullet\ Premium}{popink}{popcream}}\par
  \vspace{1pt}\popsquiggle{poppurple}\par
  \vspace{8pt}
  \begin{tcolorbox}[enhanced,colback=poporange,colframe=popink,boxrule=2.8pt,arc=13pt,
      drop shadow=popink,left=18pt,right=18pt,top=12pt,bottom=13pt,after skip=10pt]
    \poppill{#2 \textbullet\ #3}{popink}{popcream}\hspace{5pt}\poppill{#4}{white}{popink}\par\vspace{8pt}
    {\popdisplay\fontsize{14}{14}\selectfont\color{popink}LEÇON #5}\par\vspace{4pt}
    {\raggedright\hyphenpenalty=10000\exhyphenpenalty=10000\popdisplay\fontsize{25}{27}\selectfont\color{popcream}\MakeUppercase{#1}\par}
    \vspace{8pt}
    {\popxbold\fontsize{9.5}{9.5}\selectfont\color{popink}\MakeUppercase{Durée conseillée : #6}}
  \end{tcolorbox}
  \begin{tcolorbox}[enhanced,colback=white,colframe=popink,boxrule=2.2pt,arc=10pt,
      drop shadow=popink,left=16pt,right=16pt,top=10pt,bottom=10pt,after skip=8pt]
    \poppill{Dans cette leçon}{poppurple}{white}\par\vspace{5pt}
    \popsemi\fontsize{9.3}{12.6}\selectfont\color{popink}#7
  \end{tcolorbox}
  \noindent\begin{minipage}[t]{0.487\linewidth}
    \begin{tcolorbox}[enhanced,colback=popgreen,colframe=popink,boxrule=2.2pt,arc=10pt,
        drop shadow=popink,left=11pt,right=11pt,top=10pt,bottom=10pt,height=3.1cm,valign=center]
      \tcbox[colback=white,colframe=popink,boxrule=1.3pt,arc=5pt,boxsep=0pt,nobeforeafter,
        left=3pt,right=3pt,top=3pt,bottom=3pt,valign=center]{\qrcode[height=1.9cm]{https://play.google.com/store/apps/details?id=com.luvvix.onbuch&hl=fr}}%
      \hspace{8pt}\begin{minipage}[c]{0.5\linewidth}
        {\popdisplay\fontsize{11}{12.5}\selectfont\color{white}\MakeUppercase{Télécharge OnBuch}}\par\vspace{4pt}
        {\raggedright\popsemi\fontsize{8.4}{11}\selectfont\color{white}Gratuit \textbullet\ Google Play\\Tuteur IA, annales, concours, résultats\par}
      \end{minipage}
    \end{tcolorbox}
  \end{minipage}\hfill
  \begin{minipage}[t]{0.487\linewidth}
    \begin{tcolorbox}[enhanced,colback=poppurple,colframe=popink,boxrule=2.2pt,arc=10pt,
        drop shadow=popink,left=11pt,right=11pt,top=10pt,bottom=10pt,height=3.1cm,valign=center]
      \tikz[baseline=-0.7ex]\node[circle,fill=white,draw=popink,line width=1.3pt,minimum size=1.6cm,inner sep=0]{\popdisplay\fontsize{24}{24}\selectfont\color{poppurple}+};%
      \hspace{8pt}\begin{minipage}[c]{0.6\linewidth}
        {\popdisplay\fontsize{11}{12.5}\selectfont\color{white}\MakeUppercase{Passe à OnBuch+}}\par\vspace{4pt}
        {\raggedright\popsemi\fontsize{8.4}{11}\selectfont\color{white}Tous les cours signés, Tuteur IA illimité, fiches PDF — onglet OnBuch+ de l'app\par}
      \end{minipage}
    \end{tcolorbox}
  \end{minipage}
  \vspace{8pt}
  \popcontenu
  \vfill
  \signatureonbuch
  \restoregeometry
  \newpage
}

% Rangée « Ce cours contient » (page de garde)
\newcommand{\poptile}[3]{% #1 couleur #2 gros texte #3 légende
  \begin{tcolorbox}[enhanced,colback=#1,colframe=popink,boxrule=2pt,arc=9pt,shadow={2.5pt}{-2.5pt}{0pt}{popink},
      left=6pt,right=6pt,top=9pt,bottom=9pt,halign=center,valign=center,height=1.75cm,width=\linewidth,nobeforeafter]
    {\popdisplay\fontsize{12.5}{13}\selectfont\color{white}\MakeUppercase{#2}}\par\vspace{3pt}
    {\centering\popsemi\fontsize{7.8}{9.5}\selectfont\color{white}#3\par}
  \end{tcolorbox}}
\newcommand{\popcontenu}{%
  \par\noindent{\popxboldls\fontsize{7.5}{7.5}\selectfont\color{popmuted}CE COURS SIGNÉ CONTIENT}\par\vspace{7pt}
  \noindent\begin{minipage}[t]{0.235\linewidth}\poptile{popblue}{Cours}{complet, illustré, pas à pas}\end{minipage}\hfill
  \begin{minipage}[t]{0.235\linewidth}\poptile{poporange}{Méthodes}{et exercices résolus}\end{minipage}\hfill
  \begin{minipage}[t]{0.235\linewidth}\poptile{poppink}{Exercices}{type examen + corrigés}\end{minipage}\hfill
  \begin{minipage}[t]{0.235\linewidth}\poptile{popgreen}{Fiche bilan}{l'essentiel à retenir}\end{minipage}\par}

% Sommaire
\newtcolorbox{sommaire}{enhanced,colback=white,colframe=popink,boxrule=2.2pt,arc=10pt,
  drop shadow=popink,left=18pt,right=18pt,top=13pt,bottom=13pt,before skip=6pt,after skip=20pt,
  before upper={\poppill{Sommaire}{poppurple}{white}\par\vspace{9pt}\popsemi\fontsize{10.6}{14.5}\selectfont\color{popink}}}

% Signature de fin de cours — poussée tout en bas de la dernière page
\newcommand{\findecours}{%
  \par\vfill\nopagebreak
  \begin{center}\popsquiggle{poporange}\end{center}\vspace{4pt}
  \signatureonbuch
}
\AtBeginDocument{\def\llbracket{\mathopen{[\mkern-3.5mu[}}\def\rrbracket{\mathclose{]\mkern-3.5mu]}}} % intervalles d'entiers (glyphe ⟦ absent de la police)
% Glyphes absents de Fira Math : redéfinis à partir de glyphes disponibles
\AtBeginDocument{%
  \def\setminus{\mathbin{\backslash}}\let\smallsetminus\setminus
  \def\vdots{\mathord{\vbox{\baselineskip3.2pt\lineskiplimit0pt\kern4pt\hbox{.}\hbox{.}\hbox{.}}}}%
  \def\ddots{\mathinner{\mkern1mu\raise6.5pt\hbox{.}\mkern2mu\raise3.6pt\hbox{.}\mkern2mu\raise0.7pt\hbox{.}\mkern1mu}}%
  \def\triangle{\mathord{\scalebox{0.9}{$\symup{\Delta}$}}}%
  \def\nabla{\mathord{\raisebox{\depth}{\scalebox{1}[-1]{$\symup{\Delta}$}}}}%
  \def\checkmark{\popcheck{popdark}}%
  \def\star{\mathbin{\ast}}\def\bigstar{\mathord{\ast}}%
  \def\diamond{\mathbin{\scalebox{0.6}{$\Diamond$}}}%
  \def\bigcup{\mathop{\mathchoice{\raisebox{-0.3em}{\scalebox{1.7}{$\cup$}}}{\scalebox{1.25}{$\cup$}}{\cup}{\cup}}\displaylimits}%
  \def\bigcap{\mathop{\mathchoice{\raisebox{-0.3em}{\scalebox{1.7}{$\cap$}}}{\scalebox{1.25}{$\cap$}}{\cap}{\cap}}\displaylimits}%
  \def\lesseqqgtr{\mathrel{\lessgtr}}\def\bigoplus{\mathop{\oplus}\displaylimits}%
}
\providecommand{\ssi}{si et seulement si} % abréviation fréquente
\AtBeginDocument{\providecommand{\wideparen}{\overparen}} % arc de cercle

\begin{document}

\pagegarde{Lectura y comentario : « El medioambiente y nosotros » (residuos, recogida, reciclaje)}{Espagnol}{Tle A}{Módulo 2 : Environnement, santé et bien-être}{E07}{2 h}{Cette leçon de lecture-commentaire propose un texte original en espagnol sur les déchets, la collecte et le recyclage, avec glossaire et questions de compréhension. Elle fournit le lexique thématique complet, la méthode du commentaire de texte au Bac et une production guidée de propositions écocitoyennes.
\vspace{4pt}
\begin{multicols}{2}\small
\begin{itemize}
\item À la fin de cette leçon, je sais lire et comprendre globalement un texte espagnol sur l'environnement et les déchets.
\item Je sais repérer les informations explicites et implicites d'un texte (qui, quoi, où, pourquoi).
\item Je sais utiliser le lexique des déchets et du recyclage : basura, residuos, contenedor, recogida selectiva, reciclaje, vertedero, reutilizar, reducir.
\item Je sais décrire la chaîne de collecte et de recyclage en espagnol.
\item Je sais rédiger un commentaire de texte structuré (introduction, développement, conclusion) selon la méthode du Bac.
\item Je sais reformuler une idée du texte avec mes propres mots sans copier les phrases.
\end{itemize}\end{multicols}}

\begin{sommaire}
\begin{multicols}{2}
\begin{enumerate}
\item Texte support : « El medioambiente y nosotros »
\item Compréhension du texte
\item Lexique thématique : déchets et recyclage
\item Méthode du commentaire de texte
\item Commentaire modèle
\item Production guidée : gestes écocitoyens
\item Activité d'intégration
\item Méthodes et exercices résolus
\item Exercices
\item Corrigés des exercices
\item Fiche bilan
\end{enumerate}
\end{multicols}
\end{sommaire}

\begin{objectifs}
\begin{itemize}
\item À la fin de cette leçon, je sais lire et comprendre globalement un texte espagnol sur l'environnement et les déchets.
\item Je sais repérer les informations explicites et implicites d'un texte (qui, quoi, où, pourquoi).
\item Je sais utiliser le lexique des déchets et du recyclage : basura, residuos, contenedor, recogida selectiva, reciclaje, vertedero, reutilizar, reducir.
\item Je sais décrire la chaîne de collecte et de recyclage en espagnol.
\item Je sais rédiger un commentaire de texte structuré (introduction, développement, conclusion) selon la méthode du Bac.
\item Je sais reformuler une idée du texte avec mes propres mots sans copier les phrases.
\item Je sais proposer des gestes écocitoyens en espagnol avec les structures de conseil et d'obligation (hay que, deberíamos, es importante + infinitif).
\item Je sais comparer la situation du recyclage en Espagne et au Cameroun dans un court paragraphe argumenté.
\end{itemize}
\end{objectifs}

\begin{prerequis}
\begin{itemize}
\item Lexique général de l'environnement vu en Première (naturaleza, contaminación, proteger).
\item Présent de l'indicatif des verbes réguliers et irréguliers fréquents.
\item Structures d'obligation et de conseil : hay que + infinitif, tener que, deber + infinitif.
\item Pronoms relatifs simples (que, donde) pour comprendre les phrases complexes du texte.
\item Méthode de compréhension écrite vue en Première (repérage du thème, de la source, des idées principales).
\end{itemize}
\end{prerequis}

\newpage

\coursec{Texte support : « El medioambiente y nosotros »}

\courssub{Situation d'accroche et problématique}

Imagine la scène : il est sept heures du matin à Douala, quartier Deïdo. Les \esp{bendskins} slaloment entre les voitures, les mamans ouvrent leurs étals, et devant les maisons s'entassent des sacs plastiques, des bouteilles vides, des restes de manioc et de poisson. Chaque jour, Douala et Yaoundé produisent des milliers de tonnes de déchets : bouteilles en plastique, emballages, restes de nourriture... Où vont-ils ? Certains finissent dans les caniveaux, d'autres dans des décharges sauvages, et quand la pluie tombe, tout cela bouche les égouts et provoque des inondations.

Pendant ce temps, à Madrid ou à Barcelone, une scène très différente se déroule : les citoyens sortent le soir avec plusieurs sacs de couleurs différentes. Le plastique va dans le conteneur jaune, le papier dans le bleu, le verre dans le vert. Là-bas, la \esp{recogida selectiva} — le tri sélectif — est une habitude quotidienne, presque un réflexe.

Et si tu pouvais, en espagnol, comprendre un texte sur ce sujet d'actualité brûlant et proposer toi-même des gestes écocitoyens pour ta ville ? C'est exactement ce que le Bac te demandera de faire : lire un texte de quinze à vingt lignes, le comprendre, le commenter et mobiliser un lexique précis.

\textbf{Problématique :} comment lire efficacement un texte espagnol sur l'environnement, en dégager les idées principales et s'approprier le vocabulaire des déchets et du recyclage pour pouvoir, ensuite, s'exprimer à son tour ?

\courssub{Présentation du texte}

\begin{definition}[Fiche d'identité du texte]
\begin{itemize}
\item \textbf{Source :} manuel \emph{Didáctica V}, unité 2 (module « Environnement, santé et bien-être »).
\item \textbf{Titre :} \esp{El medioambiente y nosotros} (« L'environnement et nous »).
\item \textbf{Type de texte :} texte journalistique à visée argumentative (sensibilisation écologique).
\item \textbf{Thème :} les déchets, la collecte sélective et le recyclage, en comparant le contexte camerounais et le contexte espagnol.
\item \textbf{Longueur :} environ 18 lignes, format typique de l'épreuve de compréhension de l'écrit au Bac.
\end{itemize}
\end{definition}

La lecture se fait en deux temps. D'abord une \textbf{lecture silencieuse} individuelle : tu parcours le texte seul, sans dictionnaire, pour saisir le sens global. Ensuite, l'enseignant fait une \textbf{lecture à voix haute} : tu écoutes la prononciation (attention à \esp{reciclaje} [re-\textbf{θ}i-\textbf{kl}a-\textbf{x}e] en Espagne péninsulaire, [re-\textbf{s}i-\textbf{kl}a-\textbf{s}e] en Amérique latine et en Guinée équatoriale ; \esp{recogida} [re-ko-\textbf{ð}i-\textbf{ð}a] ; \esp{vertedero} [ber-\textbf{ð}e-\textbf{ð}e-ro]) et tu vérifies tes hypothèses.

Pendant cette première lecture, ton œil de francophone doit repérer les \cle{mots transparents}, c'est-à-dire les mots qui ressemblent au français et se devinent sans effort :

\begin{center}
\begin{tabularx}{\linewidth}{|X|X|X|}\hline
\rowcolor{popblueL} \textbf{Mot espagnol} & \textbf{Français} & \textbf{Type} \\ \hline
\esp{reciclaje} & recyclage & mot transparent \\ \hline
\esp{contenedor} & conteneur & mot transparent \\ \hline
\esp{selectiva} & sélective & mot transparent \\ \hline
\esp{plástico} & plastique & mot transparent \\ \hline
\esp{orgánica} & organique & mot transparent \\ \hline
\esp{transforman} & transforment & mot transparent \\ \hline
\esp{planeta} & planète & mot transparent \\ \hline
\esp{contaminen} & contaminent & mot transparent \\ \hline
\end{tabularx}
\end{center}

\begin{attention}[Mots transparents, mais pas toujours amis]
Un mot transparent aide à deviner le sens global, mais méfie-toi : \esp{residuos} ne veut pas dire « résidus » au sens chimique, mais simplement « déchets » ; \esp{el suelo} signifie « le sol » (la terre), et non « le sol » au sens de la monnaie ; et \esp{una costumbre} est « une habitude », pas « un costume » (qui se dit \esp{un traje} ou \esp{un disfraz}). Vérifie toujours avec le contexte.
\end{attention}

\courssub{Lecture du texte intégral}

\begin{textebox}[El medioambiente y nosotros — texto de lectura]
En las grandes ciudades como Douala y Yaoundé, la basura se acumula en las calles. Cada habitante produce casi un kilo de residuos al día. Bolsas de plástico, botellas vacías, restos de comida: todo termina demasiado a menudo en las cunetas o en los ríos, y cuando llegan las lluvias, los barrios se inundan.

Sin embargo, en España la recogida selectiva es una costumbre: los ciudadanos separan el plástico, el vidrio, el papel y la materia orgánica en contenedores de colores. El contenedor amarillo recibe los envases, el azul el papel y el cartón, y el verde el vidrio. Nadie pregunta ya dónde hay que tirar una botella: todo el mundo lo sabe desde la escuela.

Después, los camiones llevan los residuos a las plantas de reciclaje, donde se transforman en nuevos productos. Una botella de plástico puede convertirse en una chaqueta; un periódico viejo, en una caja de cartón; las cáscaras de fruta, en abono para los campos.

Reciclar, reutilizar y reducir son las tres erres que pueden salvar nuestro planeta. Si todos colaboramos, evitaremos que los vertederos contaminen el suelo y el agua. En Douala y en Yaoundé, algunos jóvenes ya recogen el plástico pour venderlo o para fabricar objetos útiles: es una buena señal de esperanza.

El medioambiente somos nosotros: cuidarlo es cuidarnos.
\end{textebox}

\textbf{Consigne de lecture :} lis ce texte \textbf{deux fois}. La première fois, cherche uniquement le sens global : de quoi parle-t-on ? La deuxième fois, souligne les mots-clés (noms de déchets, verbes d'action) et les connecteurs.

\begin{definition}[Glossaire du texte]
\begin{itemize}
\item \esp{se acumula} = s'accumule (voix passive réfléchie) ; \esp{sin embargo} = cependant ; \esp{una costumbre} = une habitude.
\item \esp{separan} = trient, séparent ; \esp{el vidrio} = le verre ; \esp{la materia orgánica} = la matière organique.
\item \esp{los camiones} = les camions ; \esp{las plantas de reciclaje} = les usines de recyclage.
\item \esp{salvar} = sauver ; \esp{el suelo} = le sol (terre) ; \esp{cuidarlo} = en prendre soin (litt. « le soigner »).
\item \esp{las cunetas} = les caniveaux, fossés ; \esp{se inundan} = sont inondés (passif réfléchi) ; \esp{los envases} = les emballages.
\item \esp{el cartón} = le carton ; \esp{tirar} = jeter ; \esp{convertirse en} = se transformer en (verbe pronominal).
\item \esp{las cáscaras de fruta} = les épluchures de fruits ; \esp{el abono} = l'engrais, le compost.
\item \esp{los vertederos} = les décharges (contrôlées ou non) ; \esp{una señal de esperanza} = un signe d'espérance.
\end{itemize}
\end{definition}

\textbf{Traduction globale (pour vérification) :} « Dans les grandes villes comme Douala et Yaoundé, les ordures s'accumulent dans les rues. Chaque habitant produit presque un kilo de déchets par jour. Sacs plastiques, bouteilles vides, restes de nourriture : tout finit trop souvent dans les caniveaux ou les rivières, et quand les pluies arrivent, les quartiers sont inondés. Cependant, en Espagne, le tri sélectif est une habitude : les citoyens séparent le plastique, le verre, le papier et la matière organique dans des conteneurs de couleurs. Le conteneur jaune reçoit les emballages, le bleu le papier et le carton, et le vert le verre. Personne ne demande plus où jeter une bouteille : tout le monde le sait depuis l'école. Ensuite, les camions emportent les déchets vers les usines de recyclage, où ils se transforment en nouveaux produits. Une bouteille en plastique peut devenir une veste ; un vieux journal, une boîte en carton ; les épluchures de fruits, de l'engrais pour les champs. Recycler, réutiliser et réduire sont les trois R qui peuvent sauver notre planète. Si nous collaborons tous, nous éviterons que les décharges contaminent le sol et l'eau. À Douala et à Yaoundé, certains jeunes ramassent déjà le plastique pour le vendre ou fabriquer des objets utiles : c'est un bon signe d'espoir. L'environnement, c'est nous : en prendre soin, c'est prendre soin de nous. »

\courssub{Comment lire efficacement : la double lecture}

\begin{methode}[La stratégie de la double lecture]
\begin{enumerate}
\item \textbf{Première lecture (globale) :} lis le titre et le premier paragraphe. Demande-toi : \emph{de quoi parle le texte ?} Repère la source et le thème. Ne t'arrête sur aucun mot inconnu.
\item \textbf{Deuxième lecture (analytique) :} souligne les mots-clés (ici : \esp{basura}, \esp{residuos}, \esp{recogida selectiva}, \esp{contenedores}, \esp{reciclaje}) et encadre les connecteurs. Devine les mots inconnus grâce au contexte et aux mots transparents.
\item \textbf{Vérification :} résume le texte en une phrase en français, puis en une phrase en espagnol. Si tu y arrives, tu as compris.
\end{enumerate}
\end{methode}

\begin{attention}[Le piège de la traduction mot à mot]
Ne traduis \textbf{jamais} mot à mot pendant la lecture : tu perdrais du temps et tu bloquerais sur chaque terme inconnu. Exemple : \esp{los barrios se inundan} — si tu cherches chaque mot, tu perds le fil ; mais le contexte (\esp{cuando llegan las lluvias}, « quand les pluies arrivent ») suffit à deviner : « les quartiers sont inondés ». Au Bac, le correcteur évalue ta compréhension \emph{globale}, pas ta connaissance de chaque mot.
\end{attention}

\courssub{Les connecteurs du texte : le squelette de l'argumentation}

Un commentaire de texte repose sur la capacité à suivre la \textbf{logique} de l'auteur. Or cette logique est portée par les \cle{connecteurs}. Repérons ceux de notre texte :

\begin{center}
\begin{tabularx}{\linewidth}{|X|X|X|}\hline
\rowcolor{popblueL} \textbf{Connecteur} & \textbf{Français} & \textbf{Fonction dans le texte} \\ \hline
\esp{sin embargo} & cependant & opposition : problème au Cameroun / solution en Espagne \\ \hline
\esp{después} & ensuite & étape chronologique de la chaîne de recyclage \\ \hline
\esp{donde} & où & relatif de lieu : situe la transformation (usines) \\ \hline
\esp{si} & si & condition : \esp{Si todos colaboramos...} (subjonctif \esp{contaminen} après) \\ \hline
\esp{cuando} & quand & circonstance temporelle (déclencheur des inondations) \\ \hline
\end{tabularx}
\end{center}

Ces quatre ou cinq petits mots dessinent le \textbf{plan du texte} : constat du problème (\esp{la basura se acumula}), contraste avec l'étranger (\esp{sin embargo}), déroulement du processus (\esp{después}, \esp{donde}), appel à l'action conditionnel (\esp{si}), conclusion slogan (\esp{El medioambiente somos nosotros}). Retiens-les : tu les réutiliseras dans tes propres productions écrites.

\begin{propriete}[La structure argumentative du texte]
\begin{enumerate}
\item \textbf{Constat :} accumulation des déchets dans les villes camerounaises (\esp{la basura se acumula}).
\item \textbf{Contraste :} en Espagne, le tri sélectif est une habitude culturelle (\esp{sin embargo}).
\item \textbf{Processus :} collecte, transport, transformation en nouveaux produits (\esp{después}, \esp{donde}).
\item \textbf{Solution et appel :} les « tres erres » et la condition \esp{si todos colaboramos}.
\item \textbf{Conclusion-chute :} formule frappante, \esp{cuidarlo es cuidarnos}.
\end{enumerate}
\end{propriete}

\courssub{Carte mentale du texte}

\begin{popfigure}
\begin{tikzpicture}[mindmap, grow cyclic, every node/.style={concept, align=center, font=\small},
root concept/.append style={concept color=popblue, font=\small\bfseries},
level 1/.append style={level distance=3.6cm, sibling angle=90},
level 2/.append style={level distance=2.6cm, sibling angle=45}]
\node[root concept, align=center]{El medioambiente\\y nosotros}
  child[concept color=poporange] { node {problema}
    child { node[align=center]{basura en\\las calles} }
    child { node[align=center]{un kilo de\\residuos al día} }
  }
  child[concept color=popgreen] { node {solución}
    child { node[align=center]{recogida\\selectiva} }
    child { node[align=center]{contenedores\\de colores} }
  }
  child[concept color=popgold] { node {proceso}
    child { node {camiones} }
    child { node[align=center]{plantas de\\reciclaje} }
  }
  child[concept color=poppurple] { node {conclusión}
    child { node[align=center]{las tres erres:\\reciclar, reutilizar,\\reducir} }
    child { node[align=center]{cuidarlo es\\cuidarnos} }
  };
\end{tikzpicture}
\legende{Carte mentale : les quatre idées-forces du texte « El medioambiente y nosotros ».}
\end{popfigure}

\begin{savaistu}[Les conteneurs de couleurs en Espagne et en Guinée équatoriale]
En Espagne, le code couleur du tri est une véritable culture nationale : le conteneur \textbf{jaune} (\esp{amarillo}) est réservé au plastique et aux emballages métalliques (\esp{envases}), le \textbf{bleu} (\esp{azul}) au papier et au carton, le \textbf{vert} (\esp{verde}) au verre (\esp{vidrio}). Il existe aussi le conteneur \textbf{marron} (\esp{marrón}) pour la matière organique et le \textbf{gris} (\esp{gris}) pour le reste (\esp{restos}). Les enfants apprennent ce code dès l'école primaire, au point que \esp{¿dónde se tira esto?} (« où jette-t-on ceci ? ») est une question que tout Espagnol résout en une seconde. En Guinée équatoriale, voisine hispanophone du Cameroun, des campagnes de sensibilisation en espagnol (\esp{¡Recicla!}, \esp{Separa tus residuos}) encouragent aussi le tri dans les écoles de Malabo et Bata, souvent avec le même code couleur.
\end{savaistu}

\courssub{Lexique thématique : déchets et recyclage}

\begin{definition}[Vocabulaire essentiel à mémoriser (niveau Bac)]
\begin{center}
\begin{tabularx}{\linewidth}{|>{\bfseries}X|X|X|}\hline
\rowcolor{popblueL} \textbf{Catégorie} & \textbf{Espagnol} & \textbf{Français} \\ \hline
\textbf{Types de déchets} & \esp{la basura} & les ordures (mélangées) \\ \cline{2-3}
& \esp{los residuos} & les déchets (terme technique/général) \\ \cline{2-3}
& \esp{los residuos orgánicos} & les déchets organiques / biodégradables \\ \cline{2-3}
& \esp{el plástico} & le plastique \\ \cline{2-3}
& \esp{el vidrio} & le verre \\ \cline{2-3}
& \esp{el papel y el cartón} & le papier et le carton \\ \cline{2-3}
& \esp{los envases} & les emballages \\ \hline
\textbf{Verbes d'action} & \esp{separar / clasificar} & trier / classer \\ \cline{2-3}
& \esp{reciclar} & recycler \\ \cline{2-3}
& \esp{reutilizar} & réutiliser \\ \cline{2-3}
& \esp{reducir} & réduire \\ \cline{2-3}
& \esp{tirar / echar} & jeter \\ \cline{2-3}
& \esp{recoger} & ramasser, collecter \\ \cline{2-3}
& \esp{transformar (en)} & transformer (en) \\ \cline{2-3}
& \esp{convertirse (en)} & se transformer (en) \\ \hline
\textbf{Noms du processus} & \esp{la recogida selectiva} & la collecte sélective / le tri sélectif \\ \cline{2-3}
& \esp{el contenedor} & le conteneur, la poubelle collective \\ \cline{2-3}
& \esp{la planta de reciclaje} & l'usine de recyclage \\ \cline{2-3}
& \esp{el vertedero} & la décharge (souvent sauvage) \\ \cline{2-3}
& \esp{el relleno sanitario} & la décharge contrôlée / centre d'enfouissement \\ \cline{2-3}
& \esp{el abono / el compost} & l'engrais / le compost \\ \hline
\textbf{Adjectifs / Expressions} & \esp{ecológico / ecológica} & écologique \\ \cline{2-3}
& \esp{sostenible} & durable / soutenable \\ \cline{2-3}
& \esp{la huella ecológica} & l'empreinte écologique \\ \cline{2-3}
& \esp{las tres erres (3R)} & les trois R (Réduire, Réutiliser, Recycler) \\ \hline
\end{tabularx}
\end{center}
\end{definition}

\begin{exemplebox}[Emploi de \esp{basura} vs \esp{residuos}]
\esp{La basura huele mal.} (Les ordures sentent mauvais — concret, quotidien). \\
\esp{La gestión de residuos es un reto.} (La gestion des déchets est un défi — terme administratif, technique).
\end{exemplebox}

\courssub{Méthode du commentaire de texte}

\begin{methode}[Commentaire de texte : les 4 étapes au Bac]
\begin{enumerate}
\item \textbf{Introduction (3-4 lignes) :} Présenter le texte (titre, source, thème, type). Annoncer la problématique (ex: \emph{Comment ce texte montre-t-il que le recyclage est une solution globale aux déchets urbains ?}). Annoncer le plan (2 ou 3 parties).
\item \textbf{Développement (2 parties équilibrées) :} Chaque partie = une idée-force + citations courtes intégrées + analyse (vocabulaire, connecteurs, temps verbaux). Ex: I. Le constat alarmant de la pollution locale. II. Le modèle espagnol et l'appel à l'action citoyenne.
\item \textbf{Ouverture / Expression personnelle (3-4 lignes) :} Lien avec ton contexte (Cameroun), proposition concrète, avis argumenté. Utilise \esp{En mi opinión}, \esp{Creo que}, \esp{Deberíamos} + infinitif / subjonctif.
\item \textbf{Conclusion (1 phrase) :} Réponse brève à la problématique + phrase choc (reprise de la chute du texte ou slogan personnel).
\end{enumerate}
\end{methode}

\begin{attention}[Critères de notation Bac (Compréhension / Expression)]
\begin{itemize}
\item \textbf{Compréhension (6-8 pts) :} Fidélité au texte, citations exactes (guillemets \esp{« »}), reformulation correcte.
\item \textbf{Organisation (3-4 pts) :} Plan visible, connecteurs logiques (\esp{En primer lugar}, \esp{Por un lado... por otro}, \esp{En conclusión}).
\item \textbf{Langue (8-10 pts) :} Correction grammaticale (subjonctif, \esp{ser/estar}, accords), richesse lexicale (vocabulaire du thème), orthographe (accents, \esp{ñ}, \esp{ü}).
\end{itemize}
\end{attention}

\courssub{Commentaire modèle}

\begin{exoresolu}[Commentaire guidé : « El medioambiente y nosotros »]
\textbf{Énoncé.} Commente le texte en suivant la méthode ci-dessus. Longueur visée : 150-180 mots.
\tcblower
\textbf{Proposition de commentaire.}

\textbf{Introducción}
El texto « El medioambiente y nosotros », extraído del manual \emph{Didáctica V}, es un artículo periodístico de sensibilización ecológica. Aborda el problema de la acumulación de residuos en ciudades africanas como Douala y Yaoundé y lo contrasta con el modelo de recogida selectiva en España. La pregunta central es: ¿puede el reciclaje salvar nuestro planeta?

\textbf{Desarrollo}
\textbf{1. Un diagnóstico alarmante: la basura invade las ciudades.}
El primer párrafo describe una realidad cruda: \esp{« la basura se acumula en las calles »} y \esp{« cada habitante produce casi un kilo de residuos al día »}. La enumeración \esp{« Bolsas de plástico, botellas vacías, restos de comida »} ilustra la diversidad de los desechos. La consecuencia inmediata es dramática: \esp{« cuando llegan las lluvias, los barrios se inundan »}. El verbo \esp{inundarse} (pasiva refleja) subraya la vulnerabilidad de la población.

\textbf{2. La solución: la recogida selectiva y las « tres erres ».}
El conector \esp{sin embargo} introduce el contra-modelo español. La costumbre de \esp{separar} en \esp{contenedores de colores} (amarillo, azul, verde) es un hábito escolarizado. El proceso técnico sigue una secuencia lógica marcada por \esp{después} y \esp{donde}: camiones $\rightarrow$ plantas $\rightarrow$ transformación. Ejemplos concretos (\esp{botella $\rightarrow$ chaqueta}, \esp{cáscaras $\rightarrow$ abono}) demuestran la economía circular. La frase final, \esp{« El medioambiente somos nosotros: cuidarlo es cuidarnos »}, usa la identificación sujeto-atributo para hacer responsable a cada ciudadano. La condición \esp{« Si todos colaboramos »} + subjuntivo (\esp{contaminen}) expresa la necesidad de una acción colectiva.

\textbf{Conclusión y apertura}
En definitiva, el texto demuestra que la gestión de residuos no es solo técnica, sino cultural y ciudadana. En Camerún, iniciativas juveniles (\esp{« algunos jóvenes ya recogen el plástico »}) son semillas de esperanza. \textbf{En mi opinión, deberíamos imitar el código de colores en nuestros barrios y exigir más contenedores.} \esp{¡Reciclar es vivir mejor!}
\end{exoresolu}

\courssub{Questions de compréhension et entraînement}

\begin{exoresolu}[Compréhension globale du texte]
\textbf{Énoncé.} Réponds en espagnol, avec des phrases complètes :
\begin{enumerate}
\item ¿Cuántos residuos produce cada habitante al día?
\item ¿Qué separan los ciudadanos españoles y dónde?
\item ¿Qué son las « tres erres »?
\end{enumerate}
\tcblower
\textbf{Solution détaillée.}
\begin{enumerate}
\item \esp{Cada habitante produce casi un kilo de residuos al día.} (On reprend la formulation exacte du texte, premier paragraphe.)
\item \esp{Los ciudadanos españoles separan el plástico, el vidrio, el papel y la materia orgánica en contenedores de colores.} (Deuxième paragraphe : on reformule en gardant les quatre éléments de l'énumération.)
\item \esp{Las tres erres son reciclar, reutilizar y reducir.} (Quatrième paragraphe : réponse courte mais complète, avec le verbe \esp{ser}.)
\end{enumerate}
\textbf{Méthode :} au Bac, une bonne réponse de compréhension reprend les mots du texte sans les copier servilement, et forme toujours une phrase complète avec sujet et verbe conjugué.
\end{exoresolu}

\begin{exercice}[À ton tour — Compréhension et lexique]
\begin{enumerate}
\item Trouve dans le texte trois mots transparents non listés dans le tableau de la page 2 et donne leur équivalent français.
\item Relève les cinq connecteurs du texte et explique, en français, le rôle de chacun dans l'argumentation.
\item Résume le texte en une seule phrase espagnole commençant par \esp{El texto explica que...}
\item \textbf{Version :} Traduis en espagnol : « Au Cameroun, il faut réduire les déchets plastiques pour éviter les inondations. »
\item \textbf{Thème :} Traduis en français : \esp{« Si separamos la basura, los vertederos no contaminarán el suelo. »}
\end{enumerate}
\end{exercice}

\begin{corrige}[Exercice « À ton tour »]
\begin{enumerate}
\item Par exemple : \esp{habitante} (habitant), \esp{ciudadanos} (citoyens), \esp{productos} (produits), \esp{esperanza} (espérance), \esp{inundan} (inondent), \esp{colaboramos} (collaborons).
\item \esp{Sin embargo} introduit l'opposition Cameroun/Espagne ; \esp{después} marque l'étape suivante du processus (le transport) ; \esp{donde} situe le lieu de la transformation (relatif de lieu) ; \esp{si} exprime la condition nécessaire pour sauver la planète (si clause + présent indicatif $\rightarrow$ futur / subjonctif) ; \esp{cuando} situe le moment des inondations (complément circonstanciel de temps).
\item Modèle : \esp{El texto explica que la basura es un gran problema en las ciudades, pero que la recogida selectiva y el reciclaje pueden salvar el planeta si todos colaboramos.}
\item \esp{En Camerún, hay que reducir los residuos plásticos para evitar las inundaciones.} (ou \esp{se deben reducir} / \esp{es necesario reducir}). \emph{Attention :} \esp{evitar} + infinitif ; \esp{inundaciones} (fém. pl.).
\item « Si nous trions les ordures, les décharges ne contamineront pas le sol. » (Futur simple \esp{contaminarán} car phrase principale après \esp{si} + indicatif présent ; \esp{la basura} = les ordures ménagères ; \esp{el suelo} = le sol).
\end{enumerate}
\end{corrige}

\courssub{Production guidée : gestes écocitoyens}

\begin{experience}[Débat oral : « ¿Qué hacemos en nuestro barrio? »]
\textbf{Consigne :} En groupes de 4, préparez une affiche ou une courte intervention orale (1 min 30) pour sensibiliser les élèves de ton lycée au tri des déchets. Utilise \textbf{obligatoirement} :
\begin{itemize}
\item 5 mots du lexique thématique (p. 4).
\item 3 connecteurs logiques (\esp{por eso}, \esp{además}, \esp{por tanto}, \esp{en conclusión}).
\item Une structure \esp{Si + présent indicatif, futur simple} (ex: \esp{Si reciclamos, ahorraremos energía}).
\item Une proposition avec \esp{Deberíamos / Hay que / Es necesario} + infinitif.
\end{itemize}
\textbf{Exemple de début :} \esp{« Compañeros, en nuestro instituto hay mucha basura en el patio. Por eso, deberíamos poner contenedores amarillos y azules. Si separamos el plástico y el papel, evitaremos que terminen en las cunetas. ¡Colaboremos! »}
\end{experience}

\begin{exercice}[Expression écrite : Lettre au maire / Article de blog]
Choisis \textbf{un} des deux sujets. 120-150 mots. Soigne la structure (intro, développement, conclusion) et le vocabulaire du thème.
\begin{enumerate}
\item \textbf{Lettre formelle :} Tu écris au maire de ta commune (Douala, Yaoundé, ou ta ville) pour proposer la mise en place de la \esp{recogida selectiva}. Structure : \esp{Estimado Señor Alcalde:} ... \esp{Le saluda atentamente, [Ton nom]}.
\item \textbf{Article de blog :} Titre : \esp{« Mi barrio, más limpio: ¿cómo lograrlo? »}. Décris la situation actuelle, propose 3 mesures concrètes (conteneurs, éducation, sanctions), termine par un slogan.
\end{enumerate}
\end{exercice}

\begin{corrige}[Pistes de correction — Expression écrite]
\textbf{Sujet 1 (Lettre) — Éléments attendus :}
\begin{itemize}
\item Formules de politesse : \esp{Estimado Sr. Alcalde:} / \esp{Atentamente,}.
\item Vocabulaire : \esp{solicito}, \esp{la instalación de contenedores}, \esp{la concienciación ciudadana}, \esp{los vertederos ilegales}, \esp{la salud pública}.
\item Grammaire : \esp{Le ruego que instale} (subjonctif), \esp{Si el ayuntamiento proporciona contenedores, los vecinos reciclarán} (si + ind. présent, futur), \esp{Es urgente que se tomen medidas} (il est urgent que + subjonctif).
\end{itemize}
\textbf{Sujet 2 (Blog) — Éléments attendus :}
\begin{itemize}
\item Titre accrocheur.
\item Connecteurs : \esp{En primer lugar}, \esp{Además}, \esp{Por último}, \esp{En definitiva}.
\item Mesures : \esp{poner contenedores de colores}, \esp{organizar talleres en las escuelas}, \esp{multar a quien tire basura en la calle}.
\item Slogan final : \esp{« Un barrio limpio, una vida sana »} / \esp{« Reciclar hoy, vivir mañana »}.
\end{itemize}
\end{corrige}

\begin{aretenir}[L'essentiel de la leçon E07]
\begin{itemize}
\item Le texte \esp{El medioambiente y nosotros} (Didáctica V, U2) est un texte argumentatif type Bac : \textbf{problème} $\rightarrow$ \textbf{contraste} $\rightarrow$ \textbf{processus} $\rightarrow$ \textbf{appel à l'action} $\rightarrow$ \textbf{chute}.
\item \textbf{Lexique clé :} \esp{basura / residuos}, \esp{recogida selectiva}, \esp{contenedor (amarillo/azul/verde/marrón)}, \esp{reciclar / reutilizar / reducir (3R)}, \esp{vertedero / relleno sanitario}, \esp{abono / compost}.
\item \textbf{Lecture :} double lecture (globale + analytique) ; mots transparents = alliés ; traduction mot à mot = ennemie.
\item \textbf{Commentaire :} Intro (présentation + problématique + plan) — Développement (2 parties, citations, analyse) — Ouverture personnelle (contexte Cameroun, \esp{Deberíamos} + inf.) — Conclusion (slogan).
\item \textbf{Grammaire repérée :} Passif réfléchi (\esp{se acumula, se inundan, se transforman}), Subjonctif après \esp{evitaremos que / es necesario que}, \esp{Si} + indicatif présent $\rightarrow$ Futur (condition réelle), \esp{convertirse en} (pronominal).
\item \textbf{Contexte camerounais :} Douala/Yaoundé, jeunes ramasseurs de plastique, bendskins, caniveaux bouchés, inondations, espoir par l'initiative locale.
\end{itemize}
\end{aretenir}

\coursec{Compréhension du texte}

Après avoir lu et découvert le texte \esp{« El medioambiente y nosotros »} dans la section précédente, nous allons maintenant vérifier et approfondir notre compréhension. C'est une étape décisive : au baccalauréat, les questions de compréhension portent toujours sur trois niveaux de lecture — la compréhension globale, la compréhension du détail et l'interprétation. Savoir repérer ces trois niveaux, c'est déjà gagner des points.

\courssub{Les trois niveaux de questions de compréhension}

\begin{definition}[Trois niveaux de lecture]
\begin{itemize}
\item La \cle{compréhension globale} : de quoi parle le texte ? Quel est le thème, le problème, la thèse de l'auteur ?
\item La \cle{compréhension du détail} : chiffres, faits précis, actions, lieux mentionnés dans le texte.
\item L'\cle{interprétation} : sens implicite, intention de l'auteur, signification d'une image ou d'une formule.
\end{itemize}
\end{definition}

\begin{methode}[Comment répondre à une question de compréhension]
\begin{enumerate}
\item Lis la question et souligne le \cle{mot interrogatif} (\esp{qué}, \esp{cuántos}, \esp{dónde}, \esp{por qué}...) : il annonce le type de réponse attendu.
\item Repère dans le texte le passage qui contient la réponse.
\item Rédige une \textbf{phrase complète} en reprenant les mots de la question : \esp{¿Qué separan los ciudadanos?} $\rightarrow$ \esp{Los ciudadanos separan...}
\item Justifie par une \textbf{courte citation} entre guillemets (3 à 5 mots maximum) : \esp{El texto dice: « separan sus residuos ».}
\item Vérifie les accords, la conjugaison et les accents avant de rendre ta copie.
\end{enumerate}
\end{methode}

\begin{attention}[Erreur classique au Bac]
Une réponse \textbf{copiée intégralement} du texte, sans reformulation, fait perdre des points : l'examinateur veut vérifier que tu as \emph{compris}, pas que tu sais recopier. La bonne stratégie : une phrase complète avec tes propres mots + une courte citation de 3 à 5 mots entre guillemets comme preuve. Évite aussi de répondre par un seul mot (\esp{Sí}, \esp{Los camiones}) : toute réponse doit être une phrase.
\end{attention}

\courssub{Les mots interrogatifs : petits mots, grands accents}

Avant de répondre, il faut comprendre la question. Observons :

\esp{¿Qué solución propone el texto?} « Quelle solution le texte propose-t-il ? »

\esp{¿Cuántos residuos produce cada habitante?} « Combien de déchets chaque habitant produit-il ? »

\esp{¿Adónde llevan los camiones los residuos?} « Où les camions transportent-ils les déchets ? »

\begin{propriete}[L'accent des mots interrogatifs]
En espagnol, tous les mots interrogatifs (et exclamatifs) portent un \cle{accent écrit}, même si le même mot existe sans accent avec une autre fonction. Cet accent est \textbf{obligatoire}, y compris dans les questions indirectes : \esp{No sé qué hacer.} « Je ne sais pas quoi faire. »
\begin{center}
\begin{tabularx}{\linewidth}{|X|X|X|}\hline
\rowcolor{popblueL} Interrogatif (accentué) & Relatif / conjonction (sans accent) & Traduction \\ \hline
\esp{qué} & \esp{que} & quoi, que / que \\ \hline
\esp{cuántos, -as} & \esp{cuanto} & combien de / autant que \\ \hline
\esp{dónde / adónde} & \esp{donde} & où \\ \hline
\esp{por qué} (2 mots) & \esp{porque} (1 mot) & pourquoi / parce que \\ \hline
\esp{cuál, cuáles} & \esp{cual} & quel, lequel / lequel \\ \hline
\esp{cuándo} & \esp{cuando} & quand / lorsque \\ \hline
\esp{quién, quiénes} & \esp{quien} & qui \\ \hline
\end{tabularx}
\end{center}
Attention au piège : \esp{¿Por qué reciclas?} « Pourquoi recycles-tu ? » mais \esp{Reciclo porque quiero cuidar el planeta.} « Je recycle \emph{parce que} je veux protéger la planète. »
\end{propriete}

\begin{exemplebox}[Questions et réponses]
\esp{¿Qué es la recogida selectiva?} « Qu'est-ce que le tri sélectif ? »

\esp{¿Cuándo pasan los camiones de la basura?} « Quand passent les camions à ordures ? »

\esp{¿Quién separa los residuos en casa?} « Qui trie les déchets à la maison ? »
\end{exemplebox}

\courssub{Questions de compréhension globale et réponses modèles}

Voici les trois questions globales essentielles sur notre texte, avec des réponses modèles rédigées en phrases complètes. Remarque la construction : reprise des mots de la question + courte citation.

\textbf{1.} \esp{¿De qué habla el texto?}

\esp{El texto habla del problema de los residuos y de la importancia del reciclaje para proteger el medioambiente. El autor explica que « el medioambiente somos nosotros ».}

\textbf{2.} \esp{¿Cuál es el problema presentado?}

\esp{El problema presentado es la gran cantidad de basura que producimos cada día y que amenaza la naturaleza si no la tratamos bien. El texto dice que los residuos « contaminan el suelo y el agua ».}

\textbf{3.} \esp{¿Qué solución propone el texto?}

\esp{El texto propone como solución la recogida selectiva y las « tres erres » : reducir, reutilizar y reciclar.}

\begin{exemplebox}[Exemples traduits à retenir]
\esp{La recogida selectiva es una costumbre.} « Le tri sélectif est une habitude. »

\esp{Reciclar, reutilizar y reducir son las tres erres.} « Recycler, réutiliser et réduire sont les trois R. »

\esp{Los ciudadanos separan sus residuos en contenedores de colores.} « Les citoyens trient leurs déchets dans des conteneurs de couleurs. »
\end{exemplebox}

\courssub{Questions de détail : chercher l'information précise}

Les questions de détail exigent de retrouver un chiffre, un fait, un lieu. La réponse doit être exacte mais formulée en phrase complète.

\textbf{1.} \esp{¿Cuántos residuos produce cada habitante al día?}

\esp{Cada habitante produce más de un kilo de residuos al día, según el texto.}

\textbf{2.} \esp{¿Qué separan los ciudadanos españoles?}

\esp{Los ciudadanos españoles separan el papel, el vidrio, los envases de plástico y la materia orgánica, y los depositan en « contenedores de colores ».}

\textbf{3.} \esp{¿Adónde llevan los camiones los residuos?}

\esp{Los camiones llevan los residuos a las plantas de reciclaje, donde se transforman en nuevos productos.}

\begin{attention}[Faux amis et pièges dans les questions]
\esp{¿Adónde?} demande une \emph{destination} (avec un verbe de mouvement : \esp{llevar}, \esp{ir}) ; \esp{¿dónde?} demande un \emph{lieu} (sans mouvement : \esp{estar}). Compare : \esp{¿Dónde está el contenedor?} « Où est le conteneur ? » / \esp{¿Adónde llevas la bolsa?} « Où portes-tu le sac ? »

Autres pièges fréquents : \esp{actualmente} signifie « actuellement, en ce moment » (et non « en fait ») ; \esp{los residuos} est le mot courant pour « les déchets, les ordures » (et non seulement « les résidus » au sens chimique) ; \esp{la recogida} signifie « la collecte, le ramassage » (et non « la récolte », qui se dit \esp{la cosecha}).
\end{attention}

\courssub{Questions d'interprétation : lire entre les lignes}

L'interprétation est le niveau le plus valorisé au Bac : il faut expliquer le sens profond d'une formule.

\textbf{1.} \esp{¿Qué significan « las tres erres »?}

\esp{« Las tres erres » son las palabras reducir, reutilizar y reciclar, que empiezan por la letra erre. Representan los tres gestos fundamentales para consumir menos y proteger la naturaleza.}

\textbf{2.} \esp{¿Por qué dice el autor « el medioambiente somos nosotros »?}

\esp{El autor dice « el medioambiente somos nosotros » porque nuestra salud y nuestra vida dependen de la naturaleza : si ensuciamos el aire, el agua y el suelo, nos hacemos daño a nosotros mismos. Cuidar el planeta es, por lo tanto, cuidarnos.}

\courssub{La chaîne du recyclage : schéma du texte}

Le texte décrit un parcours précis, du citoyen au produit neuf. Visualisons-le :

\begin{popfigure}
\begin{center}
\begin{tikzpicture}[node distance=0.55cm and 0.9cm,
etape/.style={rectangle, rounded corners=4pt, draw=popblue, fill=popblueL, align=center, minimum height=0.9cm, font=\small},
fleche/.style={-{Stealth[length=2.5mm]}, thick, draw=poporange}]
\node[etape, align=center] (a){ciudadanos\\ (separan)};
\node[etape, right=of a, align=center] (b){contenedores\\ de colores};
\node[etape, right=of b, align=center] (c){camiones\\ (recogida)};
\node[etape, right=of c, align=center] (d){plantas de\\ reciclaje};
\node[etape, right=of d, align=center] (e){nuevos\\ productos};
\draw[fleche] (a) -- (b);
\draw[fleche] (b) -- (c);
\draw[fleche] (c) -- (d);
\draw[fleche] (d) -- (e);
\end{tikzpicture}
\end{center}
\legende{La chaîne de la collecte et du recyclage décrite dans le texte : du tri chez soi jusqu'au produit neuf.}
\end{popfigure}

Ce schéma est aussi une aide à l'oral : apprends à le commenter en espagnol — \esp{Primero, los ciudadanos separan los residuos ; después, los camiones los llevan a las plantas de reciclaje ; finalmente, se fabrican nuevos productos.} « D'abord, les citoyens trient les déchets ; ensuite, les camions les transportent vers les usines de recyclage ; enfin, on fabrique de nouveaux produits. »

\courssub{Vrai ou faux justifié : l'exercice type Bac}

\begin{exoresolu}[Vrai / Faux justifié]
Énoncé. Lis les affirmations suivantes. Écris \esp{Verdadero} ou \esp{Falso} et justifie par une courte citation du texte.

1. \esp{Cada habitante produce muy pocos residuos al día.}

2. \esp{Los ciudadanos separan el papel y el vidrio.}

3. \esp{Los residuos van directamente al mar.}

4. \esp{Reciclar es una de las tres erres.}
\tcblower
Solution détaillée.

1. \textbf{Falso.} \esp{El texto dice que cada habitante produce « más de un kilo » de residuos al día : no es poco.} (Faux : le texte dit que chaque habitant produit « plus d'un kilo » de déchets par jour.)

2. \textbf{Verdadero.} \esp{El texto afirma que los ciudadanos « separan el papel, el vidrio » y otros materiales.} (Vrai : citation exacte du texte.)

3. \textbf{Falso.} \esp{Los camiones llevan los residuos « a las plantas de reciclaje », no al mar.} (Faux : les déchets vont aux usines de recyclage.)

4. \textbf{Verdadero.} \esp{El texto dice : « reciclar, reutilizar y reducir son las tres erres ».} (Vrai : recycler fait partie des trois R.)

Méthode : chaque justification tient en une phrase complète + une citation de 3 à 5 mots entre guillemets.
\end{exoresolu}

\courssub{Reformuler : dire autrement la même idée}

La reformulation est la compétence reine du commentaire. Prenons l'idée centrale du texte : \esp{cuidar el medioambiente es cuidarnos} « protéger l'environnement, c'est nous protéger nous-mêmes ». Voici plusieurs reformulations possibles, toutes correctes :

\begin{itemize}
\item \esp{Si protegemos la naturaleza, protegemos nuestra propia vida.} « Si nous protégeons la nature, nous protégeons notre propre vie. »
\item \esp{Nuestra salud depende de la salud del planeta.} « Notre santé dépend de la santé de la planète. »
\item \esp{Al ensuciar el medioambiente, nos hacemos daño a nosotros mismos.} « En salissant l'environnement, nous nous faisons du mal à nous-mêmes. »
\item \esp{El bienestar de las personas está unido al bienestar de la naturaleza.} « Le bien-être des personnes est lié au bien-être de la nature. »
\end{itemize}

\begin{methode}[Technique de la reformulation en trois gestes]
\begin{enumerate}
\item Remplace les mots par des \textbf{synonymes} : \esp{cuidar} $\rightarrow$ \esp{proteger} ; \esp{medioambiente} $\rightarrow$ \esp{naturaleza, planeta}.
\item Change la \textbf{structure} de la phrase : affirmation $\rightarrow$ conditionnelle avec \esp{si}, ou phrase passive $\rightarrow$ active.
\item Garde le \textbf{sens exact} : ne rajoute pas d'idée qui n'est pas dans le texte.
\end{enumerate}
\end{methode}

\courssub{Pour s'entraîner : un texte parallèle}

Pour vérifier que tu sais appliquer la méthode à un texte nouveau, voici un court texte original sur le contexte camerounais, suivi de questions.

\begin{textebox}[La basura en nuestras ciudades]
En Douala y en Yaoundé, la basura es un problema visible : bolsas de plástico en las calles, botellas vacías junto a los mercados, residuos quemados al aire libre. Sin embargo, la situación cambia poco a poco. En varios barrios, asociaciones de jóvenes organizan la recogida de residuos cada sábado por la mañana. Separan el plástico, el metal y la materia orgánica, y venden los materiales reciclables a pequeñas empresas. Con el plástico recuperado, algunos artesanos fabrican macetas, baldosas e incluso muebles escolares. « La basura es un tesoro mal colocado », repite Amina, una joven voluntaria de Douala. Gracias a estas iniciativas, los habitantes comprenden que reducir, reutilizar y reciclar no son palabras extranjeras, sino gestos cotidianos al alcance de todos. El medioambiente de nuestras ciudades depende de nosotros : cada bolsa recogida, cada botella reutilizada es una victoria para la salud de todos.
\end{textebox}

\textbf{Glossaire} : \esp{las bolsas} « les sacs » ; \esp{quemados al aire libre} « brûlés à l'air libre » ; \esp{los barrios} « les quartiers » ; \esp{la recogida} « la collecte » ; \esp{las macetas} « les pots de fleurs » ; \esp{las baldosas} « les carreaux » ; \esp{un tesoro mal colocado} « un trésor mal placé » ; \esp{al alcance de todos} « à la portée de tous ».

\begin{exercice}[À toi de jouer]
Réponds par des phrases complètes, avec une courte citation pour chaque réponse.

1. \esp{¿Qué problema presenta el texto?}

2. \esp{¿Qué organizan las asociaciones de jóvenes?}

3. \esp{¿Qué fabrican los artesanos con el plástico?}

4. \esp{¿Qué significa « la basura es un tesoro mal colocado »?} Reformule l'idée avec tes propres mots, en espagnol.

5. Vrai/Faux : \esp{Las iniciativas de reciclaje no existen en Camerún.}
\end{exercice}

\begin{corrige}[Exercice « La basura en nuestras ciudades »]
1. \esp{El texto presenta el problema de la basura en las calles de Douala y Yaoundé : « bolsas de plástico en las calles ».}

2. \esp{Las asociaciones de jóvenes organizan « la recogida de residuos » cada sábado por la mañana.}

3. \esp{Con el plástico recuperado, los artesanos fabrican « macetas, baldosas e incluso muebles escolares ».}

4. \esp{La frase significa que los residuos tienen valor si los reciclamos : lo que parece basura puede convertirse en un recurso útil.} « La phrase signifie que les déchets ont de la valeur si on les recycle : ce qui semble être un déchet peut devenir une ressource utile. »

5. \textbf{Falso.} \esp{El texto describe « estas iniciativas » de jóvenes y artesanos en Douala y Yaoundé.}
\end{corrige}

\begin{aretenir}[L'essentiel de la section]
\begin{itemize}
\item Trois niveaux de questions : globale, détail, interprétation — repère-les grâce au mot interrogatif.
\item Réponse complète = reprise de la question + reformulation + citation courte (3 à 5 mots) entre guillemets.
\item Les mots interrogatifs portent toujours un accent : \esp{qué, cuántos, dónde, por qué, cuál}.
\item Reformuler, c'est changer les mots et la structure sans changer le sens : \esp{cuidar el medioambiente es cuidarnos} $\rightarrow$ \esp{nuestra salud depende de la salud del planeta}.
\end{itemize}
\end{aretenir}

\coursec{Lexique thématique : déchets et recyclage}

Dans la section précédente, nous avons lu et compris le texte \esp{El medioambiente y nosotros}. Pour aller plus loin et pouvoir \emph{commenter} ce texte avec précision, il nous faut maintenant maîtriser son vocabulaire : celui des déchets, de leur collecte, de leur traitement et des gestes écocitoyens. Un bon commentaire de texte, en Terminale A, se reconnaît d'abord à la richesse et à l'exactitude du lexique mobilisé. Observons donc ce vocabulaire en action, dans un texte original, avant de l'étudier champ par champ.

\begin{textebox}[Un barrio que aprende a reciclar]
En el barrio de Briqueterie, en Yaoundé, la basura era un problema diario: las bolsas de plástico se acumulaban en las calles y, cuando llegaban las lluvias, los residuos bloqueaban las canaletas. Una asociación de jóvenes decidió actuar. Primero, colocaron cubos de basura delante de cada casa y explicaron a los vecinos la importancia de separar los residuos: la materia orgánica a un lado, el plástico y las latas al otro.

Después, organizaron una recogida selectiva todos los sábados por la mañana. Los niños recogían el papel y el cartón, que vendían a una pequeña planta de reciclaje instalada cerca del mercado. Las botellas de vidrio se reutilizaban para conservar la miel y el aceite de palma. La materia orgánica, en lugar de terminar en un vertedero, se compostaba detrás de la escuela para producir un abono natural que los agricultores usaban en sus campos de maíz.

« Antes tirábamos todo al mismo sitio », explica Mama Ngo Bell, una vendedora del mercado. « Ahora sabemos que reducir, reutilizar y reciclar son las tres erres que cuidan nuestro barrio ». Gracias a estos gestos sencillos, el barrio está más limpio, el agua está menos contaminada y los jóvenes han comprendido que proteger el medioambiente empieza delante de su propia puerta. El ejemplo de Briqueterie demuestra que no hace falta esperar grandes medios para evitar la contaminación: basta con voluntad, organización y un poco de educación ambiental.
\end{textebox}

\textbf{Glossaire :} \esp{el barrio} : le quartier ; \esp{las bolsas} : les sacs ; \esp{se acumulaban} : s'accumulaient ; \esp{las canaletas} : les caniveaux ; \esp{colocaron} : ils ont placé ; \esp{el cartón} : le carton ; \esp{el abono} : l'engrais ; \esp{los campos de maíz} : les champs de maïs ; \esp{la puerta} : la porte ; \esp{basta con} : il suffit de.

\textbf{Questions de vérification lexicale :}
\begin{enumerate}
\item ¿Qué residuos separan los vecinos de Briqueterie ?
\item ¿Qué hacen con la materia orgánica ?
\item ¿Cuáles son las « tres erres » mencionadas en el texto ?
\end{enumerate}

\courssub{Champ 1 : les déchets (los residuos)}

\begin{definition}[Les noms des déchets]
\esp{La basura} désigne les \cle{ordures ménagères}, ce qu'on jette à la maison. \esp{Los residuos} est le terme plus général et plus technique : les \cle{déchets} (ménagers, industriels, hospitaliers). \esp{El desecho} évoque le \cle{rebut}, ce qui est mis au rebut. On y ajoute les matières : \esp{la materia orgánica} (la matière organique), \esp{el plástico} (le plastique), \esp{el vidrio} (le verre), \esp{el papel} (le papier), \esp{el cartón} (le carton), \esp{la lata} (la canette).
\end{definition}

\begin{exemplebox}[Les déchets en contexte]
\esp{La basura de la cocina es materia orgánica.} « Les ordures de la cuisine sont de la matière organique. »

\esp{Los residuos industriales son muy peligrosos.} « Les déchets industriels sont très dangereux. »

\esp{En el mercado de Douala hay latas y botellas de vidrio por todas partes.} « Au marché de Douala, il y a des canettes et des bouteilles en verre partout. »
\end{exemplebox}

\begin{attention}[basura ou residuos ?]
\esp{Basura} se dit des ordures ménagères, du quotidien : \esp{sacar la basura} (sortir la poubelle). \esp{Residuos} est le terme technique des discours officiels, des lois et du tri : \esp{residuos tóxicos}, \esp{gestión de residuos}. Au singulier, \esp{un residuo} désigne souvent un résidu chimique : ne traduis pas « les déchets » par \esp{las basuras} (pluriel rare et incorrect ici), mais par \esp{los residuos} ou \esp{la basura}.
\end{attention}

\begin{attention}[el vidrio, el cristal, el vaso]
\esp{El vidrio} = le verre, la \emph{matière} : \esp{una botella de vidrio} (une bouteille en verre). \esp{El cristal} = le cristal ou la vitre. Le verre \emph{à boire} se dit \esp{el vaso} : \esp{un vaso de agua} (un verre d'eau). Ne dis jamais \esp{un vidrio de agua} !
\end{attention}

\courssub{Champ 2 : la collecte (la recogida)}

\begin{definition}[recogida selectiva]
La \esp{recogida selectiva} est la \cle{collecte sélective}, c'est-à-dire le tri des déchets par catégorie (plastique, papier, verre, matière organique) avant leur ramassage, afin de permettre leur recyclage.
\end{definition}

Vocabulaire de la collecte : \esp{el contenedor} (le conteneur, la grande benne publique), \esp{el cubo de basura} (la poubelle, à la maison), \esp{el camión de basura} (le camion à ordures). Les verbes essentiels : \esp{recoger} (ramasser, collecter), \esp{tirar} (jeter), \esp{echar} (jeter, synonyme très fréquent : \esp{echar la basura}), \esp{separar} (trier, séparer).

\begin{propriete}[Genre et pluriel : el contenedor, la recogida]
\begin{itemize}
\item Les noms en \esp{-dor / -dora} sont souvent masculins quand ils désignent un objet ou une machine : \esp{el contenedor}, pluriel \esp{los contenedores} (les noms se terminant par une consonne forment leur pluriel en \esp{-es}).
\item \esp{La recogida} est un \cle{nom déverbal} féminin formé sur le verbe \esp{recoger} : il désigne l'\emph{action} de ramasser. De même : \esp{reciclar} $\rightarrow$ \esp{el reciclaje} ; \esp{reducir} $\rightarrow$ \esp{la reducción} ; \esp{reutilizar} $\rightarrow$ \esp{la reutilización}.
\end{itemize}
\end{propriete}

\begin{exemplebox}[La collecte en phrases]
\esp{El camión de basura pasa los lunes por la mañana.} « Le camion à ordures passe le lundi matin. »

\esp{Hay que separar los residuos antes de tirarlos.} « Il faut trier les déchets avant de les jeter. »

\esp{Los alumnos recogen el papel de las aulas cada viernes.} « Les élèves ramassent le papier des salles de classe chaque vendredi. »
\end{exemplebox}

\begin{attention}[tirar : jeter, pas « tirer »]
Faux ami classique : \esp{tirar la basura} signifie \emph{jeter} les ordures, et non « tirer ». Pour « tirer » (une porte, une corde), l'espagnol utilise \esp{tirar de} ou \esp{jalar} (Amérique). Expression à retenir : \esp{tirar la basura al contenedor} = « jeter les ordures dans le conteneur ». Note aussi la tournure impersonnelle \esp{hay que + infinitif} = « il faut » : \esp{Hay que separar los residuos.}
\end{attention}

\courssub{Champ 3 : le traitement (el tratamiento)}

Une fois collectés, les déchets sont traités. Trois destins possibles : le \cle{recyclage} (\esp{el reciclaje}, verbe \esp{reciclar}) dans une \esp{planta de reciclaje} (usine de recyclage) ; l'enfouissement dans un \esp{vertedero} (décharge) ; l'incinération (\esp{incinerar}). Enfin, la matière organique peut être \esp{compostada} (\esp{compostar} : composter) pour produire de l'engrais.

\begin{exemplebox}[Exemples traduits]
\esp{El vertedero contamina el agua.} « La décharge pollue l'eau. »

\esp{En esta planta de reciclaje se reciclan las latas de aluminio.} « Dans cette usine de recyclage, on recycle les canettes en aluminium. »

\esp{Incinerar los residuos produce humo tóxico.} « Incinérer les déchets produit de la fumée toxique. »

\esp{Compostamos la materia orgánica para el huerto del colegio.} « Nous compostons la matière organique pour le potager du collège. »
\end{exemplebox}

\courssub{Champ 4 : les gestes écocitoyens}

\begin{definition}[Las tres erres]
\esp{Las tres erres} (les trois R) résument la conduite écocitoyenne : \cle{reducir} (réduire sa consommation), \cle{reutilizar} (réutiliser les objets), \cle{reciclar} (recycler les matériaux). On y ajoute : \esp{ahorrar} (économiser : eau, énergie), \esp{cuidar} (prendre soin de), \esp{proteger} (protéger), \esp{evitar} (éviter) et leur contraire \esp{contaminar} (polluer).
\end{definition}

\begin{exemplebox}[Des gestes en espagnol]
\esp{Debemos reducir el consumo de plástico.} « Nous devons réduire la consommation de plastique. »

\esp{Reutilizo las bolsas cuando voy al mercado.} « Je réutilise les sacs quand je vais au marché. »

\esp{Hay que ahorrar agua y cuidar el medioambiente.} « Il faut économiser l'eau et prendre soin de l'environnement. »

\esp{Evitemos contaminar los ríos.} « Évitons de polluer les rivières. »
\end{exemplebox}

\begin{attention}[cuidar, proteger : constructions et conjugaison]
\esp{Cuidar} se construit sans préposition devant un nom quand il signifie « prendre soin de » : \esp{cuidar el medioambiente} (et non \esp{cuidar de el medioambiente} dans ce sens ; \esp{cuidar de} existe, mais au sens de « s'occuper de quelqu'un » : \esp{cuida de sus hijos}).

\esp{Proteger} présente une alternance consonantique \esp{g $\rightarrow$ j} devant \esp{a} et \esp{o}, pour conserver le son : présent de l'indicatif \esp{protejo, proteges, protege, protegemos, protegéis, protegen} ; présent du subjonctif \esp{proteja, protejas, proteja, protejamos, protejáis, protejan}. Écris \esp{yo protejo el bosque}, jamais \esp{protego}. Même phénomène avec \esp{recoger} : \esp{yo recojo}, \esp{recoja}.
\end{attention}

\courssub{Tableau récapitulatif du lexique obligatoire}

\begin{center}
\begin{tabularx}{\linewidth}{|l|X|X|X|}
\hline
\rowcolor{popblueL} Champ lexical & Español & Français & Ejemplo \\
\hline
Déchets & la basura & les ordures & \esp{Saco la basura cada noche.} \\
\hline
Déchets & los residuos & les déchets & \esp{Separamos los residuos en casa.} \\
\hline
Collecte & el contenedor & le conteneur & \esp{Tira la botella al contenedor verde.} \\
\hline
Collecte & la recogida selectiva & le tri sélectif & \esp{La recogida selectiva es los sábados.} \\
\hline
Traitement & el reciclaje & le recyclage & \esp{El reciclaje salva recursos.} \\
\hline
Traitement & el vertedero & la décharge & \esp{El vertedero contamina el agua.} \\
\hline
Gestes & reutilizar & réutiliser & \esp{Reutilizamos los frascos de vidrio.} \\
\hline
Gestes & reducir & réduire & \esp{Debemos reducir el consumo de plástico.} \\
\hline
\end{tabularx}
\end{center}

\courssub{Les conteneurs de couleurs : un code à connaître}

En Espagne et dans de nombreux pays hispanophones, chaque conteneur a une couleur qui indique la matière à y jeter. Ce code est très utile pour le commentaire de texte et pour la production orale.

\begin{popfigure}
\begin{center}
\begin{tikzpicture}[font=\small]
\node[draw=popgold, fill=popgold!35, rounded corners, minimum width=2.6cm, minimum height=2.2cm, align=center] (a) at (0,0) {\textbf{Amarillo}\\ plástico, latas};
\node[draw=popblue, fill=popblue!25, rounded corners, minimum width=2.6cm, minimum height=2.2cm, align=center] (b) at (3.6,0) {\textbf{Azul}\\ papel, cartón};
\node[draw=popgreen, fill=popgreen!25, rounded corners, minimum width=2.6cm, minimum height=2.2cm, align=center] (c) at (7.2,0) {\textbf{Verde}\\ vidrio};
\node[draw=popdark, fill=poporange!30, rounded corners, minimum width=2.6cm, minimum height=2.2cm, align=center] (d) at (10.8,0) {\textbf{Marrón}\\ materia\\ orgánica};
\node[above=0.15cm of a, popink] {\esp{contenedor}};
\node[above=0.15cm of b, popink] {\esp{contenedor}};
\node[above=0.15cm of c, popink] {\esp{contenedor}};
\node[above=0.15cm of d, popink] {\esp{contenedor}};
\end{tikzpicture}
\end{center}
\legende{Le code couleur des conteneurs : jaune = plastique et canettes ; bleu = papier et carton ; vert = verre ; marron = matière organique.}
\end{popfigure}

\begin{exemplebox}[Utiliser le code couleur]
\esp{Las botellas de vidrio van al contenedor verde.} « Les bouteilles en verre vont dans le conteneur vert. »

\esp{El periódico viejo se tira al contenedor azul.} « Le vieux journal se jette dans le conteneur bleu. »
\end{exemplebox}

\courssub{Carte mentale du champ lexical}

\begin{popfigure}
\begin{center}
\begin{tikzpicture}[font=\small, level 1/.style={sibling distance=3.4cm, level distance=1.9cm}, every node/.style={align=center}]
\node[draw=popink, fill=popblueL, rounded corners, font=\bfseries] {LOS RESIDUOS}
child { node[draw=popblue, fill=popblue!15, rounded corners, align=center]{la basura\\ el desecho} }
child { node[draw=popgreen, fill=popgreen!15, rounded corners, align=center]{la recogida\\ el contenedor\\ el camión} }
child { node[draw=poporange, fill=poporange!20, rounded corners, align=center]{el reciclaje\\ el vertedero\\ compostar} }
child { node[draw=poppurple, fill=poppurple!15, rounded corners, align=center]{reducir\\ reutilizar\\ reciclar} };
\end{tikzpicture}
\end{center}
\legende{Carte mentale : les quatre champs lexicaux autour du mot central \esp{los residuos}.}
\end{popfigure}

\begin{exoresolu}[Compléter avec le mot juste]
Énoncé : complète chaque phrase avec le mot qui convient : \esp{basura, contenedor, reciclaje, vertedero, reutilizar, recogida}.

1. \esp{El camión pasa el lunes para la \ldots\ de la basura.}
2. \esp{Las latas se tiran al \ldots\ amarillo.}
3. \esp{El \ldots\ de la ciudad está cerca del río y contamina el agua.}
4. \esp{Debemos \ldots\ las botellas de vidrio en casa.}
5. \esp{El \ldots\ del papel ahorra muchos árboles.}
6. \esp{Saqué la \ldots\ después de cenar.}

\tcblower
Solution détaillée :

1. \esp{recogida} : le camion assure la \emph{collecte} (nom déverbal féminin de \esp{recoger}).
2. \esp{contenedor} : le conteneur jaune reçoit le plastique et les canettes.
3. \esp{vertedero} : la décharge pollue l'eau du fleuve.
4. \esp{reutilizar} : infinitif après \esp{debemos} ; on réutilise les bouteilles.
5. \esp{reciclaje} : nom masculin en \esp{-aje} formé sur \esp{reciclar}.
6. \esp{basura} : expression figée \esp{sacar la basura} (sortir la poubelle).
\end{exoresolu}

\begin{savaistu}[La recogida de la basura au Cameroun]
À Yaoundé et à Douala, c'est l'entreprise HYSACAM qui assure la \esp{recogida de la basura} : ses camions verts sont familiers de tous les habitants. Mais le \esp{reciclaje} et la \esp{recogida selectiva} restent peu développés au Cameroun : la plupart des déchets finissent dans des \esp{vertederos}. Ce contraste avec l'Espagne, où chaque quartier dispose de conteneurs de couleurs, est un excellent argument pour ton commentaire de texte : il te permet de comparer les réalités hispanophone et camerounaise, compétence très valorisée à l'examen.
\end{savaistu}

\begin{aretenir}[L'essentiel du lexique]
\begin{itemize}
\item \esp{Basura} = ordures ménagères ; \esp{residuos} = déchets (terme technique, employé au pluriel dans ce sens).
\item \esp{El contenedor} (pluriel \esp{los contenedores}) est la benne publique ; \esp{el cubo de basura} est la poubelle de la maison.
\item \esp{La recogida} est le nom déverbal de \esp{recoger} ; \esp{el reciclaje} celui de \esp{reciclar}.
\item Le code couleur : amarillo (plástico), azul (papel), verde (vidrio), marrón (materia orgánica).
\item Les trois erres : \esp{reducir, reutilizar, reciclar} ; le contraire de \esp{proteger} est \esp{contaminar}.
\item Alternance \esp{g $\rightarrow$ j} : \esp{yo protejo}, \esp{yo recojo}.
\end{itemize}
\end{aretenir}

\begin{exercice}[À toi de jouer]
1. Traduis en espagnol : « Nous devons trier les déchets et jeter le verre dans le conteneur vert. »
2. Trouve le nom déverbal correspondant : \esp{recoger}, \esp{reciclar}, \esp{reducir}, \esp{reutilizar}.
3. Rédige trois phrases proposant des gestes écocitoyens pour ton lycée, avec \esp{reducir}, \esp{ahorrar} et \esp{evitar}.
\end{exercice}

\begin{corrige}[Exercice « À toi de jouer »]
1. \esp{Debemos separar los residuos y tirar el vidrio al contenedor verde.}
2. \esp{recoger $\rightarrow$ la recogida ; reciclar $\rightarrow$ el reciclaje ; reducir $\rightarrow$ la reducción ; reutilizar $\rightarrow$ la reutilización.}
3. Modèle : \esp{Debemos reducir el uso de plástico en el colegio. Hay que ahorrar agua en los baños. Evitemos tirar papeles al suelo del patio.}
\end{corrige}

\coursec{Méthode du commentaire de texte}

Dans la section précédente, nous avons constitué notre boîte à outils lexicale : \esp{basura}, \esp{residuos}, \esp{contenedor}, \esp{recogida selectiva}, \esp{reciclaje}, \esp{vertedero}, \esp{reutilizar}, \esp{reducir}. Mais au baccalauréat, connaître le vocabulaire ne suffit pas : il faut savoir \cle{organiser} sa pensée pour rendre compte d'un texte de manière claire, structurée et personnelle. C'est l'objet du \cle{comentario de texto}, l'exercice roi de la Terminale A. Cette section vous donne la méthode complète, étape par étape, avec les phrases témoins à mémoriser.

\courssub{Qu'est-ce qu'un commentaire de texte ?}

\begin{definition}[Le comentario de texto]
Le \cle{commentaire de texte} (\esp{comentario de texto}) est un exercice écrit en espagnol qui consiste à : (1) \cle{présenter} le texte (source, thème), (2) \cle{résumer et analyser} ses idées principales et ses procédés d'expression, (3) \cle{conclure} par un bilan et une opinion personnelle justifiée. Il se rédige entièrement en espagnol, en trois grandes parties : \esp{introducción}, \esp{desarrollo}, \esp{conclusión}.
\end{definition}

Observez d'abord ce mini-extrait de commentaire rédigé par un élève de Terminale, avant d'en dégager les règles :

\begin{textebox}[Début d'un commentaire modèle]
Este texto, extraído de una revista de divulgación, trata de la contaminación causada por los residuos en las grandes ciudades. El autor denuncia la falta de recogida selectiva y propone soluciones sencillas. En primer lugar, explica que la basura se acumula en las calles; sin embargo, afirma que reciclar es posible si cada ciudadano colabora. Según el autor, reciclar es un deber de todos.
\end{textebox}

\textbf{Glossaire :} \esp{divulgación} = vulgarisation ; \esp{falta} = manque ; \esp{se acumula} = s'accumule ; \esp{ciudadano} = citoyen ; \esp{colaborar} = collaborer.

Qu'observe-t-on ? Le texte est \cle{présenté} (source + thème), l'idée principale est \cle{résumée} sans recopie, les procédés sont \cle{signalés} (\esp{sin embargo}), et le tout est relié par des \cle{connecteurs} (\esp{en primer lugar}, \esp{según el autor}). C'est exactement ce que l'examinateur attend.

\courssub{Étape 1 — L'introduction : présenter le texte}

L'introduction est courte (3 à 5 lignes) mais décisive : elle donne la première impression. Elle doit contenir trois éléments, dans cet ordre :

\begin{enumerate}
\item \cle{La source} : d'où vient le texte (manuel, presse, revue, discours...).
\item \cle{Le thème} : de quoi parle le texte.
\item \cle{La problématique} : quelle question le texte pose-t-il ? (\esp{¿Cuál es la idea principal?})
\end{enumerate}

\begin{propriete}[Phrases témoins pour l'introduction]
\begin{itemize}
\item \esp{Este texto, extraído de..., trata de...} « Ce texte, extrait de..., traite de... »
\item \esp{El texto presenta / describe / analiza...} « Le texte présente / décrit / analyse... »
\item \esp{El autor denuncia / defiende / propone...} « L'auteur dénonce / défend / propose... »
\item \esp{El tema principal es...} « Le thème principal est... »
\item \esp{La idea principal del texto es que...} « L'idée principale du texte est que... »
\end{itemize}
\end{propriete}

\begin{exemplebox}[Exemples traduits]
\esp{Este texto trata de la contaminación causada por los residuos.} « Ce texte traite de la pollution causée par les déchets. »

\esp{El autor denuncia la falta de contenedores en los barrios populares.} « L'auteur dénonce le manque de conteneurs dans les quartiers populaires. »

\esp{El tema principal es la recogida selectiva en las ciudades africanas.} « Le thème principal est la collecte sélective dans les villes africaines. »
\end{exemplebox}

\begin{attention}[Ne pas confondre « tratar de » et « tratar »]
\esp{Tratar de} + nom signifie « traiter de, avoir pour sujet » : \esp{El texto trata de la contaminación.} « Le texte traite de la pollution. » Ne dites jamais \esp{el texto trata la contaminación} sans \esp{de}.

Autre piège : \esp{tratar de} + \textbf{infinitif} signifie « essayer de » : \esp{Trata de entender.} « Essaie de comprendre. » Enfin, \esp{tratar} seul signifie « traiter quelqu'un » : \esp{El médico trata al enfermo.} « Le médecin traite le malade. » Trois constructions, trois sens : vérifiez toujours la présence du \esp{de} et la nature du mot qui suit.
\end{attention}

\courssub{Étape 2 — Le développement : résumer, analyser, citer}

Le développement est le cœur du commentaire. Il comporte \cle{deux ou trois parties}, chacune annoncée par un connecteur. On y fait trois choses :

\begin{enumerate}
\item \cle{Résumer} les idées principales du texte, avec ses propres mots (jamais de recopie intégrale).
\item \cle{Analyser} les procédés : le lexique thématique (champ lexical des \esp{residuos}), les connecteurs (\esp{sin embargo}, \esp{por eso}, \esp{además}), les chiffres éventuels, le ton (dénonciation, espoir, appel à l'action).
\item \cle{Dégager le point de vue de l'auteur} : est-il pessimiste, optimiste, critique, militant ?
\end{enumerate}

\begin{propriete}[Phrases témoins pour le développement]
\begin{itemize}
\item \esp{En primer lugar... / En segundo lugar... / Por último...} « En premier lieu... / En second lieu... / Enfin... »
\item \esp{Por un lado... por otro lado...} « D'un côté... de l'autre côté... »
\item \esp{Según el autor...} « Selon l'auteur... »
\item \esp{El texto subraya que...} « Le texte souligne que... »
\item \esp{El autor afirma / explica / muestra que...} « L'auteur affirme / explique / montre que... »
\item \esp{Sin embargo, el autor matiza...} « Cependant, l'auteur nuance... »
\item \esp{Como dice el texto: « ... »} « Comme le dit le texte : « ... » » (citation brève)
\end{itemize}
\end{propriete}

\begin{exemplebox}[Exemples traduits]
\esp{Según el autor, reciclar es un deber de todos.} « Selon l'auteur, recycler est un devoir de tous. »

\esp{Por un lado, describe el problema de los vertederos; por otro lado, propone la reutilización.} « D'un côté, il décrit le problème des décharges ; de l'autre, il propose la réutilisation. »

\esp{El texto subraya que reducir es más importante que reciclar.} « Le texte souligne que réduire est plus important que recycler. »

\esp{Sin embargo, el autor no es pesimista: cree que los jóvenes pueden cambiar la situación.} « Cependant, l'auteur n'est pas pessimiste : il croit que les jeunes peuvent changer la situation. »
\end{exemplebox}

\begin{attention}[« Según » : pas d'article partitif, pas de « de »]
En espagnol, \esp{según} est suivi directement du nom, \textbf{sans} préposition : \esp{según el autor}, \esp{según el texto}, \esp{según los expertos}. Ne traduisez jamais « selon » par \esp{según de} — ce \esp{de} est un calque du français qui n'existe pas en espagnol. De même, on dit \esp{según él} (et non \esp{según de él}).
\end{attention}

\begin{attention}[Citer brièvement, sans recopier]
Une citation doit être \cle{courte} (une phrase maximum) et \cle{introduite} : \esp{Como afirma el autor: « la basura no espera ».} Recopier trois lignes du texte est sanctionné : cela montre que vous ne savez pas résumer. Privilégiez la paraphrase : \esp{El autor insiste en que debemos actuar ya.} « L'auteur insiste sur le fait que nous devons agir maintenant. » Remarquez la construction : \esp{insistir en que} + indicatif quand on rapporte un fait affirmé.
\end{attention}

\courssub{Étape 3 — La conclusion : bilan et ouverture personnelle}

La conclusion fait le \cle{bilan} en une ou deux phrases, puis s'ouvre sur votre \cle{opinion personnelle}, souvent par une comparaison avec le Cameroun. C'est le seul moment où « je » est attendu.

\begin{propriete}[Phrases témoins pour la conclusion]
\begin{itemize}
\item \esp{En conclusión... / Para concluir... / En resumen...} « En conclusion... / Pour conclure... / En résumé... »
\item \esp{En mi opinión... / A mi juicio... / Pienso que... / Creo que...} « À mon avis... / Je pense que... »
\item \esp{A mi juicio, en Camerún deberíamos...} « À mon avis, au Cameroun nous devrions... »
\item \esp{Este texto me ha hecho reflexionar sobre...} « Ce texte m'a fait réfléchir sur... »
\end{itemize}
\end{propriete}

\begin{propriete}[Opinion : indicatif ou subjonctif ?]
Après \esp{en mi opinión}, \esp{pienso que}, \esp{creo que} (affirmation d'une croyance), on emploie l'\cle{indicatif} : \esp{Pienso que el reciclaje es necesario.} « Je pense que le recyclage est nécessaire. »

Mais à la forme \cle{négative}, le doute s'installe et le \cle{subjonctif} est exigé : \esp{No pienso que sea fácil.} « Je ne pense pas que ce soit facile. » \esp{No creo que el problema tenga una solución rápida.} « Je ne crois pas que le problème ait une solution rapide. »

Rappel de formation : \esp{ser} $\rightarrow$ \esp{sea} ; \esp{tener} $\rightarrow$ \esp{tenga} ; \esp{haber} $\rightarrow$ \esp{haya} ; \esp{poder} $\rightarrow$ \esp{pueda}.
\end{propriete}

\begin{exemplebox}[Exemples traduits]
\esp{En conclusión, el texto nos invita a reducir, reutilizar y reciclar.} « En conclusion, le texte nous invite à réduire, réutiliser et recycler. »

\esp{En mi opinión, en Yaoundé y Douala deberíamos instalar más contenedores.} « À mon avis, à Yaoundé et à Douala, nous devrions installer plus de conteneurs. »

\esp{No pienso que la situación sea desesperada, pero hay que actuar.} « Je ne pense pas que la situation soit désespérée, mais il faut agir. »
\end{exemplebox}

\courssub{Le schéma en entonnoir : du général au personnel}

Le commentaire se rétrécit comme un entonnoir : on part du texte (général) pour arriver à soi (personnel).

\begin{popfigure}
\begin{center}
\begin{tikzpicture}[font=\small]
\node[draw=popblue, fill=popblueL, rounded corners, align=center, text width=10.5cm, minimum height=1.2cm] (intro) at (0,0) {\textbf{INTRODUCCIÓN}\\ Fuente + tema + problemática\\ \esp{Este texto trata de... El autor denuncia...}};
\node[draw=poporange, fill=poporangeL, rounded corners, align=center, text width=8.5cm, minimum height=1.6cm] (des) at (0,-2.3) {\textbf{DESARROLLO (2-3 partes)}\\ Resumen de ideas + análisis del léxico y los conectores + citas breves\\ \esp{En primer lugar... Según el autor... Sin embargo...}};
\node[draw=popgreen, fill=popgreenL, rounded corners, align=center, text width=6.5cm, minimum height=1.2cm] (conc) at (0,-4.6) {\textbf{CONCLUSIÓN}\\ Balance + opinión personal (Camerún)\\ \esp{En mi opinión, deberíamos...}};
\draw[-{Stealth[length=3mm]}, very thick, popdark] (intro) -- (des);
\draw[-{Stealth[length=3mm]}, very thick, popdark] (des) -- (conc);
\end{tikzpicture}
\end{center}
\legende{La structure en entonnoir du commentaire de texte : du texte vers soi.}
\end{popfigure}

\courssub{Plan type du commentaire au Bac}

\begin{methode}[Plan type du comentario de texto]
\begin{enumerate}
\item \textbf{I. Presentación del texto} : source, thème, problématique. Rédiger 3 à 5 lignes avec \esp{Este texto trata de... / El tema principal es...}.
\item \textbf{II. Ideas principales y análisis} : deux ou trois paragraphes. Pour chaque partie : annoncer avec un connecteur (\esp{en primer lugar...}), résumer l'idée, citer brièvement, analyser un procédé (lexique, chiffre, \esp{sin embargo}), dégager le point de vue de l'auteur.
\item \textbf{III. Conclusión y opinión personal} : bilan en une phrase (\esp{en conclusión...}), puis avis justifié (\esp{en mi opinión... / a mi juicio, en Camerún deberíamos...}), avec indicatif après \esp{pienso que} et subjonctif après \esp{no pienso que}.
\item \textbf{Relecture} : vérifier les accents, les connecteurs, l'absence de recopie, et que l'avis personnel n'apparaît qu'à la fin.
\end{enumerate}
\end{methode}

\courssub{Tableau récapitulatif des phrases témoins}

\begin{center}
\begin{tabularx}{\linewidth}{|X|X|X|}
\hline
\rowcolor{popblueL} \textbf{Étape} & \textbf{Español} & \textbf{Français} \\
\hline
Introduction & Este texto, extraído de..., trata de... & Ce texte, extrait de..., traite de... \\
\hline
Introduction & El autor denuncia / propone... & L'auteur dénonce / propose... \\
\hline
Développement & En primer lugar... / Por un lado... por otro lado... & En premier lieu... / D'un côté... de l'autre... \\
\hline
Développement & Según el autor... / El texto subraya que... & Selon l'auteur... / Le texte souligne que... \\
\hline
Développement & Sin embargo, el autor matiza... & Cependant, l'auteur nuance... \\
\hline
Conclusion & En conclusión... / En resumen... & En conclusion... / En résumé... \\
\hline
Conclusion & En mi opinión... / A mi juicio, en Camerún deberíamos... & À mon avis... / Au Cameroun, nous devrions... \\
\hline
Opinion nuancée & No pienso que sea fácil (subjonctif) & Je ne pense pas que ce soit facile \\
\hline
\end{tabularx}
\end{center}

\courssub{Texte d'entraînement et commentaire guidé}

\begin{textebox}[Un barrio que aprendió a reciclar]
En el barrio de Briqueterie, en Yaoundé, los vecinos han decidido cambiar sus hábitos. Antes, la basura se acumulaba en las calles y los vertederos improvisados contaminaban el agua de los pozos. Hoy, gracias a una asociación de jóvenes, cada familia separa sus residuos: el plástico va a un contenedor amarillo, el vidrio a uno verde y la materia orgánica sirve para hacer compost.

Según la presidenta de la asociación, la recogida selectiva ha transformado el barrio. « Antes nadie pensaba en reutilizar; ahora vendemos el plástico a una empresa de reciclaje », explica. Sin embargo, quedan muchos problemas: los camiones no pasan todos los días y algunos vecinos todavía tiran la basura en cualquier sitio.

Los especialistas recuerdan la regla de las tres erres: reducir, reutilizar y reciclar. Reducir significa comprar menos productos con embalajes inútiles; reutilizar consiste en dar una segunda vida a los objetos; reciclar permite fabricar cosas nuevas con materiales viejos. En mi opinión, este ejemplo demuestra que el cambio empieza en cada casa. Si un barrio de Yaoundé puede hacerlo, ¿por qué no toda la ciudad? El medioambiente no espera: actuar es un deber de todos.
\end{textebox}

\textbf{Glossaire :} \esp{vecinos} = voisins ; \esp{pozos} = puits ; \esp{vidrio} = verre ; \esp{materia orgánica} = matière organique ; \esp{compost} = compost ; \esp{embalajes} = emballages ; \esp{tiran} = jettent ; \esp{en cualquier sitio} = n'importe où.

\textbf{Questions de préparation au commentaire :}
\begin{enumerate}
\item \esp{¿Cuál es la fuente y el tema principal del texto?}
\item \esp{¿Qué problemas denuncia el texto? ¿Qué soluciones propone?}
\item \esp{¿Qué papel juega el conector « sin embargo »?}
\item \esp{¿Cuál es el punto de vista del autor: pesimista u optimista?}
\item \esp{¿Y en tu barrio? ¿Qué deberíamos hacer?}
\end{enumerate}

\begin{exoresolu}[Rédiger l'introduction du commentaire]
Énoncé : rédigez en espagnol l'introduction (3 phrases) du commentaire du texte « Un barrio que aprendió a reciclar », en utilisant deux phrases témoins de la méthode.
\tcblower
Solution détaillée : on suit l'ordre source $\rightarrow$ thème $\rightarrow$ problématique.

\esp{Este texto, extraído de un artículo de prensa, trata del reciclaje en un barrio de Yaoundé. El tema principal es la recogida selectiva y la regla de las tres erres: reducir, reutilizar y reciclar. El autor muestra que el cambio es posible y se pregunta: ¿por qué no toda la ciudad?}

Traduction : « Ce texte, extrait d'un article de presse, traite du recyclage dans un quartier de Yaoundé. Le thème principal est la collecte sélective et la règle des trois R : réduire, réutiliser et recycler. L'auteur montre que le changement est possible et se demande : pourquoi pas toute la ville ? »

Points de méthode : la source est présentée (\esp{extraído de un artículo de prensa}), le thème est nommé (\esp{trata de...}), et la problématique est formulée sous forme de question — une technique très appréciée des correcteurs. Notez la contraction obligatoire \esp{trata del reciclaje} (\esp{de} + \esp{el} = \esp{del}).
\end{exoresolu}

\courssub{Les règles d'or du commentaire}

\begin{aretenir}[Règles d'or]
\begin{itemize}
\item \cle{Ne jamais recopier} le texte : résumer avec ses propres mots, citer brièvement entre guillemets.
\item \cle{Annoncer chaque partie} avec un connecteur : \esp{en primer lugar}, \esp{por un lado}, \esp{sin embargo}, \esp{en conclusión}.
\item \cle{Garder son avis pour la fin} : l'introduction et le développement parlent du texte et de l'auteur, pas de vous.
\item \cle{Surveiller la grammaire de l'opinion} : indicatif après \esp{pienso que}, subjonctif après \esp{no pienso que}.
\item \cle{Soigner la langue} : accents (\esp{opinión}, \esp{reciclaje}, \esp{contaminación}), ponctuation espagnole (¿...? ¡...!), pas de calque du français (\esp{según el autor}, jamais \esp{según de}).
\end{itemize}
\end{aretenir}

\begin{savaistu}[Comment le commentaire est-il noté au Bac ?]
Au baccalauréat camerounais, le commentaire de texte est évalué selon trois critères : la \cle{compréhension} (avez-vous saisi l'idée principale et le point de vue de l'auteur ?), la \cle{langue} (correction grammaticale, richesse du lexique, accents) et l'\cle{organisation} (plan visible, connecteurs, progression logique). Beaucoup d'élèves soignent leurs phrases mais perdent des points faute de plan. Retenez ceci : un plan clair en trois parties vaut autant que de belles phrases. Les correcteurs lisent vite : si vos connecteurs (\esp{en primer lugar}, \esp{sin embargo}, \esp{en conclusión}) balisent votre copie, vous avez déjà gagné la moitié du chemin.
\end{savaistu}

\begin{exercice}[À vous de jouer]
\begin{enumerate}
\item Rédigez le premier paragraphe du développement du commentaire sur « Un barrio que aprendió a reciclar » : commencez par \esp{En primer lugar}, résumez le problème décrit, citez brièvement la presidenta de la asociación, puis analysez le rôle de \esp{sin embargo}.
\item Rédigez la conclusion : bilan en une phrase (\esp{En conclusión...}), puis votre avis sur la situation à Douala ou dans votre ville, en utilisant \esp{a mi juicio} et une phrase avec \esp{no pienso que} + subjonctif.
\item Traduisez en espagnol : « Selon l'auteur, recycler est un devoir de tous, mais je ne pense pas que ce soit facile au Cameroun. »
\end{enumerate}
\end{exercice}

\begin{corrige}[Exercice : indications]
1. \esp{En primer lugar, el texto describe el problema de la basura en el barrio: antes se acumulaba en las calles y contaminaba el agua. Como explica la presidenta: « ahora vendemos el plástico a una empresa de reciclaje ». Sin embargo, el autor matiza este optimismo: los camiones no pasan todos los días.} (Traduction : « En premier lieu, le texte décrit le problème des déchets dans le quartier : avant, ils s'accumulaient dans les rues et contaminaient l'eau. Comme l'explique la présidente : "maintenant nous vendons le plastique à une entreprise de recyclage". Cependant, l'auteur nuance cet optimisme : les camions ne passent pas tous les jours. »)

2. Modèle : \esp{En conclusión, el texto demuestra que reciclar es posible si los vecinos colaboran. A mi juicio, en Douala deberíamos instalar más contenedores y educar a los jóvenes. No pienso que sea fácil, pero el ejemplo de la Briqueterie muestra que el cambio empieza en cada casa.} (Traduction : « En conclusion, le texte démontre que recycler est possible si les voisins collaborent. À mon avis, à Douala nous devrions installer plus de conteneurs et éduquer les jeunes. Je ne pense pas que ce soit facile, mais l'exemple de la Briqueterie montre que le changement commence dans chaque maison. »)

3. \esp{Según el autor, reciclar es un deber de todos, pero no pienso que sea fácil en Camerún.} Attention : \esp{según el autor} sans « de », et subjonctif \esp{sea} après \esp{no pienso que}.
\end{corrige}

\coursec{Commentaire modèle}

Dans la section précédente, nous avons construit pas à pas la méthode du commentaire de texte : lire deux fois, repérer le thème, dégager le plan, rédiger une introduction, un développement en deux ou trois parties et une conclusion avec un avis personnel. Il est temps maintenant de voir cette méthode \emph{en action}. Nous allons lire ensemble un commentaire complet, rédigé selon les exigences du baccalauréat, puis le disséquer phrase par phrase pour comprendre \emph{pourquoi} il fonctionne. L'objectif n'est pas de l'apprendre par cœur, mais de vous en approprier les mécanismes pour produire le vôtre.

\courssub{Lecture du commentaire modèle}

Lisez d'abord le commentaire ci-dessous en silence, crayon à la main. Consigne de classe : \textbf{soulignez tous les connecteurs} (mots de liaison) et \textbf{encadrez les citations courtes} reprises du texte. Vous vérifierez ensuite vos annotations avec l'analyse qui suit.

\begin{textebox}[Comentario modelo : « El medioambiente y nosotros »]
Este texto, titulado « El medioambiente y nosotros », trata del problema de la basura en las grandes ciudades y de las soluciones que existen, como la recogida selectiva y el reciclaje.

En primer lugar, el autor describe la situación de ciudades como Douala y Yaoundé, donde la basura se acumula en las calles. Cada habitante produce casi un kilo de residuos al día. Los vertederos están llenos y la salud de los vecinos está en peligro.

Por otro lado, el texto presenta el ejemplo de España, donde los ciudadanos separan sus residuos en contenedores de colores: amarillo para los envases, verde para el vidrio y azul para el papel. Después, los camiones de la recogida los llevan a las plantas de reciclaje, donde los materiales se transforman en productos nuevos.

Finalmente, el autor defiende las tres erres — reducir, reutilizar, reciclar — y afirma que « el medioambiente somos nosotros ». Con esta frase quiere decir que la naturaleza no es algo exterior: somos parte de ella.

En conclusión, este texto nos invita a cambiar nuestros hábitos. En mi opinión, en Camerún deberíamos organizar la recogida selectiva en todas las ciudades, porque cuidar el medioambiente es cuidar nuestra salud.
\end{textebox}

\textbf{Glossaire et remarques de traduction}

\begin{center}
\begin{tabularx}{\linewidth}{|X|X|}\hline
\rowcolor{popblueL} \textbf{Español} & \textbf{Français} \\ \hline
\esp{trata del problema de...} & traite du problème de... \\ \hline
\esp{la recogida selectiva} & la collecte sélective / le tri sélectif \\ \hline
\esp{se acumula en las calles} & s'accumule dans les rues \\ \hline
\esp{los vertederos están llenos} & les décharges sont pleines \\ \hline
\esp{los envases} & les emballages \\ \hline
\esp{el vidrio} & le verre \\ \hline
\esp{las plantas de reciclaje} & les centres / usines de recyclage \\ \hline
\esp{las tres erres} & les trois « R » \\ \hline
\esp{cambiar nuestros hábitos} & changer nos habitudes \\ \hline
\esp{deberíamos organizar} & nous devrions organiser \\ \hline
\end{tabularx}
\end{center}

\begin{exemplebox}[Deux phrases clés traduites]
\esp{El texto nos invita a cambiar nuestros hábitos.} « Le texte nous invite à changer nos habitudes. » Remarquez la construction \esp{invitar a + infinitif}, identique au français.

\esp{Cuidar el medioambiente es cuidar nuestra salud.} « Prendre soin de l'environnement, c'est prendre soin de notre santé. » Ici, l'infinitif \esp{cuidar} est \emph{sujet} de la phrase, exactement comme l'infinitif français.
\end{exemplebox}

\courssub{Analyse du modèle : l'architecture du commentaire}

Comptons d'abord : le modèle contient environ \cle{170 mots}. C'est exactement la longueur attendue au baccalauréat.

\begin{methode}[Compter les mots de son commentaire]
\begin{enumerate}
\item Un commentaire de Bac fait environ \cle{150 à 200 mots} : ni une simple réponse de trois lignes, ni une dissertation de deux pages.
\item Pour estimer sans tout compter : comptez les mots d'une ligne complète (environ 8 à 12 mots selon votre écriture), puis multipliez par le nombre de lignes.
\item Vérifiez la répartition : introduction courte (1--2 phrases), développement (3 paragraphes), conclusion courte avec avis personnel.
\item Si vous dépassez 220 mots, supprimez les répétitions ; si vous êtes sous 130 mots, ajoutez une citation courte commentée ou un exemple du texte.
\end{enumerate}
\end{methode}

Observons maintenant la structure. Chaque partie du modèle remplit une \cle{fonction} précise, et chaque fonction est signalée par un \cle{connecteur} :

\begin{center}
\begin{tabularx}{\linewidth}{|X|X|X|}\hline
\rowcolor{popblueL} \textbf{Partie du commentaire} & \textbf{Fonction} & \textbf{Connecteurs utilisés} \\ \hline
Introduction (§1) & Introduire le texte et annoncer le thème & \esp{Este texto, titulado..., trata de...} \\ \hline
Développement §2 & Décrire le problème & \esp{En primer lugar} \\ \hline
Développement §3 & Présenter les solutions & \esp{Por otro lado}, \esp{Después} \\ \hline
Développement §4 & Dégager le message de l'auteur & \esp{Finalmente}, \esp{Con esta frase quiere decir que...} \\ \hline
Conclusion (§5) & Conclure et donner un avis personnel & \esp{En conclusión}, \esp{En mi opinión}, \esp{porque} \\ \hline
\end{tabularx}
\end{center}

Trois techniques rendent ce modèle solide :

\begin{itemize}
\item \textbf{Les citations courtes.} Le commentateur cite deux fois le texte : le titre \esp{« El medioambiente y nosotros »} dans l'introduction, et la phrase-clé \esp{« el medioambiente somos nosotros »} dans la troisième partie. Une citation courte, bien choisie et \emph{commentée} (\esp{Con esta frase quiere decir que...}), prouve que vous avez réellement lu le texte.
\item \textbf{La progression logique.} Problème (\esp{En primer lugar}) $\rightarrow$ solutions (\esp{Por otro lado}) $\rightarrow$ message (\esp{Finalmente}) $\rightarrow$ avis (\esp{En conclusión}). Le correcteur suit votre raisonnement sans effort.
\item \textbf{L'avis personnel en conclusion.} \esp{En mi opinión, en Camerún deberíamos organizar...} : l'opinion est placée \emph{à la fin}, jamais au début, et elle est justifiée par \esp{porque}. C'est exactement ce que le jury attend.
\end{itemize}

\begin{propriete}[L'infinitif sujet : un outil précieux pour le commentaire]
En espagnol comme en français, l'\cle{infinitif} peut occuper la fonction de \textbf{sujet} de la phrase. Le verbe principal s'accorde alors à la 3\ieme personne du singulier. C'est une structure simple, correcte et très élégante pour exprimer une idée générale dans un commentaire.

\textbf{Forme :} \esp{Infinitivo + verbo (3\ieme pers. sing.) + complément}.

\begin{itemize}
\item \esp{Reciclar \textbf{es} importante.} « Recycler \textbf{est} important. » (verbe \esp{ser})
\item \esp{Reducir los residuos \textbf{es} un deber de todos.} « Réduire les déchets \textbf{est} un devoir de tous. » (verbe \esp{ser})
\item \esp{Reutilizar las bolsas \textbf{ayuda} al planeta.} « Réutiliser les sacs \textbf{aide} la planète. » (verbe \esp{ayudar})
\item \esp{Tirar la basura al suelo \textbf{es} peligroso para la salud.} « Jeter les ordures par terre \textbf{est} dangereux pour la santé. » (verbe \esp{ser})
\end{itemize}

\textbf{Nuance importante :} quand le sujet est une proposition subordonnée avec un verbe conjugué, l'espagnol exige souvent le subjonctif après des expressions impersonnelles : \esp{Es importante que reciclemos.} « Il est important que nous recyclions. » Mais avec l'infinitif sujet, \textbf{pas de subjonctif} : \esp{Reciclar es importante.}
\end{propriete}

\begin{attention}[Pièges fréquents des francophones dans le commentaire]
\begin{itemize}
\item Ne traduisez pas « le texte parle de » par \esp{el texto habla de} de façon systématique : préférez \esp{el texto trata de} ou \esp{el autor describe / presenta / defiende}, plus soutenus.
\item « Donner son avis » ne se dit pas \esp{dar su aviso} : \esp{un aviso} est une « annonce, un avertissement ». Dites \esp{dar su opinión} ou \esp{en mi opinión}.
\item Attention au conditionnel pour l'avis poli : \esp{deberíamos} (nous devrions) est le conditionnel de \esp{deber}, pas le futur. N'écrivez pas \esp{deberemos}.
\item N'oubliez pas les guillemets et la cohérence : une citation doit être \emph{exacte} et courte ; ne recopiez jamais trois lignes du texte.
\item \esp{Por otro lado} s'écrit en trois mots ; ne collez pas \esp{porotrolado}.
\end{itemize}
\end{attention}

\begin{popfigure}
\begin{center}
\begin{tikzpicture}[mindmap, grow cyclic, every node/.style={concept, align=center}, concept color=popblue!40, text=popdark, root concept/.append style={concept color=poporange!50, font=\bfseries}, level 1/.append style={level distance=3.4cm, sibling angle=90}, level 2/.append style={level distance=2.6cm, sibling angle=45}]
\node[root concept]{Comentario modelo (170 palabras)}
  child[concept color=popgreen!40]{ node[align=center]{Introducción\\ \esp{trata de...}} }
  child[concept color=popblue!30]{ node[align=center]{Problema\\ \esp{En primer lugar}} }
  child[concept color=poppurple!30]{ node[align=center]{Soluciones\\ \esp{Por otro lado}} }
  child[concept color=popgold!40]{ node[align=center]{Conclusión\\ \esp{En mi opinión}} };
\end{tikzpicture}
\end{center}
\legende{Carte mentale de la structure du commentaire modèle : quatre blocs, quatre connecteurs repères.}
\end{popfigure}

\courssub{Entraînement guidé}

\begin{exoresolu}[Repérer et compléter]
\textbf{Énoncé.} Voici trois phrases extraites d'un commentaire d'élève. Complétez-les avec le connecteur qui convient : \esp{en primer lugar}, \esp{por otro lado}, \esp{en conclusión}, \esp{en mi opinión}, \esp{finalmente}.

\begin{enumerate}
\item \esp{............, el autor explica que la basura contamina el agua de los barrios.}
\item \esp{............, propone crear contenedores de colores en cada mercado de Douala.}
\item \esp{............, este texto es una llamada a la acción ; ............, todos podemos hacer algo por nuestra ciudad.}
\end{enumerate}

\tcblower
\textbf{Solution détaillée.}
\begin{enumerate}
\item \esp{\textbf{En primer lugar}, el autor explica que la basura contamina el agua de los barrios.} — La phrase décrit le \emph{problème}, première étape du développement : on ouvre avec \esp{en primer lugar}.
\item \esp{\textbf{Por otro lado}, propone crear contenedores de colores en cada mercado de Douala.} — On passe du problème aux \emph{solutions} : le connecteur d'alternance \esp{por otro lado} marque ce changement de perspective.
\item \esp{\textbf{En conclusión}, este texto es una llamada a la acción ; \textbf{en mi opinión}, todos podemos hacer algo por nuestra ciudad.} — Dernière partie : \esp{en conclusión} annonce le bilan, puis \esp{en mi opinión} introduit l'avis personnel, toujours placé en fin de commentaire. (\esp{Finalmente} aurait pu convenir pour la dernière partie du développement, mais ici nous sommes déjà dans la conclusion.)
\end{enumerate}
\end{exoresolu}

\begin{exercice}[À vous de commenter]
À partir du texte « El medioambiente y nosotros » étudié dans cette leçon, rédigez votre propre commentaire en 150 à 200 mots, en imitant l'architecture du modèle mais avec vos propres phrases :
\begin{enumerate}
\item une introduction avec \esp{Este texto, titulado..., trata de...} ;
\item un développement en trois paragraphes (problème / solutions / message) ouverts par \esp{En primer lugar}, \esp{Por otro lado}, \esp{Finalmente} ;
\item au moins une citation courte du texte, commentée avec \esp{Con esta frase quiere decir que...} ;
\item une conclusion avec \esp{En conclusión} et un avis personnel justifié (\esp{En mi opinión... porque...}) ;
\item au moins une phrase avec un infinitif sujet, par exemple \esp{Reciclar es...}.
\end{enumerate}
\end{exercice}

\begin{corrige}[Exercice : éléments de correction attendus]
Un bon commentaire d'élève doit contenir : une introduction qui nomme le titre et le thème (\esp{trata del problema de la basura y de las soluciones como el reciclaje}) ; un premier paragraphe sur l'accumulation des déchets à Douala et Yaoundé ; un deuxième sur l'exemple espagnol des conteneurs de couleurs et les usines de recyclage ; un troisième sur les trois erres avec la citation \esp{« el medioambiente somos nosotros »} ; une conclusion avec un avis personnel ancré au Cameroun, par exemple : \esp{En mi opinión, las escuelas de Yaoundé deberían enseñar la recogida selectiva, porque reducir la basura es proteger la salud de todos.} Vérifiez : connecteurs présents, citation exacte, avis en conclusion, 150--200 mots, accents corrects.
\end{corrige}

\begin{experience}[Production orale : Gestes écocitoyens dans mon quartier]
\textbf{Objectif :} Réinvestir le lexique du recyclage et les structures d'opinion (\esp{deberíamos, es necesario que + subj., en mi opinión}) à l'oral.

\textbf{Déroulement (15 min) :}
\begin{enumerate}
\item \textbf{Préparation (5 min) :} En binôme, imaginez une mini-campagne de sensibilisation pour votre quartier (Essos, Nkolbisson, Akwa, Bonapriso...). Choisissez 3 gestes concrets (ex. : séparer les plastiques, composter les épluchures, réutiliser les sacs « Ghana must go »).
\item \textbf{Rédigez un court script (5 phrases) :} Utilisez \esp{En primer lugar...}, \esp{Por otro lado...}, \esp{En conclusión...}. Intégrez au moins un infinitif sujet (\esp{Separar la basura es fácil.}) et le vocabulaire : \esp{contenedor amarillo/verde/azul, vertedero, recogida selectiva, reducir, reutilizar, reciclar}.
\item \textbf{Présentation (5 min) :} Chaque binôme présente sa campagne devant la classe (1 min). Les autres notent les connecteurs entendus et votent pour la proposition la plus réaliste.
\end{enumerate}

\textbf{Exemple de début de script :}
\esp{« Vecinos de Nkolbisson: En primer lugar, nuestro barrio está sucio. Por otro lado, podemos cambiar esto. Separar la basura es fácil. En mi opinión, deberíamos pedir contenedores al ayuntamiento. ¡Reciclar es vivir mejor! »}
\end{experience}

\begin{aretenir}[L'essentiel du commentaire modèle]
Un commentaire réussi tient en quatre gestes : \textbf{introduire} (\esp{Este texto trata de...}), \textbf{décrire le problème} (\esp{En primer lugar}), \textbf{présenter les solutions et le message} (\esp{Por otro lado}, \esp{Finalmente} + une citation courte commentée), \textbf{conclure avec un avis justifié} (\esp{En conclusión}, \esp{En mi opinión... porque...}). Longueur idéale : 150 à 200 mots. Et n'oubliez pas l'infinitif sujet, structure simple et élégante : \esp{Cuidar el medioambiente es cuidar nuestra salud.}
\end{aretenir}

\coursec{Production guidée : gestes écocitoyens}

Dans la section précédente, nous avons commenté le texte \esp{El medioambiente y nosotros} et dégagé sa problématique : la responsabilité de chacun face aux déchets. Mais un bon élève de Terminale ne s'arrête pas à l'analyse : il passe à l'\cle{action}. Cette dernière étape de la leçon vous apprend à \textbf{exprimer des conseils et des obligations} en espagnol et à \textbf{rédiger un paragraphe argumenté} proposant des \cle{gestes écocitoyens} pour votre ville, que ce soit Yaoundé, Douala ou votre village.

\courssub{Observer avant d'écrire : un texte modèle}

Avant d'énoncer les règles, lisons ce court texte original, écrit dans le style d'un article de journal de lycée. Repérez les verbes soulignés par leur structure : que remarquez-vous ?

\begin{textebox}[Un barrio más limpio — artículo de un periódico escolar]
En nuestro barrio de Yaoundé, la basura es un problema serio: los montones de residuos se acumulan en las esquinas y, cuando llueve, el plástico obstruye las canaletas. Sin embargo, cada uno de nosotros puede actuar. \textbf{Hay que} cambiar nuestros hábitos desde hoy.

Primero, \textbf{hay que} separar la basura en casa: el plástico, el papel y los residuos orgánicos no deben mezclarse. Además, \textbf{deberíamos} usar bolsas reutilizables cuando vamos al mercado, en lugar de aceptar bolsas de plástico cada vez. \textbf{Es importante} no tirar nada al suelo, ni siquiera un papel pequeño, porque todo termina en los ríos.

Segundo, \textbf{tenemos que} ahorrar agua y electricidad: cerrar el grifo mientras nos cepillamos los dientes y apagar las luces son gestos sencillos pero eficaces. \textbf{Deberíamos} también reciclar el papel de los cuadernos viejos y reutilizar los envases de vidrio.

Finalmente, \textbf{es importante} plantar árboles en el patio del colegio y sensibilizar a los vecinos. \textbf{Hay que} recordar que un barrio limpio es un barrio sano: menos moscas, menos enfermedades, menos inundaciones.

Por eso, queridos compañeros, \textbf{tenemos que} actuar juntos. Si cada alumno adopta tres gestos ecociudadanos, nuestra ciudad entera cambiará. El medioambiente somos nosotros: \textbf{hay que} protegerlo hoy para vivir mejor mañana.
\end{textebox}

\textbf{Glossaire :} \esp{los montones} = les tas ; \esp{obstruye} = bouche ; \esp{las canaletas} = les caniveaux ; \esp{los hábitos} = les habitudes ; \esp{los envases de vidrio} = les emballages en verre ; \esp{el grifo} = le robinet ; \esp{cepillarse los dientes} = se brosser les dents ; \esp{sensibilizar} = sensibiliser ; \esp{las inundaciones} = les inondations ; \esp{los compañeros} = les camarades.

\textbf{Questions d'observation :}
\begin{enumerate}
\item Relevez toutes les expressions suivies d'un infinitif. Combien y en a-t-il ?
\item Quelle est la différence de sens entre \esp{hay que separar} et \esp{deberíamos usar} ?
\item Quels connecteurs organisent le texte ?
\end{enumerate}

Vous avez observé quatre structures différentes. Passons maintenant à la règle.

\courssub{Les structures du conseil et de l'obligation}

\begin{propriete}[Hay que + infinitif : l'obligation impersonnelle]
\esp{Hay que + infinitivo} exprime une \cle{obligation générale et impersonnelle}, valable pour tout le monde. Elle correspond au français « il faut ».
\begin{itemize}
\item \esp{Hay que separar los residuos.} = « Il faut trier les déchets. »
\item \esp{Hay que proteger el medioambiente.} = « Il faut protéger l'environnement. »
\end{itemize}
\esp{Deberíamos + infinitivo} (conditionnel de \esp{deber}, 1\textsuperscript{re} personne du pluriel) exprime un \cle{conseil collectif}, plus doux : « nous devrions ».
\begin{itemize}
\item \esp{Deberíamos reciclar más.} = « Nous devrions recycler davantage. »
\item \esp{Deberíamos usar bolsas reutilizables.} = « Nous devrions utiliser des sacs réutilisables. »
\end{itemize}
\end{propriete}

\begin{propriete}[Les autres structures utiles]
\begin{itemize}
\item \esp{Es importante + infinitivo} = « il est important de » : \esp{Es importante ahorrar agua.}
\item \esp{Tenemos que + infinitivo} = « nous devons » (obligation personnelle, forte) : \esp{Tenemos que cuidar nuestro barrio.}
\item \esp{Es necesario + infinitivo} = « il est nécessaire de » : \esp{Es necesario reducir el plástico.}
\item \esp{No hay que + infinitivo} = « il ne faut pas » : \esp{No hay que tirar basura al suelo.}
\end{itemize}
\end{propriete}

\begin{exemplebox}[Exemples traduits]
\esp{Es importante no tirar plástico al suelo.} = « Il est important de ne pas jeter de plastique par terre. »

\esp{Tenemos que cuidar nuestro barrio.} = « Nous devons prendre soin de notre quartier. »

\esp{No hay que quemar la basura en la calle.} = « Il ne faut pas brûler les ordures dans la rue. »

\esp{Deberíamos plantar más árboles en el colegio.} = « Nous devrions planter plus d'arbres au lycée. »
\end{exemplebox}

\begin{attention}[Pièges des francophones]
\begin{itemize}
\item Après \esp{hay que}, on utilise \textbf{toujours l'infinitif}, jamais le subjonctif : on écrit \esp{Hay que reciclar}, et non une forme conjuguée. (Comparez : le français « il faut \emph{que} + subjonctif » n'a pas d'équivalent ici.)
\item Ne confondez pas \esp{hay que} (« il faut ») et \esp{hay} (« il y a ») : \esp{Hay mucha basura} = « Il y a beaucoup de déchets » ; \esp{Hay que recoger la basura} = « Il faut ramasser les déchets ».
\item Attention à l'accent : \esp{deberíamos} porte un accent sur le \emph{i} (conditionnel). Sans accent, \esp{debiamos} est une faute.
\item \esp{Tirar} = jeter (au sol) ; ne dites pas \esp{echar al suelo} par calque de « jeter par terre » dans ce contexte scolaire, préférez \esp{tirar}.
\end{itemize}
\end{attention}

\begin{center}
\begin{tabularx}{\linewidth}{|X|X|X|}
\hline
\rowcolor{popblueL} Español & Français & Nuance \\
\hline
Hay que + inf. & Il faut + inf. & Obligation impersonnelle, générale \\
\hline
Tenemos que + inf. & Nous devons + inf. & Obligation forte, personnelle \\
\hline
Deberíamos + inf. & Nous devrions + inf. & Conseil collectif, atténué \\
\hline
Es importante + inf. & Il est important de + inf. & Mise en valeur d'une idée \\
\hline
No hay que + inf. & Il ne faut pas + inf. & Interdiction générale \\
\hline
\end{tabularx}
\end{center}

\begin{exoresolu}[Traduisez en espagnol]
Énoncé : Traduisez : 1) Il faut réduire les déchets. 2) Nous devrions réutiliser les bouteilles en verre. 3) Il est important d'économiser l'eau. 4) Il ne faut pas jeter de papiers par terre.
\tcblower
Solution détaillée : 1) Obligation impersonnelle : \esp{Hay que reducir los residuos.} 2) Conseil collectif, conditionnel avec accent : \esp{Deberíamos reutilizar las botellas de vidrio.} 3) \esp{Es importante ahorrar agua.} (attention : \esp{ahorrar}, pas de calque de « économiser » avec \esp{economizar}, qui existe mais est rare dans ce sens). 4) \esp{No hay que tirar papeles al suelo.} (la négation se place avant \esp{hay}).
\end{exoresolu}

\courssub{La production guidée en trois étapes}

\begin{methode}[Comment construire votre paragraphe]
Avant de rédiger, écrivez votre \textbf{plan en trois lignes} au brouillon. Règle d'or : \textbf{une idée = une phrase = un connecteur}.
\begin{enumerate}
\item \textbf{Problemas} : listez trois problèmes de votre ville (ordures dans les rues, plastique dans les caniveaux, décharges sauvages, gaspillage de l'eau...). Utilisez \esp{hay} : \esp{En mi ciudad hay mucha basura en las calles.}
\item \textbf{Soluciones} : proposez trois solutions avec \esp{hay que} et \esp{deberíamos} : \esp{Hay que separar la basura. Deberíamos usar bolsas reutilizables.}
\item \textbf{Párrafo final} : rédigez un paragraphe de 8 à 10 lignes en reliant vos idées avec les connecteurs \esp{primero}, \esp{además}, \esp{finalmente}, \esp{por eso}.
\end{enumerate}
\end{methode}

\begin{popfigure}
\begin{tikzpicture}[node distance=0.9cm, every node/.style={font=\small}]
\node[draw=popblue, fill=popblueL, rounded corners, align=center, minimum width=4.2cm, minimum height=1.1cm] (p1) {\textbf{1. Problemas}\\ \esp{En mi ciudad hay...} (3)};
\node[draw=poporange, fill=poporangeL, rounded corners, align=center, minimum width=4.2cm, minimum height=1.1cm, below=of p1] (p2) {\textbf{2. Soluciones}\\ \esp{Hay que... / Deberíamos...} (3)};
\node[draw=popgreen, fill=popgreenL, rounded corners, align=center, minimum width=4.2cm, minimum height=1.1cm, below=of p2] (p3) {\textbf{3. Párrafo final}\\ \esp{Primero... Además...}\\ \esp{Finalmente... Por eso...}};
\draw[-{Stealth}, thick, popdark] (p1) -- (p2);
\draw[-{Stealth}, thick, popdark] (p2) -- (p3);
\end{tikzpicture}
\legende{Organigramme de la production guidée : des problèmes aux solutions, puis au paragraphe rédigé.}
\end{popfigure}

\textbf{Banque d'idées} (à réutiliser dans vos phrases) :

\begin{center}
\begin{tabularx}{\linewidth}{|X|X|}
\hline
\rowcolor{popgreenL} Español & Français \\
\hline
separar la basura & trier les ordures \\
\hline
usar bolsas reutilizables & utiliser des sacs réutilisables \\
\hline
no tirar plástico al suelo & ne pas jeter de plastique par terre \\
\hline
plantar árboles & planter des arbres \\
\hline
ahorrar agua y electricidad & économiser l'eau et l'électricité \\
\hline
reciclar el papel & recycler le papier \\
\hline
reutilizar los envases de vidrio & réutiliser les emballages en verre \\
\hline
limpiar las canaletas & nettoyer les caniveaux \\
\hline
\end{tabularx}
\end{center}

\begin{exemplebox}[Modèle de paragraphe rédigé]
\esp{En mi ciudad, Douala, hay tres problemas graves: hay mucha basura en las calles, el plástico obstruye las canaletas y la gente quema los residuos. Primero, hay que separar la basura en casa y llevar el plástico a los puntos de recogida. Además, deberíamos usar bolsas reutilizables cuando vamos al mercado y reciclar el papel de los cuadernos. Finalmente, es importante plantar árboles y sensibilizar a los vecinos. Por eso, tenemos que actuar juntos: si cada habitante adopta gestos ecociudadanos, Douala será una ciudad más limpia y más sana.}

« Dans ma ville, Douala, il y a trois problèmes graves : il y a beaucoup d'ordures dans les rues, le plastique bouche les caniveaux et les gens brûlent les déchets. D'abord, il faut trier les ordures à la maison et apporter le plastique aux points de collecte. De plus, nous devrions utiliser des sacs réutilisables quand nous allons au marché et recycler le papier des cahiers. Enfin, il est important de planter des arbres et de sensibiliser les voisins. C'est pourquoi nous devons agir ensemble : si chaque habitant adopte des gestes écocitoyens, Douala sera une ville plus propre et plus saine. »
\end{exemplebox}

\courssub{Grille d'auto-évaluation}

Après avoir rédigé votre paragraphe, vérifiez chaque critère et cochez :

\begin{center}
\begin{tabularx}{\linewidth}{|X|X|}
\hline
\rowcolor{popgoldL} Critère & Oui / Non \\
\hline
J'ai utilisé au moins 5 mots du lexique du recyclage (\esp{basura}, \esp{residuos}, \esp{reciclar}, \esp{reutilizar}, \esp{contenedor}...) & \\
\hline
J'ai utilisé au moins 2 structures de conseil (\esp{hay que}, \esp{deberíamos}, \esp{es importante}, \esp{tenemos que}) & \\
\hline
J'ai utilisé au moins 3 connecteurs (\esp{primero}, \esp{además}, \esp{finalmente}, \esp{por eso}) & \\
\hline
J'ai vérifié les accents (\esp{deberíamos}, \esp{plástico}, \esp{árboles}, \esp{electricidad}) & \\
\hline
Mon paragraphe compte entre 8 et 10 lignes & \\
\hline
\end{tabularx}
\end{center}

\begin{savaistu}[Un mot à connaître : ecociudadano]
Le mot \esp{ecociudadano/a} existe vraiment en espagnol : il désigne un citoyen qui adopte des gestes respectueux \esp{del medioambiente}. C'est un mot composé du préfixe \esp{eco-} (de l'écologie) et de \esp{ciudadano} (citoyen). Dans votre copie, écrire \esp{gestos ecociudadanos} montre un lexique riche et actuel — les correcteurs adorent !
\end{savaistu}

\courssub{Variante orale : le débat}

\begin{experience}[Débat en binômes : ¿Es posible la recogida selectiva en Douala ?]
Par groupes de deux, l'un défend le POUR, l'autre le CONTRE. Préparez trois arguments chacun, puis débattez devant la classe.

\textbf{Arguments posibles A FAVOR :} \esp{Hay que reducir los vertederos salvajes. Deberíamos crear empleos verdes para los jóvenes. Es importante proteger la salud de los habitantes.}

\textbf{Arguments posibles EN CONTRA :} \esp{No hay suficientes contenedores en los barrios. La gente no está sensibilizada. Los camiones de recogida no pasan regularmente.}

\textbf{Phrases pour débattre :} \esp{Estoy de acuerdo porque...} (je suis d'accord parce que...) ; \esp{No estoy de acuerdo, porque...} (je ne suis pas d'accord, parce que...) ; \esp{En mi opinión...} (à mon avis...) ; \esp{Sin embargo...} (cependant...).
\end{experience}

\begin{exercice}[À vous d'écrire]
Rédigez votre propre paragraphe de 8 à 10 lignes sur les problèmes environnementaux de votre quartier à Yaoundé ou de votre ville. Suivez l'organigramme en trois étapes, utilisez la banque d'idées, puis auto-évaluez-vous avec la grille avant de rendre votre copie.
\end{exercice}

\begin{corrige}[Exercice : éléments attendus]
Il n'y a pas une seule bonne réponse, mais une bonne copie contient : (1) trois problèmes introduits par \esp{hay} ou \esp{en mi barrio hay...} ; (2) au moins deux structures de conseil correctement suivies d'un infinitif ; (3) au moins trois connecteurs ; (4) cinq mots du lexique ; (5) des accents corrects. Exemple de phrase attendue : \esp{Hay que reciclar el papel y deberíamos plantar más árboles.}
\end{corrige}

\begin{aretenir}[L'essentiel de la production guidée]
Pour proposer des gestes écocitoyens en espagnol : \esp{hay que + infinitivo} (il faut), \esp{deberíamos + infinitivo} (nous devrions), \esp{es importante + infinitivo}, \esp{tenemos que + infinitivo}. Structurez toujours votre paragraphe : problèmes, solutions, conclusion, reliés par \esp{primero}, \esp{además}, \esp{finalmente}, \esp{por eso}. Et n'oubliez jamais : après \esp{hay que}, l'infinitif, rien d'autre.
\end{aretenir}

\newpage

\coursec{Activité d'intégration}

\courssub{Objectifs de l'activité}

Cette activité d'intégration clôt la leçon \esp{El medioambiente y nosotros}. Elle permet à l'élève de réinvestir, dans une tâche authentique et créative, l'ensemble des acquis de la séquence :

\begin{itemize}
\item \cle{Compréhension} : s'approprier un texte sur les déchets et le recyclage pour en extraire les idées utiles à sa propre production ;
\item \cle{Lexique} : employer correctement au moins cinq mots du champ lexical de la leçon (\esp{basura, residuos, contenedor, recogida selectiva, reciclaje, vertedero, reutilizar, reducir}) ;
\item \cle{Grammaire} : mobiliser les structures de conseil et d'obligation impersonnelle (\esp{hay que + infinitivo}, \esp{deberíamos + infinitivo}, \esp{es importante + infinitivo}) ;
\item \cle{Expression écrite} : rédiger un slogan, cinq gestes écocitoyens et un court paragraphe explicatif ;
\item \cle{Expression orale} : présenter une affiche en deux minutes avec fluidité et conviction ;
\item \cle{Compétence citoyenne} : prendre conscience du rôle de chacun dans la gestion des déchets à Douala ou à Yaoundé.
\end{itemize}

\courssub{Situation}

La mairie de votre quartier (Bali à Douala ou Mvog-Ada à Yaoundé) lance, en partenariat avec le lycée, une campagne de sensibilisation intitulée \esp{Mi barrio sin basura}. Les rues sont souvent encombrées de sachets plastiques, de bouteilles vides et de déchets organiques qui bouchent les caniveaux et provoquent des inondations à la saison des pluies. Le comité de quartier a publié l'appel à participation suivant :

\begin{textebox}[Appel à participation — Ayuntamiento del barrio]
\textbf{CAMPAÑA « MI BARRIO SIN BASURA »}\par
\emph{Convocatoria para los alumnos de Terminale}\par
¿Estás cansado de ver la basura en las calles de tu barrio? ¿Quieres actuar?\par
El comité del barrio invita a los jóvenes a crear un \textbf{cartel de sensibilización} en español. Cada cartel debe contener:\par
— un eslogan original y pegadizo;\par
— cinco gestos ecociudadanos para reducir, reutilizar y reciclar los residuos;\par
— un párrafo breve que explique el recorrido de un residuo: del contenedor al camión de recogida y a la planta de reciclaje.\par
Los mejores carteles se expondrán en el mercado del barrio y en la entrada del liceo. ¡Recicla hoy para vivir mañana!
\end{textebox}

\textbf{Traduction de l'appel :} « Tu en as assez de voir les ordures dans les rues de ton quartier ? Tu veux agir ? Le comité de quartier invite les jeunes à créer une affiche de sensibilisation en espagnol. Chaque affiche doit contenir : un slogan original et accrocheur ; cinq gestes écocitoyens pour réduire, réutiliser et recycler les déchets ; un court paragraphe qui explique le parcours d'un déchet : du conteneur au camion de ramassage puis à l'usine de recyclage. Les meilleures affiches seront exposées au marché du quartier et à l'entrée du lycée. Recycle aujourd'hui pour vivre demain ! »

\textbf{Glossaire :} \esp{pegadizo} = accrocheur (qui reste en tête) ; \esp{el recorrido} = le parcours ; \esp{expondrán} = exposeront (futur de \esp{exponer}) ; \esp{la entrada} = l'entrée.

Par groupes de quatre, vous êtes une équipe de jeunes volontaires du quartier. Vous créez votre affiche, puis vous la présentez à l'oral devant la classe, qui joue le rôle du jury du comité. Le groupe gagnant verra son affiche « exposée au marché ».

\courssub{Consignes}

\begin{enumerate}
\item \textbf{Relecture ciblée (10 min).} Relisez le texte \esp{El medioambiente y nosotros} et soulignez tous les mots du lexique des déchets. Recopiez-en au moins huit dans votre cahier de brouillon avec leur traduction française.
\item \textbf{Le slogan (10 min).} En groupe, inventez un slogan court, rythmé et impératif en espagnol. Utilisez l'impératif affirmatif (\esp{¡Recicla!}, \esp{¡Reduce!}, \esp{¡Protege!}) ou une formule avec \esp{para + infinitivo}. Vérifiez les accents et les signes ¡ !
\item \textbf{Les cinq gestes écocitoyens (15 min).} Rédigez cinq phrases de conseil en variant les structures : deux avec \esp{hay que + infinitivo}, deux avec \esp{deberíamos + infinitivo}, une avec \esp{es importante + infinitivo}. Chaque phrase doit contenir au moins un mot du lexique de la leçon.
\item \textbf{Le parcours d'un déchet (15 min).} Rédigez un paragraphe de quatre à six lignes qui décrit la chaîne : \esp{contenedor $\rightarrow$ camión de recogida $\rightarrow$ planta de reciclaje}. Utilisez les connecteurs \esp{primero, luego, después, finalmente} et le présent de l'indicatif.
\item \textbf{Fabrication de l'affiche (10 min).} Disposez le slogan en haut, les cinq gestes au centre, le paragraphe en bas. Soignez la lisibilité : pas de faute d'accent ni d'orthographe.
\item \textbf{Présentation orale (2 min par groupe).} Chaque membre du groupe prend la parole : un pour le slogan, deux pour les gestes, un pour le paragraphe. Parlez sans lire mot à mot.
\item \textbf{Vote du jury (5 min).} La classe écoute chaque présentation et remplit la grille d'évaluation, puis vote pour la meilleure campagne.
\end{enumerate}

\begin{methode}[Grille d'évaluation du jury]
Pour chaque groupe, notez de 0 à 2 chaque critère :
\begin{enumerate}
\item \textbf{Lexique} : au moins cinq mots de la leçon, employés correctement ;
\item \textbf{Grammaire} : structures de conseil correctes (infinitif après \esp{hay que}, \esp{deberíamos}, \esp{es importante}) ;
\item \textbf{Orthographe} : accents, ñ, signes ¡ ¿ ;
\item \textbf{Fluidité orale} : débit, prononciation, répartition des rôles ;
\item \textbf{Impact} : slogan accrocheur, affiche claire et convaincante.
\end{enumerate}
Total sur 10. Le groupe qui obtient la meilleure moyenne gagne la campagne.
\end{methode}

\courssub{Ressources à mobiliser}

\begin{aretenir}[Lexique de la leçon]
\esp{la basura} (les ordures) ; \esp{los residuos} (les déchets) ; \esp{el contenedor} (le conteneur) ; \esp{la recogida selectiva} (le tri sélectif, la collecte sélective) ; \esp{el reciclaje} (le recyclage) ; \esp{el vertedero} (la décharge) ; \esp{reutilizar} (réutiliser) ; \esp{reducir} (réduire) ; \esp{reciclar} (recycler) ; \esp{el camión de recogida} (le camion de ramassage) ; \esp{la planta de reciclaje} (l'usine de recyclage) ; \esp{el medioambiente} (l'environnement) ; \esp{ensuciar} (salir) ; \esp{tirar} (jeter) ; \esp{el embalaje} (l'emballage) ; \esp{el vidrio} (le verre).
\end{aretenir}

\begin{propriete}[Structures de conseil et d'obligation]
\begin{itemize}
\item \esp{hay que + infinitivo} : obligation générale, impersonnelle. \esp{Hay que separar los residuos.} « Il faut trier les déchets. »
\item \esp{deberíamos + infinitivo} : conseil collectif (conditionnel présent de \esp{deber}, 1\iere{} personne du pluriel). \esp{Deberíamos reutilizar las bolsas.} « Nous devrions réutiliser les sacs. »
\item \esp{es importante + infinitivo} : mise en valeur d'une action nécessaire. \esp{Es importante reciclar el plástico.} « Il est important de recycler le plastique. »
\item Impératif affirmatif pour le slogan : \esp{¡Recicla!} (tu), \esp{¡Reciclemos!} (nous : « recyclons ! »).
\end{itemize}
\textbf{Exception et nuance :} si le conseil vise une personne précise et non une vérité générale, l'espagnol passe au subjonctif : \esp{Es importante que recicles el plástico.} « Il est important que tu recycles le plastique. » Dans cette activité, on s'adresse à tous en général : on garde donc l'infinitif.
\end{propriete}

\begin{attention}[Pièges fréquents des francophones]
\begin{itemize}
\item Après \esp{hay que}, \esp{deberíamos} et \esp{es importante}, toujours l'\textbf{infinitif}, jamais de forme conjuguée : on écrit \esp{hay que reducir}, pas \esp{hay que reducimos}.
\item \esp{reciclaje} s'écrit avec un \emph{c} puis un \emph{j} ; \esp{recogida} avec un \emph{g} (ne pas confondre avec \esp{recolección}, qui désigne plutôt la récolte).
\item Ne traduisez pas « déchets » par \esp{desechos} dans tous les contextes : dans le texte de la leçon, le mot courant est \esp{residuos}.
\item Faux ami : \esp{el cartel} signifie « l'affiche », et non « le cartel » au sens français d'organisation criminelle.
\item Le slogan exige les signes doubles, collés à la phrase : \esp{¡Recicla hoy para vivir mañana!}
\end{itemize}
\end{attention}

\courssub{Proposition de production et corrigé commenté}

\begin{exoresolu}[Affiche modèle — Grupo « Bali Limpio », Douala]
Énoncé : rédigez l'affiche complète (slogan, cinq gestes, paragraphe) du groupe lauréat.
\tcblower
\textbf{Slogan :}\par
\esp{¡Bali sin basura, barrio con futuro! ¡Recicla hoy para vivir mañana!}\par
« Bali sans ordures, un quartier avec un avenir ! Recycle aujourd'hui pour vivre demain ! »\par
\textbf{Commentaire :} slogan bref, rime \esp{basura / futuro}, double impératif \esp{recicla} et structure de but \esp{para + infinitivo} vue en classe. Les signes ¡ ! encadrent correctement chaque formule.\par\medskip
\textbf{Cinco gestos ecociudadanos :}\par
\esp{1. Hay que separar los residuos orgánicos de los residuos de plástico.}\par
« Il faut séparer les déchets organiques des déchets plastiques. »\par
\esp{2. Hay que tirar la basura en el contenedor y nunca en la calle.}\par
« Il faut jeter les ordures dans le conteneur et jamais dans la rue. »\par
\esp{3. Deberíamos reutilizar las bolsas y las botellas de plástico.}\par
« Nous devrions réutiliser les sachets et les bouteilles en plastique. »\par
\esp{4. Deberíamos reducir el consumo de productos con mucho embalaje.}\par
« Nous devrions réduire la consommation de produits très emballés. »\par
\esp{5. Es importante participar en la recogida selectiva del barrio.}\par
« Il est important de participer au tri sélectif du quartier. »\par
\textbf{Commentaire :} les trois structures de conseil sont bien variées et toujours suivies de l'infinitif ; six mots du lexique de la leçon sont réinvestis (\esp{residuos, basura, contenedor, reutilizar, reducir, recogida selectiva}).\par\medskip
\textbf{El recorrido de un residuo :}\par
\esp{Primero, el ciudadano tira el residuo en el contenedor del barrio. Luego, el camión de recogida pasa cada martes y lleva los residuos a la planta de reciclaje. Después, en la planta, las máquinas separan el plástico, el vidrio y el papel. Finalmente, estos materiales se transforman en productos nuevos. Así, la basura no termina en el vertedero ni en los desagües de Douala.}\par
« D'abord, le citoyen jette le déchet dans le conteneur du quartier. Ensuite, le camion de ramassage passe chaque mardi et emporte les déchets à l'usine de recyclage. Puis, à l'usine, les machines trient le plastique, le verre et le papier. Enfin, ces matériaux sont transformés en produits nouveaux. Ainsi, les ordures ne finissent ni à la décharge ni dans les caniveaux de Douala. »\par
\textbf{Commentaire :} la chaîne \esp{contenedor $\rightarrow$ camión $\rightarrow$ planta de reciclaje} est respectée ; les quatre connecteurs chronologiques structurent le texte ; le présent de l'indicatif décrit un processus habituel ; la chute ancre le propos dans le contexte local (les caniveaux de Douala), ce qui renforce l'impact de la campagne.
\end{exoresolu}

\courssub{Bilan de l'activité}

À l'issue de cette campagne, l'élève a mobilisé simultanément les quatre savoir-faire de la leçon : il a \textbf{relu} le texte support pour y puiser le lexique, il a \textbf{réemployé} au moins cinq mots du champ des déchets dans un contexte nouveau, il a \textbf{appliqué} les structures de conseil à l'infinitif sans calquer le français, et il a \textbf{présenté à l'oral} une production collective en respectant un temps de parole limité. Le vote du jury développe en outre l'écoute active et l'esprit critique : chaque élève a dû évaluer la correction linguistique et la force persuasive des affiches de ses camarades. Cette tâche finale, ancrée dans la réalité des quartiers de Douala et de Yaoundé, montre que l'espagnol n'est pas seulement une langue d'examen : c'est un outil pour agir sur son environnement. \esp{¡Recicla hoy para vivir mañana!}

\newpage

\coursec{Méthodes et exercices résolus}

\courssub{Fiches méthodes et exercices résolus}

\begin{methode}[Lire et comprendre un texte sur les déchets]
\begin{enumerate}
\item \textbf{Première lecture globale.} Lis le texte une première fois sans dictionnaire. Repère le \cle{tema} (le sujet), le \cle{tipo de texto} (article, reportage, lettre) et le \cle{destinatario} (à qui il s'adresse).
\item \textbf{Repérage des mots-clés.} Souligne le champ lexical dominant : \esp{basura}, \esp{residuos}, \esp{contenedor}, \esp{recogida selectiva}, \esp{reciclaje}, \esp{vertedero}. Ces mots confirment le thème.
\item \textbf{Deuxième lecture analytique.} Pour chaque paragraphe, formule l'idée principale en une phrase : \esp{El primer párrafo presenta el problema; el segundo propone soluciones.}
\item \textbf{Exploitation du contexte.} Devine les mots inconnus par la famille (\esp{reutilizar} $\to$ \esp{utilizar} $\to$ « utiliser »), par le préfixe (\esp{re-} = « de nouveau ») ou par le contexte de la phrase.
\item \textbf{Réponse aux questions.} Réponds toujours par une \textbf{phrase complète} en espagnol, en reprenant les mots de la question, et justifie par une citation courte du texte entre guillemets.
\end{enumerate}
\textbf{Réflexes à retenir :} ne jamais répondre par « oui/non » seul ; ne jamais recopier tout un paragraphe ; citer 3 à 6 mots du texte suffit.
\end{methode}

\begin{exoresolu}[Compréhension de texte — type Bac]
\textbf{Texte :} \esp{En Douala, como en muchas ciudades africanas, la cantidad de residuos crece cada año. Sin embargo, la recogida selectiva todavía es poco frecuente: la mayoría de los ciudadanos mezcla todo en un solo contenedor. Algunas asociaciones de jóvenes han empezado a recoger botellas de plástico en los mercados para venderlas a empresas de reciclaje. « Reducir, reutilizar y reciclar: estas tres palabras pueden cambiar nuestra ciudad », explica un voluntario.}

\textbf{Questions :} 1. ¿Cuál es el problema principal mencionado en el texto? 2. ¿Qué hacen las asociaciones de jóvenes? 3. ¿Cuáles son las « tres palabras » citadas por el voluntario?
\tcblower
\textbf{Raisonnement.} Question 1 : le problème est annoncé dès la première phrase (\esp{la cantidad de residuos crece}) et précisé à la deuxième (\esp{la recogida selectiva todavía es poco frecuente}). Question 2 : la troisième phrase décrit l'action des associations. Question 3 : la citation finale donne la réponse littéralement.

\textbf{Réponses rédigées :}

1. \esp{El problema principal es que la cantidad de residuos crece cada año y que la recogida selectiva es todavía poco frecuente, porque « la mayoría de los ciudadanos mezcla todo en un solo contenedor ».}

2. \esp{Las asociaciones de jóvenes han empezado a recoger botellas de plástico en los mercados para venderlas a empresas de reciclaje.}

3. \esp{Las tres palabras son « reducir, reutilizar y reciclar ».}

\textbf{Points de grammaire.} \esp{han empezado} : passé composé (pretérito perfecto) avec l'auxiliaire \esp{haber} + participe \esp{empezado} (invariable). \esp{para venderlas} : \esp{para} + infinitif exprime le but ; le pronom \esp{las} (les bouteilles, féminin pluriel) se colle à l'infinitif. \esp{poco frecuente} : \esp{poco} est ici un \textbf{adverbe}, donc invariable ; c'est l'adjectif \esp{frecuente} qui s'accorde (\esp{una recogida poco frecuente}, \esp{unos ciudadanos poco frecuentes}).
\end{exoresolu}

\begin{methode}[Maîtriser et employer le lexique du recyclage]
\begin{enumerate}
\item \textbf{Organise le vocabulaire en réseaux.} Classe les mots en trois colonnes : les déchets (\esp{la basura, los residuos, los desechos}), les contenants et lieux (\esp{el contenedor, el cubo de basura, el vertedero}), les actions (\esp{tirar, recoger, separar, reutilizar, reducir, reciclar}).
\item \textbf{Apprends les mots avec leur article.} Le genre est indispensable : \esp{la basura}, \esp{el residuo}, \esp{el contenedor}, \esp{la recogida}, \esp{el reciclaje}, \esp{el vertedero}.
\item \textbf{Mémorise les expressions figées.} \esp{tirar la basura} (« jeter les ordures »), \esp{la recogida selectiva} (« le tri sélectif »), \esp{separar los residuos} (« trier les déchets »), \esp{proteger el medioambiente} (« protéger l'environnement »).
\item \textbf{Utilise les familles de mots.} \esp{reciclar} $\to$ \esp{el reciclaje} $\to$ \esp{reciclable} ; \esp{reutilizar} $\to$ \esp{la reutilización}. Le préfixe \esp{re-} exprime la répétition.
\item \textbf{Réemploie le lexique dans tes propres phrases.} Un mot n'est acquis que si tu sais l'utiliser : \esp{En mi barrio no hay contenedores para la recogida selectiva.}
\end{enumerate}
\end{methode}

\begin{definition}[Lexique thématique : déchets et recyclage]
\begin{center}
\begin{tabularx}{\linewidth}{|X|X|X|}\hline
\rowcolor{popblueL} \textbf{Español} & \textbf{Français} & \textbf{Remarque} \\ \hline
\esp{la basura} & les ordures, les déchets & féminin ; \esp{tirar la basura} \\ \hline
\esp{los residuos} & les déchets, les résidus & masculin ; plus soutenu que \esp{basura} \\ \hline
\esp{el contenedor} & le conteneur, la poubelle & masculin ; \esp{el cubo de basura} = la poubelle de cuisine \\ \hline
\esp{la recogida selectiva} & le tri sélectif & \esp{recoger} = ramasser, collecter \\ \hline
\esp{el reciclaje} & le recyclage & masculin, sans accent ; verbe : \esp{reciclar} \\ \hline
\esp{el vertedero} & la décharge & masculin ; \esp{verter} = verser, déverser \\ \hline
\esp{reutilizar} & réutiliser & préfixe \esp{re-} = « de nouveau » \\ \hline
\esp{reducir} & réduire & irrégulier : \esp{reduzco}, \esp{reduje} \\ \hline
\esp{el medioambiente} & l'environnement & un seul mot, masculin \\ \hline
\esp{la contaminación} & la pollution & féminin ; verbe : \esp{contaminar} \\ \hline
\end{tabularx}
\end{center}
\end{definition}

\begin{exoresolu}[Lexique — texte à trous et réemploi]
\textbf{Énoncé.} Complète avec : \esp{vertedero, reciclaje, contenedor, recogida selectiva, reutilizar, residuos}.

\esp{1. En Yaoundé, los \dots{} domésticos se mezclan a menudo en un solo \dots{}. 2. La \dots{} permite separar el plástico, el vidrio y el papel. 3. El \dots{} de botellas reduce la contaminación. 4. Podemos \dots{} las bolsas de tela muchas veces. 5. Los residuos no tratados terminan en el \dots{}.}

Puis traduis : « Le tri sélectif protège l'environnement. »
\tcblower
\textbf{Solutions des trous.}

1. \esp{los residuos domésticos} (« les déchets ménagers ») ; \esp{un solo contenedor} (« une seule poubelle »). On vérifie le genre : \esp{residuo} est masculin (\esp{los residuos}), \esp{contenedor} est masculin (\esp{un contenedor}).

2. \esp{La recogida selectiva} : sujet féminin singulier, d'où l'article \esp{la} et l'accord \esp{selectiva}.

3. \esp{El reciclaje de botellas} : \esp{reciclaje} est masculin (\esp{el}).

4. \esp{Podemos reutilizar las bolsas} : après \esp{poder}, on met l'infinitif.

5. \esp{el vertedero} : « la décharge », masculin.

\textbf{Traduction.} \esp{La recogida selectiva protege el medioambiente.}

\textbf{Justifications.} « Le tri sélectif » se traduit par l'expression figée \esp{la recogida selectiva} (et non un calque comme \esp{la selección}). « Protège » : présent de \esp{proteger}, 3\up{e} personne du singulier : \esp{protege}. Attention à l'irrégularité du verbe : \esp{g} $\to$ \esp{j} devant \esp{a} et \esp{o} (\esp{yo protejo}, \esp{que yo proteja}). « L'environnement » : \esp{el medioambiente} (un seul mot, masculin).
\end{exoresolu}

\begin{methode}[Commenter un texte au Baccalauréat]
\begin{enumerate}
\item \textbf{Introduction (2-3 phrases).} Présente le texte : \esp{Este texto, extraído de una revista, trata del problema de los residuos en las grandes ciudades.}
\item \textbf{Résumé bref.} Résume les idées principales en 2 ou 3 phrases avec tes propres mots : \esp{El autor explica que... / Según el texto...}
\item \textbf{Analyse.} Repère : le champ lexical (\esp{el vocabulario de la basura y el reciclaje}), le ton (\esp{el tono es informativo / alarmista / optimista}), les procédés (\esp{la cita de un voluntario da un ejemplo concreto}).
\item \textbf{Opinion personnelle argumentée.} Donne ton avis avec des connecteurs : \esp{En mi opinión... / Me parece que... / Estoy de acuerdo porque...} Illustre avec le contexte camerounais : \esp{En Douala, por ejemplo, ...}
\item \textbf{Conclusion.} Une phrase d'ouverture : \esp{En conclusión, el reciclaje es un reto para todos los ciudadanos.}
\end{enumerate}
\textbf{Structures à retenir :} \esp{El texto trata de...} ; \esp{El autor quiere demostrar que...} ; \esp{Según el texto...} ; \esp{En primer lugar / Además / Sin embargo / En conclusión}.
\end{methode}

\begin{exoresolu}[Commentaire guidé — rédaction]
\textbf{Énoncé.} À partir du texte suivant, rédige un commentaire de 8 à 10 lignes : présentation, résumé, ton, avis personnel.

\esp{Cada día, los camiones de higiene recorren los barrios de Yaoundé para recoger la basura. Pero el vertedero de la ciudad está casi lleno. La solución existe: separar los residuos en casa, reutilizar los objetos y reciclar el plástico. Algunos alumnos de un liceo han creado un club ecológico que organiza la limpieza del barrio cada sábado.}
\tcblower
\textbf{Démarche.} On suit les 5 étapes : présentation (type, thème), résumé (problème + solutions), analyse (champ lexical, ton), avis personnel ancré au Cameroun, conclusion.

\textbf{Commentaire rédigé (modèle) :}

\esp{Este texto es un artículo informativo que trata del problema de la basura en Yaoundé y de las soluciones posibles. El autor explica que el vertedero de la ciudad está casi lleno porque se recoge mucha basura cada día. Según el texto, la solución consiste en separar los residuos, reutilizar los objetos y reciclar el plástico. El vocabulario pertenece al campo léxico del medioambiente: basura, vertedero, residuos, reciclar. El tono es serio pero optimista, porque el texto da el ejemplo de los alumnos de un liceo que han creado un club ecológico. En mi opinión, esta iniciativa es muy importante: en mi barrio también podríamos organizar la limpieza cada sábado. En conclusión, el reciclaje y la educación de los ciudadanos son la clave para una ciudad limpia.}

\textbf{Points de grammaire.} \esp{está casi lleno} : \esp{estar} + adjectif pour un état (et non \esp{ser}). \esp{consiste en} + infinitif : construction figée. \esp{han creado} : pretérito perfecto pour une action récente liée au présent. \esp{podríamos} : conditionnel pour exprimer une suggestion polie.
\end{exoresolu}

\begin{methode}[Proposer des gestes écocitoyens en espagnol]
\begin{enumerate}
\item \textbf{Choisis la bonne structure selon le registre.}
\begin{itemize}
\item Conseil général : \esp{Hay que} + infinitif : \esp{Hay que separar los residuos.}
\item Obligation collective : \esp{Debemos / Tenemos que} + infinitif : \esp{Debemos reducir el plástico.}
\item Suggestion : \esp{Podríamos} + infinitif ou \esp{¿Por qué no...?} : \esp{¿Por qué no reutilizamos las bolsas?}
\item Interdiction : \esp{No se debe} + infinitif / \esp{Está prohibido} + infinitif : \esp{No se debe tirar basura en la calle.}
\end{itemize}
\item \textbf{Structure ta production en trois « R ».} \esp{Reducir} (consommer moins), \esp{Reutilizar} (se servir de nouveau), \esp{Reciclar} (transformer).
\item \textbf{Justifie chaque geste} avec \esp{para} + infinitif ou \esp{porque} + verbe conjugué : \esp{Hay que reciclar el plástico para proteger el medioambiente.}
\item \textbf{Conclus par un appel à l'action.} \esp{¡Actuemos ahora para un Camerún más limpio!}
\end{enumerate}
\end{methode}

\begin{exoresolu}[Production guidée — affiche écocitoyenne]
\textbf{Énoncé.} Ton lycée de Douala organise la « Semana verde ». Rédige 5 conseils écocitoyens pour une affiche, en utilisant au moins trois structures différentes (\esp{hay que}, \esp{debemos}, \esp{podríamos}, \esp{no se debe}).
\tcblower
\textbf{Raisonnement.} On varie les structures pour éviter la répétition, on utilise le lexique de la leçon (\esp{residuos, contenedor, reciclar, reutilizar, basura}) et on justifie au moins un conseil avec \esp{para} + infinitif.

\textbf{Affiche rédigée (modèle) :}

\esp{¡SEMANA VERDE EN NUESTRO LICEO!}

\esp{1. Hay que separar los residuos: el plástico, el papel y el vidrio en contenedores diferentes.}

\esp{2. Debemos reutilizar las bolsas de tela para reducir el plástico.}

\esp{3. Podríamos reciclar las botellas y venderlas a las empresas de reciclaje.}

\esp{4. No se debe tirar basura en la calle ni en los ríos.}

\esp{5. Tenemos que plantar árboles porque protegen el medioambiente.}

\esp{¡Actuemos juntos para una ciudad más limpia!}

\textbf{Points de grammaire.} \esp{Hay que} + infinitif : tournure impersonnelle invariable (jamais de verbe conjugué après \esp{hay que}). \esp{ni... ni} : la double négation est correcte et obligatoire en espagnol. \esp{Actuemos} : subjonctif présent de \esp{actuar}, 1\up{re} personne du pluriel, valeur d'impératif collectif (« agissons »). \esp{más limpia} : comparatif avec accent sur \esp{más} ; l'adjectif s'accorde avec \esp{ciudad} (féminin).
\end{exoresolu}

\begin{exoresolu}[Thème — traduction de phrases]
\textbf{Énoncé.} Traduis en espagnol : 1. « Nous devons jeter les déchets dans la poubelle. » 2. « Le recyclage est très important pour notre ville. » 3. « Il faut réutiliser les bouteilles en plastique. »
\tcblower
\textbf{Solutions détaillées.}

1. \esp{Debemos tirar los residuos al contenedor.} — « Nous devons » : \esp{debemos} (présent de \esp{deber} + infinitif \esp{tirar}). « Dans la poubelle » : contraction obligatoire \esp{a + el = al}. On peut aussi dire \esp{al cubo de la basura} pour la poubelle domestique.

2. \esp{El reciclaje es muy importante para nuestra ciudad.} — Ici \esp{ser} (et non \esp{estar}) car l'importance est présentée comme une caractéristique permanente. « Pour » de destination/bénéfice : \esp{para}.

3. \esp{Hay que reutilizar las botellas de plástico.} — « Il faut » : tournure impersonnelle \esp{hay que} + infinitif (jamais \esp{es necesario que reutilizamos} : après \esp{es necesario que}, il faudrait le subjonctif \esp{reutilicemos}). « Les bouteilles en plastique » : \esp{las botellas de plástico} — en espagnol, la matière s'exprime avec \esp{de} + nom sans article, contrairement au français « en ».

\textbf{Conclusion.} Ces trois phrases mobilisent les réflexes clés : \esp{deber/hay que} + infinitif, \esp{ser} pour la caractéristique, \esp{de} pour la matière.
\end{exoresolu}

\begin{attention}[Les 5 erreurs qui coûtent des points au Bac]
\begin{enumerate}
\item \textbf{Oublier les accents.} \esp{reciclaje} et \esp{vidrio} ne prennent pas d'accent, mais \esp{plástico}, \esp{árbol}, \esp{camión}, \esp{contaminación} et \esp{más} (comparatif) en prennent un. \esp{mas} sans accent signifie « mais » !
\item \textbf{Faux amis.} \esp{un residuo} = « un déchet, un résidu », pas « un reste de repas » (\esp{las sobras}) ; \esp{actualmente} = « actuellement, en ce moment », jamais « en fait » (qui se dit \esp{en realidad}) ; \esp{la basura} = « les ordures » (ne pas inventer un mot sur le français).
\item \textbf{Calquer le français.} « Jeter les ordures » se dit \esp{tirar la basura} ou \esp{echar la basura} ; le piège est d'inventer un verbe calqué sur « jeter ». « Trier les déchets » = \esp{separar los residuos} ou \esp{hacer la recogida selectiva}, pas un calque de « trier ».
\item \textbf{Confondre ser et estar.} \esp{El vertedero está lleno} (état) mais \esp{El reciclaje es importante} (caractéristique). L'accord de l'adjectif suit le nom : \esp{la ciudad está limpia}.
\item \textbf{Mal construire l'impersonnel.} « Il faut » = \esp{hay que} + infinitif, sans sujet et sans subjonctif. Après \esp{para} (but), toujours l'infinitif : \esp{para reciclar}, jamais un verbe conjugué.
\end{enumerate}
\end{attention}

\newpage

\coursec{Exercices}

\courssub{Je vérifie mes connaissances}

\begin{exercice}[QCM : vocabulaire des déchets]
Pour chaque question, choisis la bonne réponse.

\begin{enumerate}
\item \esp{El lugar donde se depositan los residuos de una ciudad es...}
\begin{enumerate}
\item \esp{el contenedor}
\item \esp{el vertedero}
\item \esp{la recogida}
\end{enumerate}
\item \esp{« Recogida selectiva » significa...}
\begin{enumerate}
\item trier les déchets selon leur matière
\item ramasser les déchets une fois par mois
\item brûler les déchets dans la rue
\end{enumerate}
\item \esp{El contenedor amarillo sirve para...}
\begin{enumerate}
\item \esp{el vidrio}
\item \esp{el papel}
\item \esp{los envases de plástico}
\end{enumerate}
\item \esp{« Reutilizar » veut dire...}
\begin{enumerate}
\item jeter
\item utiliser de nouveau
\item acheter
\end{enumerate}
\end{enumerate}
\end{exercice}

\begin{exercice}[Vrai ou faux justifié]
D'après le texte \esp{« El medioambiente y nosotros »}, dis si les affirmations suivantes sont vraies ou fausses et justifie ta réponse par une courte citation du texte.

\begin{enumerate}
\item \esp{Todos los residuos de la ciudad se reciclan.}
\item \esp{La separación de la basura en casa facilita el reciclaje.}
\item \esp{Reducir el consumo es tan importante como reciclar.}
\item \esp{El vidrio no se puede reciclar.}
\end{enumerate}
\end{exercice}

\begin{exercice}[Vocabulaire : trouve le mot]
Complète chaque phrase avec le mot espagnol qui convient : \esp{basura}, \esp{residuos}, \esp{contenedor}, \esp{reciclaje}, \esp{reutilizar}, \esp{reducir}.

\begin{enumerate}
\item En Douala, los camiones recogen la \ldots\ldots\ldots cada mañana.
\item Hay que echar las botellas de vidrio en el \ldots\ldots\ldots verde.
\item Podemos \ldots\ldots\ldots las bolsas de tela en lugar de comprar bolsas de plástico.
\item El \ldots\ldots\ldots del papel permite fabricar papel nuevo.
\item Para proteger el medioambiente, debemos \ldots\ldots\ldots el consumo de agua y de electricidad.
\item Los \ldots\ldots\ldots orgánicos pueden transformarse en abono para los campos.
\end{enumerate}
\end{exercice}

\begin{exercice}[Conjugaison à compléter]
Conjugue le verbe entre parenthèses au présent de l'indicatif.

\begin{enumerate}
\item Nosotros \ldots\ldots\ldots (separar) los residuos antes de tirarlos.
\item El ayuntamiento \ldots\ldots\ldots (recoger) la basura tres veces por semana.
\item Tú \ldots\ldots\ldots (reutilizar) las cajas de cartón, ¿verdad?
\item Los ciudadanos \ldots\ldots\ldots (poder) reducir sus residuos fácilmente.
\item Yo \ldots\ldots\ldots (querer) vivir en una ciudad limpia.
\end{enumerate}
\end{exercice}

\courssub{Je m'entraîne}

\begin{exercice}[Traduction : thème]
Traduis les phrases suivantes en espagnol.

\begin{enumerate}
\item Nous devons trier nos déchets pour protéger l'environnement.
\item À Yaoundé, la collecte des ordures est un grand problème.
\item Il faut réduire le plastique et réutiliser les sacs.
\item Le verre et le papier peuvent être recyclés.
\item Les jeunes doivent apprendre les gestes écocitoyens à l'école.
\end{enumerate}
\end{exercice}

\begin{exercice}[Reformulation]
Réécris chaque phrase avec tes propres mots, sans changer le sens (entraînement à la citation courte pour le commentaire).

\begin{enumerate}
\item \esp{La recogida selectiva es el primer paso del reciclaje.}
\item \esp{Debemos reducir la cantidad de residuos que producimos.}
\item \esp{El vertedero contamina el suelo y el agua.}
\item \esp{Reutilizar los objetos alarga su vida útil.}
\item \esp{Cada ciudadano es responsable del medioambiente.}
\end{enumerate}
\end{exercice}

\begin{exercice}[Texte à trous : la recogida selectiva]
Complète le texte suivant avec dix mots de la leçon : \esp{basura}, \esp{contenedor}, \esp{contenedores}, \esp{vertedero}, \esp{reutilizar}, \esp{reducir}, \esp{reciclaje}, \esp{recogida}, \esp{residuos}, \esp{vidrio}.

\begin{textebox}[La recogida selectiva]
En mi barrio, la \ldots\ldots\ldots de la basura se hace los lunes y los jueves. Cada familia separa sus \ldots\ldots\ldots en casa : el papel va al \ldots\ldots\ldots azul, el \ldots\ldots\ldots al contenedor verde y los envases de plástico al amarillo. Gracias a este sistema, el \ldots\ldots\ldots es más fácil y llega menos \ldots\ldots\ldots al \ldots\ldots\ldots de la ciudad. Además, intentamos \ldots\ldots\ldots el consumo de productos con mucho embalaje y \ldots\ldots\ldots los frascos de cristal para guardar alimentos. Así, los \ldots\ldots\ldots de colores nos ayudan a cuidar el medioambiente.
\end{textebox}
\end{exercice}

\begin{exercice}[Association : déchets et conteneurs]
Relie chaque déchet à son conteneur de couleur, puis rédige une phrase complète en espagnol pour chaque paire (modèle : \esp{El papel se echa en el contenedor azul.}).

\begin{center}
\begin{tabularx}{\linewidth}{|X|X|}
\hline
\rowcolor{popblueL} \textbf{Residuo} & \textbf{Contenedor} \\
\hline
\esp{el plástico y los envases} & \esp{contenedor verde} \\
\hline
\esp{el vidrio} & \esp{contenedor azul} \\
\hline
\esp{el papel y el cartón} & \esp{contenedor marrón} \\
\hline
\esp{la materia orgánica} & \esp{contenedor amarillo} \\
\hline
\end{tabularx}
\end{center}
\end{exercice}

\courssub{Je me prépare au Bac}

\begin{exercice}[Compréhension écrite type Bac]
Lis le texte suivant, puis réponds aux questions.

\begin{textebox}[El reciclaje en América Latina]
En muchas ciudades de América Latina, el reciclaje no es solo una obligación : es también una fuente de trabajo. Miles de personas, llamadas « recicladores », recorren las calles cada día para recoger cartón, plástico y metal que después venden a las empresas de reciclaje. Su labor es fundamental : gracias a ellos, una parte importante de los residuos no termina en los vertederos.

Sin embargo, estos trabajadores viven en condiciones difíciles. Muchos no tienen contrato, ni seguro, ni protección contra las enfermedades. Por eso, en varios países se han creado cooperativas que organizan a los recicladores, les ofrecen formación y mejoran sus ingresos.

Los gobiernos empiezan a comprender que el futuro del medioambiente pasa por la recogida selectiva en las casas, por la educación de los niños en las escuelas y por el reconocimiento del trabajo de los recicladores. Reducir, reutilizar y reciclar : estas tres palabras deben convertirse en hábitos diarios de todos los ciudadanos.
\end{textebox}

\textbf{Questions de compréhension} (réponds en espagnol, avec des phrases complètes) :

\begin{enumerate}
\item \esp{¿Qué hacen los « recicladores » cada día?}
\item \esp{¿Por qué es importante su trabajo para el medioambiente?}
\item \esp{¿En qué condiciones viven muchos de estos trabajadores?}
\item \esp{¿Qué hacen las cooperativas para ayudarles?}
\item \esp{Según el texto, ¿cuáles son las tres claves del futuro del medioambiente?}
\end{enumerate}

\textbf{Vrai ou faux} : justifie chaque réponse par une citation du texte.

\begin{enumerate}
\item \esp{Los recicladores venden los materiales que recogen.}
\item \esp{Todos los recicladores tienen un contrato y un seguro.}
\end{enumerate}

\textbf{Question de langue} : relève dans le texte trois mots de la famille de \esp{reciclar} et traduis-les en français.
\end{exercice}

\begin{exercice}[Thème et version]
\textbf{A. Thème.} Traduis en espagnol :

\begin{enumerate}
\item Les habitants de Douala jettent trop de déchets dans les rues.
\item Si nous trions nos ordures, le recyclage sera plus facile.
\item Il est nécessaire que les écoliers apprennent à protéger la nature.
\end{enumerate}

\textbf{B. Version.} Traduis en français :

\begin{textebox}[Un barrio limpio]
Desde hace un año, los vecinos de nuestro barrio separan la basura en casa. Los lunes, un camión recoge el papel y el cartón ; los jueves, el vidrio y el plástico. Gracias a este esfuerzo, las calles están más limpias y los niños juegan sin peligro. Todos estamos orgullosos de nuestro barrio.
\end{textebox}
\end{exercice}

\begin{exercice}[Production écrite type Bac]
\textbf{Sujet :} \esp{Escribe un párrafo (80-100 palabras) sobre los gestos ecociudadanos que podemos adoptar en tu ciudad.}

Consignes :
\begin{itemize}
\item Présente au moins \textbf{quatre gestes} concrets (tri, réutilisation, économies d'eau et d'énergie, propreté des rues...).
\item Emploie au moins \textbf{trois structures de conseil} différentes : \esp{hay que + infinitivo}, \esp{debemos + infinitivo}, \esp{es importante + infinitivo}, \esp{se debe + infinitivo}.
\item Utilise au moins \textbf{cinq mots du lexique} de la leçon (\esp{basura, residuos, contenedor, reciclaje, reutilizar, reducir, vertedero, recogida selectiva}).
\item Termine par une phrase d'appel à l'action avec le subjonctif : \esp{¡Cuidemos nuestro medioambiente!}
\end{itemize}
\end{exercice}

\newpage

\coursec{Corrigés des exercices}

\begin{corrige}[Exercice 1 — QCM : vocabulaire des déchets]
\begin{enumerate}
\item \textbf{Réponse b} : \esp{el vertedero}. Le \esp{vertedero} est la décharge, le lieu où l'on dépose les déchets d'une ville. Le \esp{contenedor} est la poubelle ou le conteneur de rue ; la \esp{recogida} est la collecte (le ramassage).
\item \textbf{Réponse a} : trier les déchets selon leur matière. \esp{Selectiva} vient de \esp{seleccionar} (sélectionner, trier) : la \esp{recogida selectiva} est le tri sélectif.
\item \textbf{Réponse c} : \esp{los envases de plástico}. Rappel du code couleur : jaune = plastique et emballages ; vert = verre ; bleu = papier et carton ; marron = matière organique.
\item \textbf{Réponse b} : utiliser de nouveau. Le préfixe \esp{re-} exprime la répétition : \esp{reutilizar} = réutiliser, comme \esp{reciclar} = recycler ou \esp{repetir} = répéter.
\end{enumerate}
\end{corrige}

\begin{corrige}[Exercice 2 — Vrai ou faux justifié]
\begin{enumerate}
\item \textbf{Faux.} Tous les déchets ne sont pas recyclés : une partie finit à la décharge. Justification possible : \esp{« una parte de los residuos termina en el vertedero »}.
\item \textbf{Vrai.} Le tri à la maison facilite le recyclage : \esp{« separar la basura en casa facilita el reciclaje »}.
\item \textbf{Vrai.} Le texte insiste sur les trois verbes : \esp{« reducir, reutilizar y reciclar »} ; réduire la consommation est présenté comme aussi important que recycler.
\item \textbf{Faux.} Le verre se recycle : on le jette dans le conteneur vert, \esp{« el vidrio se echa en el contenedor verde »}, et il peut être recyclé à l'infini.
\end{enumerate}
\textbf{Méthode :} pour un vrai/faux justifié, la citation doit être courte (5 à 10 mots), exacte, introduite par deux-points et placée entre guillemets. Une justification paraphrasée ou inventée ne rapporte aucun point.
\end{corrige}

\begin{corrige}[Exercice 3 — Vocabulaire : trouve le mot]
\begin{enumerate}
\item \esp{basura} — « À Douala, les camions ramassent les ordures chaque matin. »
\item \esp{contenedor} — « Il faut jeter les bouteilles de verre dans le conteneur vert. »
\item \esp{reutilizar} — « Nous pouvons réutiliser les sacs en tissu au lieu d'acheter des sacs en plastique. » (infinitif après \esp{podemos})
\item \esp{reciclaje} — « Le recyclage du papier permet de fabriquer du papier nouveau. »
\item \esp{reducir} — « Pour protéger l'environnement, nous devons réduire la consommation d'eau et d'électricité. » (infinitif après \esp{debemos})
\item \esp{residuos} — « Les déchets organiques peuvent être transformés en engrais pour les champs. »
\end{enumerate}
\textbf{Point de vigilance :} \esp{basura} (ordures ménagères, mot courant) et \esp{residuos} (déchets, mot plus technique, registre administratif ou journalistique) ne sont pas interchangeables dans tous les contextes. On dit \esp{tirar la basura} (jeter les ordures) mais \esp{la gestión de residuos} (la gestion des déchets).
\end{corrige}

\begin{corrige}[Exercice 4 — Conjugaison à compléter]
\begin{enumerate}
\item \esp{separamos} — \esp{separar}, 1\textsuperscript{re} personne du pluriel, présent régulier en \esp{-amos}.
\item \esp{recoge} — \esp{recoger}, 3\textsuperscript{e} personne du singulier. Attention au changement orthographique \esp{g} $\rightarrow$ \esp{j} devant \esp{o} et \esp{a} (\esp{yo recojo}, subjonctif \esp{recoja}), mais ici devant \esp{e} on garde \esp{g} : \esp{recoge}.
\item \esp{reutilizas} — \esp{reutilizar}, 2\textsuperscript{e} personne du singulier. Attention : \esp{z} $\rightarrow$ \esp{c} devant \esp{e} (\esp{yo reutilice} au subjonctif), mais devant \esp{a} on garde \esp{z} : \esp{reutilizas}.
\item \esp{pueden} — \esp{poder}, verbe à diphtongaison \esp{o} $\rightarrow$ \esp{ue} : \esp{puedo, puedes, puede, podemos, podéis, pueden}. La diphtongaison ne touche pas \esp{nosotros} ni \esp{vosotros}.
\item \esp{quiero} — \esp{querer}, verbe à diphtongaison \esp{e} $\rightarrow$ \esp{ie} : \esp{quiero, quieres, quiere, queremos, queréis, quieren}.
\end{enumerate}
\textbf{Point de vigilance :} les changements orthographiques (\esp{g/j}, \esp{z/c}) servent à conserver le son ; les diphtongaisons (\esp{o/ue}, \esp{e/ie}) concernent la racine et disparaissent aux personnes où l'accent tonique ne porte pas sur la racine.
\end{corrige}

\begin{corrige}[Exercice 5 — Traduction : thème]
\begin{enumerate}
\item \esp{Debemos separar (o clasificar) nuestros residuos para proteger el medioambiente.}
\item \esp{En Yaoundé, la recogida de la basura es un gran problema.}
\item \esp{Hay que reducir el plástico y reutilizar las bolsas.}
\item \esp{El vidrio y el papel pueden ser reciclados.} (ou : \esp{se pueden reciclar})
\item \esp{Los jóvenes deben aprender los gestos ecociudadanos en la escuela.}
\end{enumerate}
\textbf{Points de vigilance :}
\begin{itemize}
\item « Il faut » impersonnel se traduit par \esp{hay que + infinitivo}, jamais par un calque du français.
\item \esp{medioambiente} s'écrit en un seul mot.
\item « Trier » se dit \esp{separar} ou \esp{clasificar} ; ne pas utiliser \esp{triar}, qui n'appartient pas à l'espagnol standard.
\item Phrase 4 : au passif espagnol, le participe s'accorde avec le sujet : \esp{reciclados} (masculin pluriel, comme \esp{el vidrio y el papel}).
\end{itemize}
\end{corrige}

\begin{corrige}[Exercice 6 — Reformulation]
Propositions de reformulation (d'autres formulations sont possibles si le sens est conservé) :
\begin{enumerate}
\item \esp{El primer paso para reciclar es separar la basura según su materia.}
\item \esp{Tenemos que producir menos residuos.}
\item \esp{La basura del vertedero ensucia la tierra y el agua.}
\item \esp{Si volvemos a usar los objetos, duran más tiempo.}
\item \esp{Todos tenemos la obligación de cuidar el medioambiente.}
\end{enumerate}
\textbf{Méthode :} reformuler, c'est garder l'idée mais changer la structure et le vocabulaire (synonymes, transformation active/passive, verbe nominalisé). C'est exactement la compétence exigée pour citer et expliquer un passage dans un commentaire de texte. Noter l'expression utile \esp{volver a + infinitivo} = « faire de nouveau » : \esp{volvemos a usar} = « nous réutilisons ».
\end{corrige}

\begin{corrige}[Exercice 7 — Texte à trous : la recogida selectiva]
Texte complété :
\begin{enumerate}
\item \esp{recogida} (la collecte a lieu les lundis et jeudis)
\item \esp{residuos} (chaque famille trie ses déchets)
\item \esp{contenedor} (le papier va au conteneur bleu)
\item \esp{vidrio} (le verre va au conteneur vert)
\item \esp{reciclaje} (grâce à ce système, le recyclage est plus facile)
\item \esp{basura} (moins d'ordures arrivent à la décharge)
\item \esp{vertedero} (la décharge de la ville)
\item \esp{reducir} (réduire la consommation de produits très emballés)
\item \esp{reutilizar} (réutiliser les bocaux en verre)
\item \esp{contenedores} (les conteneurs de couleurs nous aident)
\end{enumerate}
\textbf{Point de vigilance :} accorder le singulier et le pluriel : \esp{el contenedor azul} (singulier) mais \esp{los contenedores de colores} (pluriel). Après une préposition (\esp{para}, \esp{sin}, \esp{de}), on emploie l'infinitif : \esp{para reducir}, \esp{para reutilizar}.
\end{corrige}

\begin{corrige}[Exercice 8 — Association : déchets et conteneurs]
Associations correctes et phrases complètes :
\begin{enumerate}
\item \esp{El plástico y los envases se echan en el contenedor amarillo.}
\item \esp{El vidrio se echa en el contenedor verde.}
\item \esp{El papel y el cartón se echan en el contenedor azul.}
\item \esp{La materia orgánica se echa en el contenedor marrón.}
\end{enumerate}
\textbf{Point de vigilance :} ne pas confondre \esp{vidrio} (le verre, matériau recyclable) et \esp{cristal} (le verre fin, les objets en verre de table). Pour le tri, on dit \esp{el vidrio}. Le verbe \esp{echar} (jeter) est très utile : \esp{echar la basura} = jeter les ordures ; on peut aussi dire \esp{tirar la basura}. Noter l'accord de l'adjectif de couleur : \esp{el contenedor marrón} (accent sur le masculin singulier).
\end{corrige}

\begin{corrige}[Exercice 9 — Compréhension écrite type Bac]
\textbf{Barème indicatif (sur 10 points) :} questions de compréhension : 5 $\times$ 1 pt ; vrai/faux justifié : 2 $\times$ 1,5 pt ; question de langue : 2 pts.

\textbf{Questions de compréhension — réponses modèles :}
\begin{enumerate}
\item \esp{Cada día recorren las calles para recoger cartón, plástico y metal que después venden a las empresas de reciclaje.}
\item \esp{Su trabajo es importante porque, gracias a ellos, una parte importante de los residuos no termina en los vertederos.}
\item \esp{Viven en condiciones difíciles : muchos no tienen contrato, ni seguro, ni protección contra las enfermedades.}
\item \esp{Las cooperativas organizan a los recicladores, les ofrecen formación y mejoran sus ingresos.}
\item \esp{Las tres claves son la recogida selectiva en las casas, la educación de los niños en las escuelas y el reconocimiento del trabajo de los recicladores.}
\end{enumerate}

\textbf{Vrai ou faux :}
\begin{enumerate}
\item \textbf{Vrai} : \esp{« recogen cartón, plástico y metal que después venden a las empresas de reciclaje »}.
\item \textbf{Faux} : \esp{« Muchos no tienen contrato, ni seguro, ni protección contra las enfermedades »}.
\end{enumerate}

\textbf{Question de langue :} trois mots de la famille de \esp{reciclar} : \esp{el reciclaje} (le recyclage), \esp{los recicladores} (les récupérateurs / recycleurs), \esp{reciclar} lui-même (recycler). On peut ajouter \esp{reciclado/a} (recyclé).

\textbf{Points de vigilance :} répondre par des phrases complètes en espagnol, sans recopier tout le paragraphe ; pour le vrai/faux, la justification doit être une citation exacte entre guillemets ; attention à l'accord \esp{condiciones difíciles} et à la négation \esp{ni... ni...} ; \esp{les ofrecen} : pronom d'objet indirect pluriel pour « leur offrent ».
\end{corrige}

\begin{corrige}[Exercice 10 — Thème et version]
\textbf{A. Thème — traductions proposées :}
\begin{enumerate}
\item \esp{Los habitantes de Douala tiran demasiados residuos (o demasiada basura) en las calles.}
\item \esp{Si separamos nuestra basura, el reciclaje será más fácil.}
\item \esp{Es necesario que los escolares aprendan a proteger la naturaleza.}
\end{enumerate}
\textbf{Points de vigilance :}
\begin{itemize}
\item Phrase 2 : après \esp{si} de condition, on emploie le présent de l'indicatif (\esp{si separamos}), jamais le futur ni le subjonctif ; le futur \esp{será} se place dans la principale.
\item Phrase 3 : \esp{es necesario que} entraîne le \textbf{subjonctif} : \esp{aprendan} (et non \esp{aprenden}). C'est la structure \esp{es + adjectif + que + subjuntivo}. De même : \esp{es importante que}, \esp{es preciso que}.
\item « Jeter » se traduit par \esp{tirar} ou \esp{echar}, pas par un calque du français.
\end{itemize}

\textbf{B. Version — traduction proposée :}

« Depuis un an, les voisins de notre quartier trient les ordures à la maison. Le lundi, un camion ramasse le papier et le carton ; le jeudi, le verre et le plastique. Grâce à cet effort, les rues sont plus propres et les enfants jouent sans danger. Nous sommes tous fiers de notre quartier. »

\textbf{Points de vigilance :}
\begin{itemize}
\item \esp{Desde hace un año} + présent = « depuis un an » + présent en français (action qui continue encore).
\item \esp{estar orgulloso de} = être fier de ; ici \esp{estamos orgullosos} exprime un état, d'où \esp{estar} et non \esp{ser}.
\item \esp{sin peligro} = sans danger ; ne pas traduire \esp{peligro} par « péril » (faux ami de registre).
\item \esp{las calles están más limpias} : \esp{estar} + adjectif pour un résultat, un état passager (« sont plus propres » grâce à l'effort).
\end{itemize}
\end{corrige}

\begin{corrige}[Exercice 11 — Production écrite type Bac]
\textbf{Barème indicatif (sur 10 points) :} respect des consignes et du nombre de mots : 2 pts ; richesse du lexique de la leçon (au moins 5 mots) : 2 pts ; variété des structures de conseil (au moins 3) : 2 pts ; correction grammaticale et orthographique : 3 pts ; cohérence et phrase finale d'appel à l'action : 1 pt.

\textbf{Proposition de rédaction modèle (environ 95 mots) :}

\esp{En mi ciudad, todos podemos adoptar gestos ecociudadanos. Primero, hay que separar la basura en casa : el papel, el vidrio y el plástico deben ir a contenedores diferentes. Además, debemos reducir el consumo de plástico y reutilizar las bolsas de tela cuando vamos al mercado. También es importante apagar la luz y cerrar el grifo para ahorrar energía y agua. Por último, se debe mantener limpias las calles y no tirar residuos por el suelo. Si cada ciudadano participa, el reciclaje será más fácil y nuestra ciudad será más sana. ¡Cuidemos nuestro medioambiente!}

\textbf{Traduction :} « Dans ma ville, nous pouvons tous adopter des gestes écocitoyens. D'abord, il faut trier les ordures à la maison : le papier, le verre et le plastique doivent aller dans des conteneurs différents. De plus, nous devons réduire la consommation de plastique et réutiliser les sacs en tissu quand nous allons au marché. Il est important aussi d'éteindre la lumière et de fermer le robinet pour économiser l'énergie et l'eau. Enfin, on doit garder les rues propres et ne pas jeter de déchets par terre. Si chaque citoyen participe, le recyclage sera plus facile et notre ville sera plus saine. Protégeons notre environnement ! »

\textbf{Points de vigilance :}
\begin{itemize}
\item Vérifier la présence des trois structures de conseil : \esp{hay que + inf.}, \esp{debemos + inf.}, \esp{es importante + inf.}, \esp{se debe + inf.}
\item La phrase finale au subjonctif exhortatif : \esp{¡Cuidemos nuestro medioambiente!} (1\textsuperscript{re} personne du pluriel du subjonctif présent de \esp{cuidar}), avec les signes ¡ !
\item Orthographe : \esp{medioambiente} en un mot ; accents sur \esp{plástico}, \esp{energía} ; \esp{grifo} (robinet) et non un calque du français.
\item Compter les mots : rester entre 80 et 100 mots, ni trop court ni trop long.
\end{itemize}
\end{corrige}

\newpage

\coursec{Fiche bilan}

\begin{popfigure}
\begin{tikzpicture}[mindmap, grow cyclic, every node/.style={concept, execute at begin node=\hspace{0pt}}, concept color=popblue!60, root concept/.append style={concept color=popblue, font=\bfseries, text=white, minimum size=2.6cm}, level 1 concept/.append style={level distance=3.4cm, sibling angle=51.5, font=\small, text=white}, text width=2.2cm, align=center]
\node[root concept]{El medioambiente y nosotros}
  child[concept color=popgreen]{ node[align=center]{Lexique des déchets\\ \esp{basura, residuos, contenedor}} }
  child[concept color=poporange]{ node[align=center]{Chaîne du recyclage\\ \esp{recogida, reciclaje, vertedero}} }
  child[concept color=poppurple]{ node[align=center]{Règle des 3 R\\ \esp{reducir, reutilizar, reciclar}} }
  child[concept color=poppink]{ node[align=center]{Grammaire\\ futur simple, subjonctif, impératif} }
  child[concept color=popgold, text=popdark]{ node[align=center]{Méthode du commentaire\\ 4 étapes} }
  child[concept color=popblue!80]{ node[align=center]{Contexte camerounais\\ Douala, Yaoundé, Hysacam} }
  child[concept color=popgreen!70!black]{ node[align=center]{Gestes écocitoyens\\ \esp{Hay que\ldots}} };
\end{tikzpicture}
\legende{Carte mentale de la leçon : les sept piliers du commentaire sur les déchets et le recyclage.}
\end{popfigure}

\begin{bilan}
\textbf{1. Lexique clé à connaître par cœur}
\begin{center}
\begin{tabularx}{\linewidth}{|X|X|}\hline
\rowcolor{popblueL} \textbf{Español} & \textbf{Français} \\ \hline
\esp{la basura / los residuos} & les ordures / les déchets \\ \hline
\esp{el contenedor (amarillo, verde, azul)} & le conteneur (plastique, verre, papier) \\ \hline
\esp{la recogida selectiva} & le tri sélectif \\ \hline
\esp{el reciclaje / reciclar} & le recyclage / recycler \\ \hline
\esp{el vertedero} & la décharge \\ \hline
\esp{reutilizar / reducir} & réutiliser / réduire \\ \hline
\esp{el medioambiente} & l'environnement \\ \hline
\esp{contaminar / la contaminación} & polluer / la pollution \\ \hline
\esp{proteger el planeta} & protéger la planète \\ \hline
\esp{los gestos ecociudadanos} & les gestes écocitoyens \\ \hline
\end{tabularx}
\end{center}

\textbf{2. Grammaire mobilisée dans ce type de texte}
\begin{center}
\begin{tabularx}{\linewidth}{|X|X|}\hline
\rowcolor{popgreenL} \textbf{Outil} & \textbf{Emploi dans le commentaire} \\ \hline
Futur simple (\esp{reciclaremos}) & prévisions, conséquences, promesses \\ \hline
\esp{hay que} + infinitif & exprimer l'obligation générale (\esp{Hay que reducir los residuos.}) \\ \hline
Subjonctif après \esp{es importante que, para que} & recommandations et buts \\ \hline
Impératif (\esp{¡Recicla! ¡No tires!}) & slogans et appels à l'action \\ \hline
\esp{se} passif (\esp{se recicla el vidrio}) & décrire les procédures de traitement \\ \hline
\end{tabularx}
\end{center}

\textbf{3. Méthode express du commentaire de texte}
\begin{itemize}
\item \cle{Étape 1} : présenter le texte (auteur, source, thème, genre) en 2 ou 3 phrases.
\item \cle{Étape 2} : dégager l'idée générale et annoncer un plan en deux ou trois axes.
\item \cle{Étape 3} : analyser le lexique (champ lexical des déchets) et les procédés (impératifs, futur, chiffres).
\item \cle{Étape 4} : conclure par un jugement personnel et une ouverture (exemple : la situation à Douala ou à Yaoundé).
\end{itemize}

\textbf{4. Réflexes de traduction}
\begin{itemize}
\item « les déchets » se dit \esp{los residuos} (registre soutenu) ou \esp{la basura} (courant), jamais \esp{los desechos} au sens de « déchets ménagers » au Bac.
\item « trier » = \esp{separar / clasificar la basura} ; « jeter » = \esp{tirar} (faux ami : \esp{jeter} n'existe pas).
\item « il faut » = \esp{hay que} + infinitif, sans sujet.
\end{itemize}
\end{bilan}

\begin{aretenir}[Checklist avant le Bac]
\begin{itemize}
\item[$\square$] Je connais le lexique : \esp{basura, residuos, contenedor, vertedero, reciclaje}.
\item[$\square$] Je sais expliquer la règle des 3 R : \esp{reducir, reutilizar, reciclar}.
\item[$\square$] Je maîtrise \esp{hay que} + infinitif pour formuler des obligations.
\item[$\square$] Je conjugue le futur simple sans erreur (\esp{reciclaré, reciclarás, reciclará\ldots}).
\item[$\square$] Je sais employer le subjonctif après \esp{es importante que} et \esp{para que}.
\item[$\square$] Je forme l'impératif affirmatif et négatif (\esp{¡Recicla! / ¡No contamines!}).
\item[$\square$] Je connais les 4 étapes du commentaire de texte.
\item[$\square$] Je sais présenter un texte en 3 phrases (auteur, source, thème).
\item[$\square$] J'évite les faux amis : \esp{tirar} = jeter, \esp{los residuos} = les déchets.
\item[$\square$] Je peux citer des exemples camerounais (Hysacam à Douala et Yaoundé).
\item[$\square$] Je sais proposer trois gestes écocitoyens en espagnol.
\item[$\square$] Je vérifie mes accents : \esp{contaminación, plástico, vidrio, energía}.
\end{itemize}
\end{aretenir}

\findecours

\end{document}
